\documentclass[14pt]{extarticle}

% 16:9 page (presentation aspect ratio) - tightened margins
\usepackage[
  paperwidth=257mm,
  paperheight=144.5mm,
  top=9mm,
  bottom=7mm,
  left=5mm,
  right=5mm,
  headheight=7mm,
  headsep=2mm,
  footskip=4mm
]{geometry}

\usepackage{amsmath,amsthm,amssymb,amsfonts}
\usepackage{booktabs,array}
\usepackage{enumitem}
\usepackage{xcolor}
\usepackage{fancyhdr}
\usepackage[most]{tcolorbox}
\tcbuselibrary{theorems,skins,breakable}
\usepackage{hyperref}
\hypersetup{colorlinks=true,linkcolor=blue!60!black,urlcolor=blue!70!black,citecolor=blue!60!black}

%%% COLOR PALETTE %%%
\definecolor{hdrBG}{HTML}{1E3A8A}        % header blue
\definecolor{accent}{HTML}{0EA5E9}       % subsection cyan
\definecolor{defBG}{HTML}{EFF6FF}        % definition: soft blue
\definecolor{defFR}{HTML}{2563EB}
\definecolor{thmBG}{HTML}{FEF3C7}        % theorem: soft amber
\definecolor{thmFR}{HTML}{B45309}
\definecolor{lemBG}{HTML}{FEF9C3}        % lemma: lighter amber
\definecolor{lemFR}{HTML}{A16207}
\definecolor{corBG}{HTML}{FDF2F8}        % corollary: soft pink
\definecolor{corFR}{HTML}{BE185D}
\definecolor{propBG}{HTML}{F3E8FF}       % proposition: soft purple
\definecolor{propFR}{HTML}{7C3AED}
\definecolor{proofBG}{HTML}{F8FAFC}      % proof: very soft gray
\definecolor{proofFR}{HTML}{64748B}
\definecolor{remBG}{HTML}{ECFDF5}        % remark: soft green
\definecolor{remFR}{HTML}{047857}
\definecolor{exBG}{HTML}{FFF7ED}         % exercise: soft orange
\definecolor{exFR}{HTML}{C2410C}

%%% HEADER / FOOTER %%%
\pagestyle{fancy}
\fancyhf{}
\renewcommand{\headrulewidth}{0.4pt}
\renewcommand{\headrule}{\color{hdrBG!30}\hrule width\textwidth height0.4pt}
\renewcommand{\footrulewidth}{0pt}
\fancyhead[L]{\small\color{hdrBG}\textbf{Classical Linear Models}}
\fancyhead[R]{\small\color{hdrBG}\leftmark}
\fancyfoot[R]{\small\color{gray}\thepage}
\renewcommand{\sectionmark}[1]{\markboth{\S\thesection~#1}{}}

%%% AMSTHM ENVIRONMENTS %%%
% Numbering as in clm.tex: all theorem-like environments (remarks included) share one
% counter per section, so every number here is the note's; exercises are numbered by
% tier (A1, ..., D4) via \exercisetier
\theoremstyle{plain}
\newtheorem{theorem}{Theorem}[section]
\newtheorem{lemma}[theorem]{Lemma}
\newtheorem{corollary}[theorem]{Corollary}
\newtheorem{proposition}[theorem]{Proposition}
\theoremstyle{definition}
\newtheorem{definition}[theorem]{Definition}
\newtheorem{example}[theorem]{Example}
\newtheorem{property}[theorem]{Property}
\newtheorem{exercise}{Exercise}
\newcommand{\extier}{}
\renewcommand{\theexercise}{\extier\arabic{exercise}}
\def\theHexercise{\extier\arabic{exercise}}% keeps hyperref anchors distinct across tiers
\newcommand{\exercisetier}[1]{\renewcommand{\extier}{#1}\setcounter{exercise}{0}}
\theoremstyle{remark}
\newtheorem{remark}[theorem]{Remark}

%%% WRAP ENVIRONMENTS IN COLORED BOXES %%%
%%% Theorem-like statements: NOT breakable - must appear whole on a single page %%%
%%% Proof: breakable - long proofs need to span pages %%%
\tcolorboxenvironment{definition}{
  enhanced,breakable=false,
  colback=defBG,colframe=defFR,
  boxrule=0.6pt,arc=2pt,
  left=3mm,right=3mm,top=1.5mm,bottom=1.5mm,
  before skip=5pt,after skip=5pt
}
\tcolorboxenvironment{theorem}{
  enhanced,breakable=false,
  colback=thmBG,colframe=thmFR,
  boxrule=0.6pt,arc=2pt,
  left=3mm,right=3mm,top=1.5mm,bottom=1.5mm,
  before skip=5pt,after skip=5pt
}
\tcolorboxenvironment{lemma}{
  enhanced,breakable=false,
  colback=lemBG,colframe=lemFR,
  boxrule=0.6pt,arc=2pt,
  left=3mm,right=3mm,top=1.5mm,bottom=1.5mm,
  before skip=5pt,after skip=5pt
}
\tcolorboxenvironment{corollary}{
  enhanced,breakable=false,
  colback=corBG,colframe=corFR,
  boxrule=0.6pt,arc=2pt,
  left=3mm,right=3mm,top=1.5mm,bottom=1.5mm,
  before skip=5pt,after skip=5pt
}
\tcolorboxenvironment{proposition}{
  enhanced,breakable=false,
  colback=propBG,colframe=propFR,
  boxrule=0.6pt,arc=2pt,
  left=3mm,right=3mm,top=1.5mm,bottom=1.5mm,
  before skip=5pt,after skip=5pt
}
\tcolorboxenvironment{remark}{
  enhanced,breakable,
  colback=remBG,colframe=remFR,
  leftrule=3pt,rightrule=0pt,toprule=0pt,bottomrule=0pt,
  arc=0pt,left=3mm,right=2mm,top=1mm,bottom=1mm,
  before skip=4pt,after skip=4pt
}
\tcolorboxenvironment{proof}{
  enhanced,breakable,
  colback=proofBG,colframe=proofFR,
  leftrule=2pt,rightrule=0pt,toprule=0pt,bottomrule=0pt,
  arc=0pt,left=3mm,right=2mm,top=1mm,bottom=1mm,
  before skip=4pt,after skip=4pt
}
\tcolorboxenvironment{exercise}{
  enhanced,breakable,
  colback=exBG,colframe=exFR,
  leftrule=3pt,rightrule=0pt,toprule=0pt,bottomrule=0pt,
  arc=0pt,left=3mm,right=2mm,top=1mm,bottom=1mm,
  before skip=4pt,after skip=4pt
}

%%% CUSTOM COMMANDS (same as clm.tex) %%%
\newcommand\expc{\mathsf{E}}
\newcommand\prb{\mathsf{P}}
\DeclareMathOperator\var{var}
\DeclareMathOperator\cov{cov}
\DeclareMathOperator\corr{corr}
\DeclareMathOperator\plim{plim}
\DeclareMathOperator\tr{tr}
\DeclareMathOperator\rank{rank}
\DeclareMathOperator\diag{diag}
\newcommand{\ds}{\displaystyle}
\newcommand{\ie}{\;\Longrightarrow\;}
\newcommand{\ifff}{\;\Longleftrightarrow\;}
\newcommand{\SST}{\mathrm{SST}}
\newcommand{\SSR}{\mathrm{SSR}}
\newcommand{\SSE}{\mathrm{SSE}}
\newcommand{\MSR}{\mathrm{MSR}}
\newcommand{\MSE}{\mathrm{MSE}}
\newcommand{\SE}{\mathrm{SE}}
\newcommand{\bX}{\mathbf{X}}
\newcommand{\bY}{\mathbf{Y}}
\newcommand{\by}{\mathbf{y}}
\newcommand{\bx}{\mathbf{x}}
\def\bb{\boldsymbol{\beta}}
\newcommand{\be}{\boldsymbol{\varepsilon}}
\newcommand{\bI}{\mathbf{I}}
\newcommand{\bH}{\mathbf{H}}
\newcommand{\bM}{\mathbf{M}}
\newcommand{\bR}{\mathbf{R}}
\newcommand{\br}{\mathbf{r}}
\def\bc{\mathbf{c}}
\newcommand{\bone}{\mathbf{1}}
\newcommand{\bzero}{\mathbf{0}}
\newcommand{\bz}{\mathbf{z}}
\newcommand{\bw}{\mathbf{w}}
\newcommand{\bA}{\mathbf{A}}
\newcommand{\bC}{\mathbf{C}}
\newcommand{\bD}{\mathbf{D}}
\newcommand{\bP}{\mathbf{P}}
\newcommand{\ba}{\mathbf{a}}
\newcommand{\bv}{\mathbf{v}}
\newcommand{\Real}{\mathbb{R}}
\newcommand{\dd}{\mathrm{d}}
\newcommand{\calC}{\mathcal{C}}
\newcommand{\bJ}{\mathbf{J}}
\newcommand{\indep}{\perp\!\!\!\perp}
\newcommand{\bu}{\mathbf{u}}
\newcommand{\VIF}{\mathrm{VIF}}

%%% SPACING %%%
\usepackage{parskip}
\setlength{\parskip}{3pt}
\setlength{\jot}{4pt}

%%% SECTION STYLING %%%
\usepackage{titlesec}
\titleformat{\section}
  {\color{hdrBG}\normalfont\Large\bfseries}{\thesection}{0.7em}{}
\titleformat{\subsection}
  {\color{accent}\normalfont\large\bfseries}{\thesubsection}{0.6em}{}
\titlespacing*{\section}{0pt}{4pt}{2pt}
\titlespacing*{\subsection}{0pt}{4pt}{1pt}

\title{\color{hdrBG}\bfseries\Huge Classical Theory of Linear Models}
\author{}
\date{}

\begin{document}

%%% TITLE PAGE %%%
\begin{titlepage}
\thispagestyle{empty}
\vspace*{\fill}
\begin{center}
{\color{hdrBG}\bfseries\fontsize{36}{44}\selectfont Classical Theory of Linear Models}\\[1.2em]
{\color{accent}\bfseries\Large Slide edition}\\[2em]
\rule{60mm}{0.8pt}\\[1em]
{\large A complete treatment of OLS, Gauss--Markov,\\[0.2em] and exact inference under normality}\\[2em]
\rule{60mm}{0.8pt}
\end{center}
\vspace*{\fill}
\end{titlepage}

%=============================================================================
\section{The Multiple Linear Regression Model}
%=============================================================================

\subsection{Model Specification and Matrix Notation}

\begin{definition}[Multiple Linear Regression Model]\label{def:model}
Given $n$ observations, each consisting of a dependent variable $Y_i$ and $k$ independent variables $X_{i1}, \ldots, X_{ik}$, the \textbf{multiple linear regression model} assumes
\begin{equation}
Y_i = \beta_0 + \beta_1 X_{i1} + \beta_2 X_{i2} + \cdots + \beta_k X_{ik} + \varepsilon_i, \quad i = 1, 2, \ldots, n,
\end{equation}
where $\beta_0$ is the intercept, $\beta_1, \ldots, \beta_k$ are slope coefficients, and $\varepsilon_i$ is the error term. Equivalently, in matrix form,
\begin{equation}\label{eq:model-matrix}
\boxed{\bY = \bX\bb + \be}
\end{equation}
where
\[
\bY = \begin{pmatrix} Y_1 \\ Y_2 \\ \vdots \\ Y_n \end{pmatrix}_{n \times 1}, \quad
\bX = \begin{pmatrix} 1 & X_{11} & \cdots & X_{1k} \\ 1 & X_{21} & \cdots & X_{2k} \\ \vdots & \vdots & \ddots & \vdots \\ 1 & X_{n1} & \cdots & X_{nk} \end{pmatrix}_{n \times (k+1)}, \quad
\bb = \begin{pmatrix} \beta_0 \\ \beta_1 \\ \vdots \\ \beta_k \end{pmatrix}_{(k+1) \times 1}, \quad
\be = \begin{pmatrix} \varepsilon_1 \\ \varepsilon_2 \\ \vdots \\ \varepsilon_n \end{pmatrix}_{n \times 1}.
\]
The matrix $\bX$ is called the \textbf{design matrix}. Its first column is the all-ones vector $\bone = (1, 1, \ldots, 1)'$, corresponding to the intercept. We write $\bx_i'$ for the $i$-th row of $\bX$, so that $Y_i = \bx_i'\bb + \varepsilon_i$.
\end{definition}

\subsection{Classical Assumptions}

\begin{definition}[Classical Assumptions]\label{def:assumptions}
The following assumptions are imposed on model~\eqref{eq:model-matrix}:
\begin{enumerate}[label=\textup{(A\arabic*)}]
\item\label{A1} \textbf{Linearity}: $\expc[\bY \mid \bX] = \bX\bb$.
\item\label{A2} \textbf{Zero conditional mean}: $\expc[\be \mid \bX] = \bzero$.
\item\label{A3} \textbf{Spherical errors}: $\var(\be \mid \bX) = \sigma^2 \bI_n$ for a constant $\sigma^2 > 0$. This comprises two sub-conditions:
\begin{itemize}
    \item \emph{Homoscedasticity}: $\var(\varepsilon_i) = \sigma^2$ for all $i$ (a constant).
    \item \emph{No serial correlation}: $\cov(\varepsilon_i, \varepsilon_j) = 0$ for $i \neq j$.
\end{itemize}
\item\label{A4} \textbf{Full rank}: $\rank(\bX) = k + 1$ (the columns of $\bX$ are linearly independent). This requires $n \geqslant k + 1$. For variance estimation and inference (Sections~\ref{sec:finite}--\ref{sec:test} and~\ref{sec:pred}, and wherever $s$ appears in later sections), we additionally require $n > k + 1$ so that $\MSE = \SSE/(n-k-1)$ (see Definition~\ref{def:sumsq}) is well-defined.
\item\label{A5} \textbf{Normality}: $\be \mid \bX \sim N(\bzero, \sigma^2 \bI_n)$ (the multivariate normal distribution $N(\boldsymbol\mu, \boldsymbol\Sigma)$ is defined in Definition~\ref{def:mvn}, Section~\ref{sec:mvn}).
\end{enumerate}
Assumptions~\ref{A1}--\ref{A4} are called the \textbf{Gauss--Markov assumptions}. Assumption~\ref{A5} is additionally needed for exact distributional results in finite samples.
\end{definition}

\begin{remark}\label{rem:fixed-X}
Throughout this article, we treat $\bX$ as fixed (non-stochastic) and condition on it. All expectations, variances, and distributions are conditional on $\bX$, but we suppress the conditioning notation for brevity.
\end{remark}

\begin{remark}[Relationship between \ref{A1} and \ref{A2}]\label{rem:A1-A2}
Given the model $\bY = \bX\bb + \be$, assumptions \ref{A1} and \ref{A2} are logically equivalent: taking conditional expectations of both sides gives $\expc[\bY\mid\bX] = \bX\bb + \expc[\be\mid\bX]$, so $\expc[\bY\mid\bX] = \bX\bb$ holds if and only if $\expc[\be\mid\bX] = \bzero$. We state both because each form is convenient in different proofs: \ref{A1} is closer to how the model is written, and \ref{A2} is closer to how residuals are handled. Conditioning on $\bX$ (previous remark) also means that inside any proof, $\expc[\be] = \bzero$ and $\var(\be) = \sigma^2\bI$ can be used without further notation---these are the conditional expectation and variance with the conditioning suppressed.
\end{remark}

\subsection{Prerequisites from Linear Algebra (assumed without proof)}

The following standard facts are used below without proof; they are the only linear-algebra results imported from outside these notes. (The probabilistic prerequisites are listed in Section~\ref{sec:prereq-prob}.) Notation: for $\bA \in \Real^{m \times p}$, $\calC(\bA) = \{\bA\bv : \bv \in \Real^p\} \subseteq \Real^m$ is the \textbf{column space} of $\bA$ and $\mathcal{N}(\bA) = \{\bv \in \Real^p : \bA\bv = \bzero\}$ its \textbf{null space}; a symmetric matrix $\bA$ is \textbf{positive semi-definite}, written $\bA \succeq \bzero$, if $\bv'\bA\bv \geqslant 0$ for all $\bv$, and \textbf{positive definite} if $\bv'\bA\bv > 0$ for all $\bv \neq \bzero$. Vectors $\bv, \bw \in \Real^m$ are \textbf{orthogonal}, written $\bv \perp \bw$, if $\bv'\bw = 0$; for a subspace $\mathcal{S} \subseteq \Real^m$, $\bv \perp \mathcal{S}$ means $\bv \perp \bw$ for every $\bw \in \mathcal{S}$, and $\mathcal{S}^\perp = \{\bv \in \Real^m : \bv \perp \mathcal{S}\}$ is the \textbf{orthogonal complement} of $\mathcal{S}$; for a single vector $\ba$ we write $\ba^\perp = \calC(\ba)^\perp = \{\bv \in \Real^m : \ba'\bv = 0\}$.
\begin{enumerate}[label=\textup{(L\arabic*)}]
\item\label{L1} \emph{Rank--nullity.} For $\bA \in \Real^{m \times p}$, $\dim\mathcal{N}(\bA) + \rank(\bA) = p$. Hence two matrices with the same number of columns and the same null space have the same rank. (Used in Theorem~\ref{thm:ols}, Step~3, and in the remark following Corollary~\ref{cor:nested-F}.)
\item\label{L2} \emph{Trace.} $\tr(\bA\mathbf{B}) = \tr(\mathbf{B}\bA)$ whenever both products are defined; in particular $\tr(\bP\boldsymbol\Lambda\bP^{-1}) = \tr(\boldsymbol\Lambda)$. (Used in Lemma~\ref{lem:idempotent} and in the solution to Exercise~\ref{ex:A2}.)
\item\label{L3} \emph{Spectral theorem.} Every symmetric $\bA \in \Real^{n \times n}$ can be written $\bA = \bP\boldsymbol\Lambda\bP'$ with $\bP$ orthogonal ($\bP'\bP = \bP\bP' = \bI$) and $\boldsymbol\Lambda$ diagonal, its entries being the eigenvalues of $\bA$. (Used in Theorem~\ref{thm:quadratic-chisq}, in~\ref{L4}, and in the solutions to Exercises~\ref{ex:A2} and~\ref{ex:A8}.)
\item\label{L4} \emph{Square roots of positive definite matrices.} If $\bA$ is symmetric positive definite, then $\bA = \bP\boldsymbol\Lambda\bP'$ as in \ref{L3} with all eigenvalues positive, and $\bA^{1/2} := \bP\boldsymbol\Lambda^{1/2}\bP'$ is symmetric positive definite with $\bA^{1/2}\bA^{1/2} = \bA$; its inverse $\bA^{-1/2} := (\bA^{1/2})^{-1}$ satisfies $\bA^{-1/2}\bA\bA^{-1/2} = \bI$. In particular $\bA = \mathbf{L}\mathbf{L}'$ for some invertible $\mathbf{L}$ (for instance $\mathbf{L} = \bA^{1/2}$, or the Cholesky factor). (Used in Theorems~\ref{thm:general-F} and~\ref{thm:ellipsoid} and in the solutions to Exercises~\ref{ex:B4} and~\ref{ex:B5}.)
\end{enumerate}

%=============================================================================
\section{Ordinary Least Squares}\label{sec:ols}
%=============================================================================

Standing assumption: from here to the end of these notes, assumption~\ref{A4} (full column rank of $\bX$) is in force. It is what makes $(\bX'\bX)^{-1}$---and hence $\hat{\bb}$, $\bH$ and $\bM$---exist, and it is restated in a theorem only when the proof uses it in an additional way. Assumptions~\ref{A1}--\ref{A3} and~\ref{A5} are invoked explicitly where needed; the results of this section use none of them.

\subsection{Definition and the Normal Equations}

\begin{definition}[Ordinary Least Squares]\label{def:ols}
The \textbf{OLS estimator} $\hat{\bb}$ is the value of $\bb$ that minimizes the sum of squared residuals:
\begin{equation}
Q(\bb) = \sum_{i=1}^{n}(Y_i - \bx_i'\bb)^2 = (\bY - \bX\bb)'(\bY - \bX\bb) = \|\bY - \bX\bb\|^2.
\end{equation}
\end{definition}

\begin{lemma}[Matrix Differentiation Formulas]\label{lem:matrix-diff}
For $\bb \in \Real^p$, constant $\ba \in \Real^p$ and constant $\bA \in \Real^{p\times p}$, write $Df(\bb) := \partial f/\partial\bb'$ (numerator layout, row vector). Then
\begin{enumerate}[label=\textup{(\roman*)}]
\item $D(\ba'\bb) = D(\bb'\ba) = \ba'$;
\item $D(\bb'\bA\bb) = \bb'(\bA + \bA')$, which equals $2\bb'\bA$ when $\bA' = \bA$.
\end{enumerate}
\end{lemma}

\begin{proof}
Both functions are polynomials in the coordinates $\beta_1, \ldots, \beta_p$ of $\bb$, so we compute the partial derivatives coordinate by coordinate; $Df(\bb)$ is the row vector $(\partial f/\partial\beta_1, \ldots, \partial f/\partial\beta_p)$.

\textbf{(i).} $\ba'\bb = \sum_{i=1}^p a_i\beta_i$, so $\partial(\ba'\bb)/\partial\beta_j = a_j$ for each $j$; assembling the row vector gives $D(\ba'\bb) = (a_1, \ldots, a_p) = \ba'$. The case $D(\bb'\ba)$ follows from $\bb'\ba = \ba'\bb$.

\textbf{(ii).} $\bb'\bA\bb = \sum_{i=1}^p\sum_{l=1}^p \beta_i A_{il}\beta_l$. The coordinate $\beta_j$ appears in the terms with $i = j$ and in the terms with $l = j$ (the term $i = l = j$ is $A_{jj}\beta_j^2$, whose derivative $2A_{jj}\beta_j$ is counted once in each group), so
\[
\frac{\partial(\bb'\bA\bb)}{\partial\beta_j} = \sum_{l=1}^p A_{jl}\beta_l + \sum_{i=1}^p \beta_i A_{ij} = (\bA\bb)_j + (\bA'\bb)_j.
\]
Hence $D(\bb'\bA\bb) = (\bA\bb + \bA'\bb)' = \bb'(\bA + \bA')$, which equals $2\bb'\bA$ when $\bA' = \bA$.
\end{proof}

\begin{remark}[The same formulas via differentials]\label{rem:differentials}
Readers who know the differential calculus of Magnus and Neudecker (developed in the Matrix Calculus notes, and \emph{not} assumed here) can obtain the same formulas from the linearity and product rule of $\dd$ together with the \emph{first identification theorem} $\dd f(\bb) = \bc(\bb)'\dd\bb \ifff Df(\bb) = \bc(\bb)'$: since $\dd\ba = 0$ and $\dd\bA = 0$, $\dd(\ba'\bb) = \ba'\dd\bb$, and, using the scalar identity $(\dd\bb)'\bA\bb = \bb'\bA'\dd\bb$,
\[
\dd(\bb'\bA\bb) = (\dd\bb)'\bA\bb + \bb'\bA\,\dd\bb = \bb'\bA'\dd\bb + \bb'\bA\,\dd\bb = \bb'(\bA+\bA')\,\dd\bb.
\]
Nothing in these notes depends on this remark.
\end{remark}

\subsection{Derivation via Differentiation}

\begin{theorem}[OLS Estimator]\label{thm:ols}
Under assumption~\ref{A4}, the OLS estimator is
\begin{equation}
\boxed{\hat{\bb} = (\bX'\bX)^{-1}\bX'\bY.}
\end{equation}
\end{theorem}

\begin{proof}
\textbf{Step 1: Expand $Q(\bb)$.}
\begin{align}
Q(\bb) &= (\bY - \bX\bb)'(\bY - \bX\bb) \notag\\
&= \bY'\bY - \bY'\bX\bb - \bb'\bX'\bY + \bb'\bX'\bX\bb \notag\\
&= \bY'\bY - 2\bb'\bX'\bY + \bb'\bX'\bX\bb.
\end{align}
In the last line, $\bY'\bX\bb$ is a scalar, hence equals its transpose $\bb'\bX'\bY$.

\textbf{Step 2: First-order condition.} Apply Lemma~\ref{lem:matrix-diff} term by term, with $\ba = \bX'\bY$ and $\bA = \bX'\bX$ (symmetric):
\[
DQ(\bb) = -2(\bX'\bY)' + 2\bb'\bX'\bX = -2\bY'\bX + 2\bb'\bX'\bX.
\]
Setting $DQ = \bzero'$ and transposing gives the \textbf{normal equation}
\begin{equation}\label{eq:normal-eq}
\bX'\bX\hat{\bb} = \bX'\bY.
\end{equation}

\textbf{Step 3: Solve.} By~\ref{A4}, $\rank(\bX) = k+1$. We show $\rank(\bX'\bX) = \rank(\bX)$: $\bX\bv = \bzero$ trivially implies $\bX'\bX\bv = \bzero$; conversely, $\bX'\bX\bv = \bzero$ gives $\bv'\bX'\bX\bv = \|\bX\bv\|^2 = 0$, hence $\bX\bv = \bzero$. So $\bX$ and $\bX'\bX$ have the same null space, hence the same rank (\ref{L1}). Thus $\bX'\bX$ is $(k+1) \times (k+1)$ of full rank, hence invertible. Multiplying~\eqref{eq:normal-eq} by $(\bX'\bX)^{-1}$:
\[
\hat{\bb} = (\bX'\bX)^{-1}\bX'\bY.
\]

\textbf{Step 4: Verify the second-order condition.} The Hessian matrix is
\[
\frac{\partial^2 Q}{\partial \bb \partial \bb'} = 2\bX'\bX.
\]
Under assumption~\ref{A4}, for any nonzero $\bv \in \Real^{k+1}$, $\bX\bv \neq \bzero$, so $\bv'\bX'\bX\bv = \|\bX\bv\|^2 > 0$. Hence $\bX'\bX$ is positive definite, so $\hat{\bb}$ is a local minimum. (Its inverse $(\bX'\bX)^{-1}$ is positive definite as well, a fact needed in Theorems~\ref{thm:general-F} and~\ref{thm:cls}: it is symmetric, being the inverse of a symmetric matrix, and for $\bv \neq \bzero$ the vector $\bw = (\bX'\bX)^{-1}\bv$ is nonzero (as $\bv = \bX'\bX\bw$), so $\bv'(\bX'\bX)^{-1}\bv = \bw'\bX'\bX\bw = \|\bX\bw\|^2 > 0$.) Positive definiteness of the Hessian by itself certifies only a \emph{local} minimum; that the minimum is global and unique follows from an exact expansion. For any $\bb \in \Real^{k+1}$, write $\bY - \bX\bb = (\bY - \bX\hat{\bb}) - \bX(\bb - \hat{\bb})$; then
\begin{equation}\label{eq:ols-global}
Q(\bb) = Q(\hat{\bb}) - 2(\bb - \hat{\bb})'\bX'(\bY - \bX\hat{\bb}) + \|\bX(\bb - \hat{\bb})\|^2 = Q(\hat{\bb}) + \|\bX(\bb - \hat{\bb})\|^2 \;\geqslant\; Q(\hat{\bb}),
\end{equation}
because $\bX'(\bY - \bX\hat{\bb}) = \bzero$ is the normal equation~\eqref{eq:normal-eq}. Equality holds iff $\bX(\bb - \hat{\bb}) = \bzero$, i.e.\ (by~\ref{A4}) iff $\bb = \hat{\bb}$. Thus $\hat{\bb}$ is the unique global minimizer of $Q$.
\end{proof}

\subsection{Derivation via Orthogonal Projection}

\begin{theorem}[OLS as Orthogonal Projection]\label{thm:ols-proj}
The fitted vector $\hat{\bY} = \bX\hat{\bb}$ is the orthogonal projection of $\bY$ onto $\calC(\bX)$, that is, $\hat{\bY} \in \calC(\bX)$ and $\bY - \hat{\bY} \perp \calC(\bX)$ (the term is defined for a general subspace, in matrix form, in Remark~\ref{rem:hat-projection}), and the OLS estimator is $\hat{\bb} = (\bX'\bX)^{-1}\bX'\bY$.
\end{theorem}

\begin{proof}
The OLS problem is to minimize $Q(\bb) = \|\bY - \bX\bb\|^2$ over $\bb \in \Real^{k+1}$, equivalently, to find the point in $\calC(\bX) = \{\bX\bb : \bb \in \Real^{k+1}\}$ closest to $\bY$ in Euclidean norm.

\textbf{Step 1: Orthogonality characterization of the minimum.} We claim $\hat{\bY} \in \calC(\bX)$ minimizes $\|\bY - \tilde{\bY}\|^2$ over $\tilde{\bY} \in \calC(\bX)$ iff $\bY - \hat{\bY} \perp \calC(\bX)$.

($\Leftarrow$) Assume the orthogonality condition and let $\tilde{\bY} \in \calC(\bX)$. Since $\calC(\bX)$ is a subspace, $\hat{\bY} - \tilde{\bY} \in \calC(\bX)$, so $(\bY - \hat{\bY}) \perp (\hat{\bY} - \tilde{\bY})$. Pythagoras applied to $\bY - \tilde{\bY} = (\bY - \hat{\bY}) + (\hat{\bY} - \tilde{\bY})$ gives
\[
\|\bY - \tilde{\bY}\|^2 = \|\bY - \hat{\bY}\|^2 + \|\hat{\bY} - \tilde{\bY}\|^2 \;\geqslant\; \|\bY - \hat{\bY}\|^2,
\]
with equality iff $\tilde{\bY} = \hat{\bY}$. So $\hat{\bY}$ is the unique minimizer.

($\Rightarrow$) Let $\hat{\bY} \in \calC(\bX)$ minimize. Suppose, for contradiction, some $\bv \in \calC(\bX)$ has $\bv'(\bY - \hat{\bY}) \neq 0$; normalize to $\|\bv\|^2 = 1$. For $\tilde{\bY} = \hat{\bY} + t\bv \in \calC(\bX)$,
\[
\|\bY - \tilde{\bY}\|^2 = \|\bY - \hat{\bY}\|^2 - 2t\,\bv'(\bY - \hat{\bY}) + t^2,
\]
minimized at $t^* = \bv'(\bY - \hat{\bY}) \neq 0$ with value $\|\bY - \hat{\bY}\|^2 - (t^*)^2 < \|\bY - \hat{\bY}\|^2$, contradicting optimality. Hence $\bY - \hat{\bY} \perp \calC(\bX)$.

\textbf{Step 2: Translate the orthogonality condition to a normal equation.} The residual is orthogonal to every vector in $\calC(\bX)$ if and only if it is orthogonal to every column of $\bX$ (since those columns span $\calC(\bX)$):
\[
\bY - \hat{\bY} \perp \calC(\bX) \;\ifff\; \bX'(\bY - \hat{\bY}) = \bzero.
\]
Writing $\hat{\bY} = \bX\hat{\bb}$ for some $\hat{\bb}$, this becomes $\bX'\bX\hat{\bb} = \bX'\bY$, the same normal equation~\eqref{eq:normal-eq} obtained by differentiation.

\textbf{Step 3: Solve.} Under assumption~\ref{A4}, $\bX'\bX$ is invertible, so $\hat{\bb} = (\bX'\bX)^{-1}\bX'\bY$, and the minimizer $\hat{\bY} = \bX\hat{\bb}$ exists and is unique.
\end{proof}

\subsection{Hat Matrix, Residual Maker, and Idempotent Matrices}

\begin{definition}[Hat Matrix and Residual Maker Matrix]\label{def:hat-matrix}
The \textbf{hat matrix} is
\begin{equation}
\bH = \bX(\bX'\bX)^{-1}\bX'.
\end{equation}
The \textbf{residual maker matrix} (or \textbf{annihilator matrix}) is
\begin{equation}
\bM = \bI_n - \bH.
\end{equation}
Then $\hat{\bY} = \bX\hat{\bb} = \bH\bY$, and the \textbf{residual vector} $\mathbf{e} := \bY - \hat{\bY}$, with components $e_i = Y_i - \hat{Y}_i$, satisfies $\mathbf{e} = \bY - \bH\bY = \bM\bY$.
\end{definition}

\begin{definition}[Idempotent Matrix]\label{def:idempotent}
A square matrix $\bA$ is \textbf{idempotent} if $\bA^2 = \bA$.
\end{definition}

\begin{lemma}[Trace Equals Rank for Idempotent Matrices]\label{lem:idempotent}
If $\bA$ is an $n \times n$ idempotent matrix, then $\tr(\bA) = \rank(\bA)$.
\end{lemma}

\begin{proof}
We show every idempotent matrix is diagonalizable with eigenvalues in $\{0, 1\}$.

\textbf{Step 1: Eigenvalues are in $\{0, 1\}$.} If $\bA\mathbf{v} = \lambda\mathbf{v}$ with $\mathbf{v} \neq \bzero$, then $\lambda\mathbf{v} = \bA\mathbf{v} = \bA^2\mathbf{v} = \lambda^2\mathbf{v}$, forcing $\lambda \in \{0, 1\}$.

\textbf{Step 2: $\Real^n = \mathcal{N}(\bA) \oplus \calC(\bA)$.} Any $\mathbf{v} \in \Real^n$ decomposes as $\mathbf{v} = (\mathbf{v} - \bA\mathbf{v}) + \bA\mathbf{v}$, with $\bA\mathbf{v} \in \calC(\bA)$ and $\mathbf{v} - \bA\mathbf{v} \in \mathcal{N}(\bA)$ (since $\bA(\mathbf{v} - \bA\mathbf{v}) = \bA\mathbf{v} - \bA^2\mathbf{v} = \bzero$). The sum is direct: if $\mathbf{v} = \bA\mathbf{w} \in \calC(\bA) \cap \mathcal{N}(\bA)$, then $\mathbf{v} = \bA\mathbf{w} = \bA^2\mathbf{w} = \bA\mathbf{v} = \bzero$.

\textbf{Step 3: Diagonalizability.} Any $\mathbf{u} \in \calC(\bA)$ satisfies $\bA\mathbf{u} = \mathbf{u}$ ($\mathbf{u} = \bA\mathbf{w}$, so $\bA\mathbf{u} = \bA^2\mathbf{w} = \bA\mathbf{w} = \mathbf{u}$). Combining a basis of $\mathcal{N}(\bA)$ (eigenvalue $0$) and a basis of $\calC(\bA)$ (eigenvalue $1$) gives an eigenvector basis of $\Real^n$ by Step~2. Thus $\bA = \bP\boldsymbol\Lambda\bP^{-1}$ with $\boldsymbol\Lambda$ diagonal and entries in $\{0, 1\}$.

\textbf{Step 4: Conclude.} Let $r = \rank(\bA) = \dim\calC(\bA)$. Then $\boldsymbol\Lambda$ has $r$ ones and $n-r$ zeros, so by the cyclic property of trace (\ref{L2}),
\[
\tr(\bA) = \tr(\bP\boldsymbol\Lambda\bP^{-1}) = \tr(\boldsymbol\Lambda) = r = \rank(\bA). \qedhere
\]
\end{proof}

\begin{theorem}[Properties of $\bH$ and $\bM$]\label{thm:HM}
The matrices $\bH$ and $\bM$ satisfy:
\begin{enumerate}[label=\textup{(\roman*)}]
\item\label{HM:i} $\bH$ and $\bM$ are both symmetric and idempotent.
\item\label{HM:ii} $\bH\bX = \bX$ and $\bM\bX = \bzero$.
\item\label{HM:iii} $\bH\bM = \bM\bH = \bzero$.
\item\label{HM:iv} $\tr(\bH) = k + 1$ and $\tr(\bM) = n - k - 1$.
\end{enumerate}
\end{theorem}

\begin{proof}
\textbf{(i) Symmetry of $\bH$:} Since $\bX'\bX$ is symmetric, so is $(\bX'\bX)^{-1}$:
\[
\bH' = [\bX(\bX'\bX)^{-1}\bX']' = \bX[(\bX'\bX)^{-1}]'\bX' = \bX(\bX'\bX)^{-1}\bX' = \bH.
\]

\textbf{Idempotency of $\bH$:}
\[
\bH^2 = \bX(\bX'\bX)^{-1}\underbrace{\bX'\bX(\bX'\bX)^{-1}}_{\bI_{k+1}}\bX' = \bX(\bX'\bX)^{-1}\bX' = \bH.
\]

\textbf{Symmetry and idempotency of $\bM$:} $\bM' = (\bI - \bH)' = \bI - \bH' = \bI - \bH = \bM$. Also, $\bM^2 = (\bI - \bH)^2 = \bI - 2\bH + \bH^2 = \bI - 2\bH + \bH = \bI - \bH = \bM$.

\textbf{(ii)} $\bH\bX = \bX(\bX'\bX)^{-1}\bX'\bX = \bX\bI_{k+1} = \bX$, and $\bM\bX = (\bI - \bH)\bX = \bX - \bX = \bzero$.

\textbf{(iii)} $\bH\bM = \bH(\bI - \bH) = \bH - \bH^2 = \bH - \bH = \bzero$, and $\bM\bH = (\bI - \bH)\bH = \bH - \bH^2 = \bzero$.

\textbf{(iv)} By Lemma~\ref{lem:idempotent}, $\tr(\bH) = \rank(\bH)$. We show $\rank(\bH) = k + 1$ in two steps.

\emph{Step 1:} $\calC(\bH) = \calC(\bX)$. The inclusion $\calC(\bH) \subseteq \calC(\bX)$ is immediate: every column of $\bH = \bX(\bX'\bX)^{-1}\bX'$ is of the form $\bX\bv$ for some $\bv \in \Real^{k+1}$. Conversely, $\calC(\bX) \subseteq \calC(\bH)$ because $\bH\bX = \bX$ by~\ref{HM:ii}, so every column of $\bX$ equals $\bH$ times a column of $\bX$, hence lies in $\calC(\bH)$.

\emph{Step 2:} $\dim\calC(\bX) = k + 1$ by assumption~\ref{A4}.

Combining, $\rank(\bH) = \dim\calC(\bH) = \dim\calC(\bX) = k + 1$, so $\tr(\bH) = k + 1$. Finally, $\tr(\bM) = \tr(\bI) - \tr(\bH) = n - (k + 1) = n - k - 1$.
\end{proof}

\begin{remark}[$\bH$ is the orthogonal projection onto $\calC(\bX)$]\label{rem:hat-projection}
For a subspace $\mathcal{S} \subseteq \Real^n$, a matrix $\bP$ is called the \textbf{orthogonal projection onto $\mathcal{S}$} if $\bP\bv \in \mathcal{S}$ and $\bv - \bP\bv \perp \mathcal{S}$ for every $\bv \in \Real^n$. By Step~1 of the proof of Theorem~\ref{thm:ols-proj} (which used only that $\calC(\bX)$ is a subspace), $\bP\bv$ is then the unique point of $\mathcal{S}$ closest to $\bv$; hence there is at most one such matrix. One always exists: if $\mathcal{S} = \{\bzero\}$, the zero matrix is the orthogonal projection onto $\mathcal{S}$ (it maps every $\bv$ to $\bzero \in \mathcal{S}$, and $\bv - \bzero = \bv \perp \{\bzero\}$ since $\bv'\bzero = 0$), and it is symmetric and idempotent with column space $\{\bzero\} = \mathcal{S}$; otherwise, writing $\mathcal{S} = \calC(\mathbf{B})$ with the columns of $\mathbf{B}$ a basis of $\mathcal{S}$ (so that $\mathbf{B}$ has full column rank), Definition~\ref{def:hat-matrix} and Theorem~\ref{thm:HM} apply with $\mathbf{B}$ in place of $\bX$, and $\mathbf{B}(\mathbf{B}'\mathbf{B})^{-1}\mathbf{B}'$ is the orthogonal projection onto $\mathcal{S}$ (it maps $\bv$ into $\calC(\mathbf{B})$, and $\mathbf{B}'[\bv - \mathbf{B}(\mathbf{B}'\mathbf{B})^{-1}\mathbf{B}'\bv] = \bzero$); it is symmetric and idempotent with column space $\mathcal{S}$. In particular $\bH$ is the orthogonal projection onto $\calC(\bX)$: $\bH\bv \in \calC(\bX)$ and $\bX'(\bv - \bH\bv) = \bX'\bM\bv = \bzero$ by~\ref{HM:ii}.

Combining~\ref{HM:i}, \ref{HM:ii}, and the identity $\calC(\bH) = \calC(\bX)$ proved above: $\bH$ is symmetric, idempotent, and has column space $\calC(\bX)$. These three properties characterize the orthogonal projection onto $\calC(\bX)$, in the sense that any symmetric idempotent $\bP$ with $\calC(\bP) = \calC(\bX)$ equals $\bH$. Indeed, for $\bv \in \calC(\bX) = \calC(\bP)$, say $\bv = \bP\bw$, we get $\bP\bv = \bP^2\bw = \bP\bw = \bv$; for $\bv \perp \calC(\bX) = \calC(\bP)$, we get $\|\bP\bv\|^2 = \bv'\bP'\bP\bv = \bv'(\bP\bv) = 0$ because $\bP\bv \in \calC(\bP)$; and every $\bv \in \Real^n$ decomposes as $\bv = \bH\bv + \bM\bv$ with $\bH\bv \in \calC(\bX)$ and $\bM\bv \perp \calC(\bX)$, so $\bP\bv = \bP(\bH\bv) + \bP(\bM\bv) = \bH\bv + \bzero = \bH\bv$. Geometrically, $\hat{\bY} = \bH\bY$ is the closest point in $\calC(\bX)$ to $\bY$ in Euclidean distance (this is the content of \S\ref{sec:ols}, specifically the orthogonal projection derivation). Analogously, $\bM = \bI - \bH$ is the orthogonal projection onto $\calC(\bX)^\perp$: $\bM\bv \perp \calC(\bX)$ as just noted, and $\bv - \bM\bv = \bH\bv \in \calC(\bX)$ is orthogonal to $\calC(\bX)^\perp$.
\end{remark}

\subsection{Errors Versus Residuals}

\begin{remark}[Errors versus residuals]\label{rem:errors-residuals}
The \textbf{error} $\varepsilon_i = Y_i - \bx_i'\bb$ is the deviation from the true regression function; computing it requires the unknown $\bb$, so it is \emph{unobservable}. The \textbf{residual} $e_i = Y_i - \bx_i'\hat{\bb} = Y_i - \hat{Y}_i$ is the deviation from the fitted regression function and is a function of the data $(\bY, \bX)$, hence \emph{observable}. In matrix form, $\mathbf{e} = \bM\bY = \bM(\bX\bb + \be) = \bM\be$ (since $\bM\bX = \bzero$ by Theorem~\ref{thm:HM}\ref{HM:ii}). This identity does not make $\be$ observable: $\bM$ is the orthogonal projection onto $\calC(\bX)^\perp$, of rank $n-k-1$, so $\mathbf{e}$ retains only the component of $\be$ orthogonal to the column space, and the discarded component $\bH\be = \bX(\hat{\bb} - \bb)$ is precisely the estimation error. The residual is thus the observable part of the error, and the two behave differently: under assumption~\ref{A3} the $\varepsilon_i$ are uncorrelated with equal variance, whereas $\var(\mathbf{e} \mid \bX) = \sigma^2\bM$ (Theorem~\ref{thm:var-resid}) makes the $e_i$ correlated in general (with an intercept and $n > k+1$, $\bM\bone = \bzero$ and $\tr\bM = n-k-1 > 0$ force a nonzero off-diagonal entry of $\bM$) and, unless the diagonal entries of $\bM$ are all equal, heteroskedastic.
\end{remark}

\subsection{Algebraic Properties of OLS Residuals}

\begin{theorem}[Implications of the Normal Equations]\label{thm:normal-eq}
The OLS residuals satisfy:
\begin{enumerate}[label=\textup{(\roman*)}]
\item\label{ne:i} $\bX'\mathbf{e} = \bzero$ (residuals are orthogonal to every column of $\bX$).
\item\label{ne:ii} $\sum_{i=1}^{n} e_i = 0$ (residuals sum to zero), provided the model includes an intercept.
\item\label{ne:iii} $\hat{\bY}'\mathbf{e} = 0$ (residuals are orthogonal to fitted values).
\end{enumerate}
\end{theorem}

\begin{proof}
\textbf{(i)} From the normal equation $\bX'\bX\hat{\bb} = \bX'\bY$, rearranging:
\[
\bX'\bY - \bX'\bX\hat{\bb} = \bzero \iff \bX'(\bY - \bX\hat{\bb}) = \bzero \iff \bX'\mathbf{e} = \bzero.
\]

\textbf{(ii)} The first column of $\bX$ is $\bone$. The first component of the vector equation $\bX'\mathbf{e} = \bzero$ is $\bone'\mathbf{e} = \sum_{i=1}^{n} e_i = 0$. This relies on the presence of the intercept column.

\textbf{(iii)} Since $\hat{\bY} = \bX\hat{\bb}$:
\[
\hat{\bY}'\mathbf{e} = (\bX\hat{\bb})'\mathbf{e} = \hat{\bb}'\bX'\mathbf{e} = \hat{\bb}' \cdot \bzero = 0.
\]
Geometrically, $\hat{\bY} \in \calC(\bX)$ and $\mathbf{e} \perp \calC(\bX)$.
\end{proof}

\subsection{Leverages}

\begin{definition}[Leverage]\label{def:leverage}
The \textbf{leverage} of the $i$-th observation is $h_{ii}$, the $i$-th diagonal entry of $\bH$.
\end{definition}

\begin{theorem}[Properties of Leverages]\label{thm:leverage}
\leavevmode
\begin{enumerate}[label=\textup{(\roman*)}]
\item\label{lev:i} $0 \leqslant h_{ii} \leqslant 1$ for all $i$.
\item $\sum_{i=1}^{n} h_{ii} = k + 1$.
\end{enumerate}
\end{theorem}

\begin{proof}
\textbf{(i)} Since $\bH$ is symmetric and idempotent, $\bH = \bH^2 = \bH'\bH$. Comparing the $(i, i)$ entry on both sides using $h_{ji} = h_{ij}$:
\[
h_{ii} = \sum_{j=1}^{n} h_{ij} h_{ji} = \sum_{j=1}^{n} h_{ij}^2 = h_{ii}^2 + \sum_{j \neq i} h_{ij}^2.
\]
The rightmost sum is $\geqslant 0$, so $h_{ii} \geqslant h_{ii}^2$, which rearranges to $h_{ii}(1 - h_{ii}) \geqslant 0$. The leftmost expression $h_{ii} = \sum_j h_{ij}^2 \geqslant 0$ is also a sum of squares. Combining $h_{ii} \geqslant 0$ with $h_{ii}(1 - h_{ii}) \geqslant 0$ forces $h_{ii} \leqslant 1$.

\textbf{(ii)} $\sum_{i=1}^{n} h_{ii} = \tr(\bH) = k + 1$ by Theorem~\ref{thm:HM}\ref{HM:iv}.
\end{proof}

\subsection{The Case \texorpdfstring{$k = 1$}{k = 1}}\label{sec:k1}

When $k = 1$, the model of Definition~\ref{def:model} is the \textbf{simple linear regression model} $Y_i = \beta_0 + \beta_1 x_i + \varepsilon_i$, where we write $x_i = X_{i1}$ (a scalar, not to be confused with the row $\bx_i' = (1, x_i)$ of $\bX$) and $\bx = (x_1, \ldots, x_n)'$, so that $\bX = [\bone \;\; \bx]$. This subsection specializes the formulas above to this case. They are expressed through the following sample quantities.

\begin{definition}[Sample Moments and Sample Correlation]\label{def:corr}
For vectors $\bu, \bv \in \Real^n$ with means $\bar u = n^{-1}\sum_i u_i$ and $\bar v = n^{-1}\sum_i v_i$, write
\[
S_{uv} = \sum_{i=1}^n (u_i - \bar u)(v_i - \bar v) = \sum_{i=1}^n u_iv_i - n\bar u\bar v, \qquad S_{uu} = \sum_{i=1}^n (u_i - \bar u)^2 = \sum_{i=1}^n u_i^2 - n\bar u^2
\]
(the second forms follow on expanding the products and using $\sum_i u_i = n\bar u$ and $\sum_i v_i = n\bar v$), and, for $n \geqslant 2$, $s_u^2 = S_{uu}/(n-1)$ for the \textbf{sample variance} of $\bu$. Since $S_{uu}$ is a sum of squares, $S_{uu} = 0$ iff $\bu$ is \textbf{constant}, i.e.\ $\bu = \bar u\bone$. If neither $\bu$ nor $\bv$ is constant, the \textbf{(Pearson) sample correlation} of $\bu$ and $\bv$ is
\[
\corr(\bu, \bv) = \frac{S_{uv}}{\sqrt{S_{uu}S_{vv}}} = \frac{S_{uv}}{(n-1)s_us_v}.
\]
In the simple linear regression model we write $r = \corr(\bx, \bY)$, $S_{YY} = \sum_i(Y_i - \bar Y)^2$ and $s_Y^2 = S_{YY}/(n-1)$.
\end{definition}

\begin{lemma}[Cauchy--Schwarz Inequality]\label{lem:cs}
For $\bu, \bv \in \Real^n$, $(\bu'\bv)^2 \leqslant \|\bu\|^2\|\bv\|^2$, with equality if and only if one of the two vectors is a scalar multiple of the other.
\end{lemma}

\begin{proof}
If $\bu = \bzero$, both sides vanish and $\bu = 0 \cdot \bv$. Let $\bu \neq \bzero$, put $c = \bu'\bv/\|\bu\|^2$ and $\bw = \bv - c\bu$. Then $\bu'\bw = \bu'\bv - c\|\bu\|^2 = 0$, so Pythagoras applied to $\bv = c\bu + \bw$ gives
\[
\|\bv\|^2 = c^2\|\bu\|^2 + \|\bw\|^2 = \frac{(\bu'\bv)^2}{\|\bu\|^2} + \|\bw\|^2 \;\geqslant\; \frac{(\bu'\bv)^2}{\|\bu\|^2},
\]
which is the inequality, with equality iff $\bw = \bzero$, i.e.\ $\bv = c\bu$. Conversely, if $\bv = a\bu$ for some $a$, then $c = a$ and $\bw = \bzero$; and if $\bu = a\bv$ with $\bu \neq \bzero$, then $a \neq 0$ and $\bv = a^{-1}\bu$.
\end{proof}

\begin{proposition}[Properties of the Sample Correlation]\label{prop:corr}
Let $\bu, \bv \in \Real^n$ be non-constant. Then:
\begin{enumerate}[label=\textup{(\roman*)}]
\item $-1 \leqslant \corr(\bu, \bv) \leqslant 1$;
\item $\corr(\bu, \bv) = 1$ (respectively $-1$) iff there are constants $\alpha$ and $\beta > 0$ (respectively $\beta < 0$) with $v_i = \alpha + \beta u_i$ for all $i$, i.e.\ iff the points $(u_i, v_i)$ lie on a line of positive (respectively negative) slope;
\item for constants $a, c$ and $b, d \neq 0$, replacing $u_i$ by $a + bu_i$ and $v_i$ by $c + dv_i$ for all $i$ leaves $\corr(\bu, \bv)$ unchanged if $bd > 0$ and changes its sign if $bd < 0$.
\end{enumerate}
\end{proposition}

\begin{proof}
Put $\tilde\bu = \bu - \bar u\bone$ and $\tilde\bv = \bv - \bar v\bone$. Then $S_{uv} = \tilde\bu'\tilde\bv$, $S_{uu} = \|\tilde\bu\|^2 > 0$ and $S_{vv} = \|\tilde\bv\|^2 > 0$, so $\corr(\bu, \bv) = \tilde\bu'\tilde\bv/(\|\tilde\bu\|\,\|\tilde\bv\|)$.

\textbf{(i)} Lemma~\ref{lem:cs} applied to $\tilde\bu$ and $\tilde\bv$ gives $\corr(\bu, \bv)^2 \leqslant 1$.

\textbf{(ii)} By the equality case of Lemma~\ref{lem:cs}, $|\corr(\bu, \bv)| = 1$ iff one of $\tilde\bu, \tilde\bv$ is a multiple of the other; as both are nonzero, this means $\tilde\bv = \beta\tilde\bu$ with $\beta \neq 0$, i.e.\ $v_i = (\bar v - \beta\bar u) + \beta u_i$ for all $i$. In that case $\corr(\bu, \bv) = \beta\|\tilde\bu\|^2/(|\beta|\,\|\tilde\bu\|^2) = \beta/|\beta|$, which is $1$ for $\beta > 0$ and $-1$ for $\beta < 0$. Conversely, if $v_i = \alpha + \beta u_i$ for all $i$, then $\beta \neq 0$ (otherwise $\bv$ would be constant), $\bar v = \alpha + \beta\bar u$, and $\tilde\bv = \beta\tilde\bu$.

\textbf{(iii)} The centered transformed vectors are $b\tilde\bu$ and $d\tilde\bv$ (the constants $a$ and $c$ cancel on centering), so the numerator $S_{uv}$ is multiplied by $bd$ and the denominator $\sqrt{S_{uu}S_{vv}}$ by $|b|\,|d|$.
\end{proof}

\begin{corollary}[OLS in Simple Linear Regression]\label{cor:slr}
Let $k = 1$ and $\bX = [\bone \;\; \bx]$.
\begin{enumerate}[label=\textup{(\roman*)}]
\item \ref{A4} holds if and only if $\bx$ is not constant, i.e.\ $S_{xx} > 0$.
\item\label{slr:ii} Under~\ref{A4},
\[
(\bX'\bX)^{-1} = \frac{1}{nS_{xx}}\begin{pmatrix} \sum_i x_i^2 & -n\bar x\\ -n\bar x & n\end{pmatrix} = \begin{pmatrix} \frac{1}{n} + \frac{\bar x^2}{S_{xx}} & -\frac{\bar x}{S_{xx}}\\[2pt] -\frac{\bar x}{S_{xx}} & \frac{1}{S_{xx}}\end{pmatrix}.
\]
\item\label{slr:iii} $\hat\beta_1 = S_{xY}/S_{xx}$ and $\hat\beta_0 = \bar Y - \hat\beta_1\bar x$; if $\bY$ is not constant, $\hat\beta_1 = r\,s_Y/s_x$.
\item\label{slr:iv} The fitted line $y = \hat\beta_0 + \hat\beta_1 x$ passes through the point of means $(\bar x, \bar Y)$, and $\hat Y_i - \bar Y = \hat\beta_1(x_i - \bar x)$ for every $i$.
\end{enumerate}
\end{corollary}

\begin{proof}
\textbf{(i)} Since $\bone \neq \bzero$, the columns $\bone$ and $\bx$ are linearly dependent iff $\bx = c\bone$ for some $c$, i.e.\ iff $\bx$ is constant, i.e.\ (Definition~\ref{def:corr}) iff $S_{xx} = 0$.

\textbf{(ii)} $\bX'\bX = \begin{pmatrix} n & n\bar x\\ n\bar x & \sum_i x_i^2\end{pmatrix}$. Multiplying it by the displayed matrix and using $\sum_i x_i^2 - n\bar x^2 = S_{xx}$,
\[
\frac{1}{nS_{xx}}\begin{pmatrix} n\sum_i x_i^2 - n^2\bar x^2 & -n^2\bar x + n^2\bar x\\ n\bar x\sum_i x_i^2 - n\bar x\sum_i x_i^2 & -n^2\bar x^2 + n\sum_i x_i^2\end{pmatrix} = \frac{1}{nS_{xx}}\begin{pmatrix} nS_{xx} & 0\\ 0 & nS_{xx}\end{pmatrix} = \bI_2,
\]
so the displayed matrix is $(\bX'\bX)^{-1}$. The second form uses $\sum_i x_i^2/(nS_{xx}) = (S_{xx} + n\bar x^2)/(nS_{xx})$.

\textbf{(iii)} Here $\bX'\bY = (n\bar Y, \sum_i x_iY_i)'$. The first row of the normal equation~\eqref{eq:normal-eq} reads $n\hat\beta_0 + n\bar x\hat\beta_1 = n\bar Y$, i.e.\ $\hat\beta_0 = \bar Y - \hat\beta_1\bar x$. The second coordinate of $\hat\bb = (\bX'\bX)^{-1}\bX'\bY$ is, by (ii),
\[
\hat\beta_1 = \frac{-n\bar x \cdot n\bar Y + n\sum_i x_iY_i}{nS_{xx}} = \frac{\sum_i x_iY_i - n\bar x\bar Y}{S_{xx}} = \frac{S_{xY}}{S_{xx}}.
\]
If $\bY$ is not constant, $S_{xY} = r\sqrt{S_{xx}S_{YY}}$ by Definition~\ref{def:corr}, so $\hat\beta_1 = r\sqrt{S_{YY}/S_{xx}} = r\,s_Y/s_x$.

\textbf{(iv)} By (iii), $\hat\beta_0 + \hat\beta_1\bar x = \bar Y$, and $\hat Y_i - \bar Y = (\hat\beta_0 + \hat\beta_1x_i) - (\hat\beta_0 + \hat\beta_1\bar x) = \hat\beta_1(x_i - \bar x)$.
\end{proof}

\begin{proposition}[The Slope as a Weighted Sum of the Responses]\label{prop:weights}
Let $k = 1$ and $S_{xx} > 0$. Then $\hat\beta_1 = \sum_{i=1}^n w_iY_i$ with $w_i = (x_i - \bar x)/S_{xx}$, and the weights satisfy
\[
\sum_{i=1}^n w_i = 0, \qquad \sum_{i=1}^n w_ix_i = 1, \qquad \sum_{i=1}^n w_i^2 = \frac{1}{S_{xx}}.
\]
(For general $k$, $\hat\bb = \bC\bY$ with the constant matrix $\bC = (\bX'\bX)^{-1}\bX'$ of Theorem~\ref{thm:ols}, so every $\hat\beta_j$ is a fixed linear combination of the responses.)
\end{proposition}

\begin{proof}
$\sum_i w_i = \sum_i (x_i - \bar x)/S_{xx} = 0$. Hence $\sum_i w_iY_i = \sum_i w_i(Y_i - \bar Y) + \bar Y\sum_i w_i = S_{xY}/S_{xx} = \hat\beta_1$ by Corollary~\ref{cor:slr}\ref{slr:iii}. Next, $\sum_i (x_i - \bar x)x_i = \sum_i (x_i - \bar x)^2 + \bar x\sum_i(x_i - \bar x) = S_{xx}$, so $\sum_i w_ix_i = 1$. Finally, $\sum_i w_i^2 = \sum_i (x_i - \bar x)^2/S_{xx}^2 = 1/S_{xx}$.
\end{proof}

\begin{corollary}[Regression Through the Origin]\label{cor:origin}
Theorem~\ref{thm:ols} and its proof do not use the intercept column, so they apply to any design of full column rank. For the model $Y_i = \beta x_i + \varepsilon_i$ without intercept, i.e.\ the single-column design $\bX = \bx$, which has full column rank iff $\bx \neq \bzero$, the OLS estimator is
\[
\hat\beta = \frac{\sum_i x_iY_i}{\sum_i x_i^2}, \qquad \hat Y_i = \hat\beta x_i.
\]
\end{corollary}

\begin{proof}
For $\bX = \bx$, $\bX'\bX = \sum_i x_i^2$, which is positive iff $\bx \neq \bzero$, and $\bX'\bY = \sum_i x_iY_i$; Theorem~\ref{thm:ols} gives the ratio. Theorem~\ref{thm:normal-eq}\ref{ne:i} and~\ref{ne:iii} still hold for this design, but~\ref{ne:ii} does not (its proof reads $\bone'\mathbf{e} = 0$ off the intercept row), so in general $\sum_i e_i \neq 0$.
\end{proof}

%=============================================================================
\section{Finite-Sample Properties Without Normality}\label{sec:finite}
%=============================================================================

In this section, only assumptions~\ref{A1}--\ref{A4} are used. No distributional assumption on $\be$ is needed.

Before deriving the finite-sample properties of $\hat{\bb}$, we collect the notation for the sum-of-squares and mean-square quantities used throughout the remaining sections.

\begin{definition}[Sums of Squares and Mean Squares]\label{def:sumsq}
Given the fitted model $\hat{\bY} = \bX\hat{\bb}$, define:
\begin{itemize}
\item \textbf{Total sum of squares}: $\SST = \sum_{i=1}^n(Y_i - \bar{Y})^2$, where $\bar Y = n^{-1}\sum_i Y_i$.
\item \textbf{Regression sum of squares}: $\SSR = \sum_{i=1}^n(\hat{Y}_i - \bar{Y})^2$.
\item \textbf{Error (residual) sum of squares}: $\SSE = \sum_{i=1}^n(Y_i - \hat{Y}_i)^2 = \mathbf{e}'\mathbf{e}$, where $\mathbf{e} = \bY - \hat{\bY}$ is the residual vector of Definition~\ref{def:hat-matrix}.
\item \textbf{Mean squared error}: $\MSE = s^2 = \SSE/(n - k - 1)$, defined when $n > k + 1$.
\item \textbf{Mean squared regression}: $\MSR = \SSR/k$, defined when $k \geqslant 1$ (at least one non-intercept regressor).
\end{itemize}
The decomposition $\SST = \SSR + \SSE$ is proved in Theorem~\ref{thm:anova} (§\ref{sec:R2}); the unbiasedness $\expc[\MSE] = \sigma^2$ is proved in Theorem~\ref{thm:mse-unbiased} below.
\end{definition}

We also record a tool used repeatedly in this and later sections.

\begin{lemma}[Variance and covariance under linear transformations]\label{lem:var-linear}
Let $\bz$ be an $n \times 1$ random vector with mean $\boldsymbol\mu = \expc[\bz]$ and covariance matrix $\var(\bz) = \expc[(\bz - \boldsymbol\mu)(\bz - \boldsymbol\mu)']$. For two random vectors $\mathbf{U}$ ($m \times 1$) and $\mathbf{V}$ ($l \times 1$) with finite second moments, the \textbf{cross-covariance matrix} is the $m \times l$ matrix
\[
\cov(\mathbf{U}, \mathbf{V}) = \expc\bigl[(\mathbf{U} - \expc[\mathbf{U}])(\mathbf{V} - \expc[\mathbf{V}])'\bigr],
\]
so that $\var(\bz) = \cov(\bz, \bz)$. Let $\bA$ ($m \times n$) and $\mathbf{B}$ ($l \times n$) be constant matrices and $\mathbf{b} \in \Real^m$, $\mathbf{d} \in \Real^l$ constant vectors. Then
\begin{enumerate}[label=\textup{(\roman*)}]
\item\label{vl:i} $\expc[\bA\bz + \mathbf{b}] = \bA\boldsymbol\mu + \mathbf{b}$ and
\begin{equation}\label{eq:var-linear}
\var(\bA\bz + \mathbf{b}) = \bA\,\var(\bz)\,\bA';
\end{equation}
\item\label{vl:ii} $\cov(\bA\bz + \mathbf{b},\; \mathbf{B}\bz + \mathbf{d}) = \bA\,\var(\bz)\,\mathbf{B}'$.
\end{enumerate}
In particular, under assumption~\ref{A3}, $\var(\bY) = \sigma^2\bI_n$.
\end{lemma}

\begin{proof}
\textbf{(i)} Since $\bA$ and $\mathbf{b}$ are constant, $\expc[\bA\bz + \mathbf{b}] = \bA\boldsymbol\mu + \mathbf{b}$. Therefore
\begin{align*}
\var(\bA\bz + \mathbf{b}) &= \expc\bigl[(\bA\bz + \mathbf{b} - \bA\boldsymbol\mu - \mathbf{b})(\bA\bz + \mathbf{b} - \bA\boldsymbol\mu - \mathbf{b})'\bigr]\\
&= \expc\bigl[\bA(\bz - \boldsymbol\mu)(\bz - \boldsymbol\mu)'\bA'\bigr]\\
&= \bA\,\expc[(\bz - \boldsymbol\mu)(\bz - \boldsymbol\mu)']\,\bA' = \bA\,\var(\bz)\,\bA',
\end{align*}
since constant matrices pull outside the expectation.

\textbf{(ii)} By the first part of (i), $\bA\bz + \mathbf{b} - \expc[\bA\bz + \mathbf{b}] = \bA(\bz - \boldsymbol\mu)$ and $\mathbf{B}\bz + \mathbf{d} - \expc[\mathbf{B}\bz + \mathbf{d}] = \mathbf{B}(\bz - \boldsymbol\mu)$, so
\[
\cov(\bA\bz + \mathbf{b},\; \mathbf{B}\bz + \mathbf{d}) = \expc\bigl[\bA(\bz - \boldsymbol\mu)(\bz - \boldsymbol\mu)'\mathbf{B}'\bigr] = \bA\,\expc[(\bz - \boldsymbol\mu)(\bz - \boldsymbol\mu)']\,\mathbf{B}' = \bA\,\var(\bz)\,\mathbf{B}'.
\]
(Taking $\mathbf{B} = \bA$, $\mathbf{d} = \mathbf{b}$ recovers~\eqref{eq:var-linear}.)

Finally, applying (i) to $\bY = \bX\bb + \be$ (with $\bA = \bI_n$, $\mathbf{b} = \bX\bb$ constant given the conditioning on $\bX$) gives $\var(\bY) = \var(\be) = \sigma^2\bI_n$ by~\ref{A3}.
\end{proof}

\subsection{Unbiasedness}

\begin{theorem}[Unbiasedness of $\hat{\bb}$]\label{thm:unbiased}
Under~\ref{A2} and~\ref{A4} (the latter is needed for $\hat{\bb}$ to exist), the OLS estimator is unbiased:
\begin{equation}
\boxed{\expc[\hat{\bb}] = \bb.}
\end{equation}
\end{theorem}

\begin{proof}
Substituting $\bY = \bX\bb + \be$ into $\hat{\bb} = (\bX'\bX)^{-1}\bX'\bY$:
\begin{align}
\hat{\bb} &= (\bX'\bX)^{-1}\bX'(\bX\bb + \be) \notag\\
&= (\bX'\bX)^{-1}\bX'\bX\bb + (\bX'\bX)^{-1}\bX'\be \notag\\
&= \bI_{k+1}\bb + (\bX'\bX)^{-1}\bX'\be \notag\\
&= \bb + (\bX'\bX)^{-1}\bX'\be. \label{eq:beta-decomp}
\end{align}
Taking expectations, using~\ref{A2} ($\expc[\be] = \bzero$):
\[
\expc[\hat{\bb}] = \bb + (\bX'\bX)^{-1}\bX'\expc[\be] = \bb + (\bX'\bX)^{-1}\bX' \cdot \bzero = \bb. \qedhere
\]
\end{proof}

\subsection{Variance--Covariance Matrix}

\begin{theorem}[Variance--Covariance Matrix of $\hat{\bb}$]\label{thm:var-beta}
Under~\ref{A3} and~\ref{A4}:
\begin{equation}
\boxed{\var(\hat{\bb}) = \sigma^2 (\bX'\bX)^{-1}.}
\end{equation}
\end{theorem}

\begin{proof}
From~\eqref{eq:beta-decomp}, $\hat{\bb} = \bb + (\bX'\bX)^{-1}\bX'\be$, which is a linear transformation of $\be$ (with a constant shift). Applying Lemma~\ref{lem:var-linear} with $\bA = (\bX'\bX)^{-1}\bX'$:
\[
\var(\hat{\bb}) = (\bX'\bX)^{-1}\bX'\,\var(\be)\,\bX(\bX'\bX)^{-1}.
\]
By~\ref{A3}, $\var(\be) = \sigma^2\bI_n$. Substituting:
\[
\var(\hat{\bb}) = (\bX'\bX)^{-1}\bX'(\sigma^2\bI)\bX(\bX'\bX)^{-1} = \sigma^2(\bX'\bX)^{-1}\bX'\bX(\bX'\bX)^{-1} = \sigma^2(\bX'\bX)^{-1},
\]
using $\bX'\bX(\bX'\bX)^{-1} = \bI_{k+1}$. The diagonal entry $[\var(\hat{\bb})]_{jj} = \sigma^2[(\bX'\bX)^{-1}]_{jj}$ gives $\var(\hat{\beta}_j)$.
\end{proof}

\subsection{The Gauss--Markov Theorem}

\begin{theorem}[Gauss--Markov Theorem]\label{thm:gm}
Under~\ref{A1}--\ref{A4}, the OLS estimator $\hat{\bb}$ is the \textbf{Best Linear Unbiased Estimator} (BLUE) of $\bb$. That is, for any other linear unbiased estimator $\tilde{\bb} = \bC\bY$ (where $\bC$ is a $(k+1) \times n$ constant matrix), the difference $\var(\tilde{\bb}) - \var(\hat{\bb})$ is positive semi-definite. In particular, $\var(\tilde{\beta}_j) \geqslant \var(\hat{\beta}_j)$ for every $j$.
\end{theorem}

\begin{proof}
\textbf{Step 1: Necessary condition for unbiasedness.} By~\ref{A1}, $\expc[\tilde{\bb}] = \bC\bX\bb = \bb$ for all $\bb \in \Real^{k+1}$ forces $\bC\bX = \bI_{k+1}$ (setting $\bv = \boldsymbol\iota_j$ in $\bC\bX\bv = \bv$ reads off the $j$-th column; here and below $\boldsymbol\iota_j$ denotes the $j$-th standard basis vector, i.e.\ the $j$-th column of the identity matrix of the relevant dimension).

\textbf{Step 2: Decompose $\bC$.} Write $\bC = (\bX'\bX)^{-1}\bX' + \bD$, where $\bD$ is the excess beyond the OLS coefficient matrix. From $\bC\bX = \bI_{k+1}$:
\[
\bD\bX = \bC\bX - (\bX'\bX)^{-1}\bX'\bX = \bI_{k+1} - \bI_{k+1} = \bzero.
\]

\textbf{Step 3: Compute $\var(\tilde{\bb})$.} Since $\tilde{\bb} = \bC\bY$ with $\bC$ constant, Lemma~\ref{lem:var-linear} gives
\begin{align*}
\var(\tilde{\bb}) = \bC\var(\bY)\bC' = \sigma^2\bC\bC' &= \sigma^2\bigl[(\bX'\bX)^{-1}\bX' + \bD\bigr]\bigl[\bX(\bX'\bX)^{-1} + \bD'\bigr].
\end{align*}
Expanding yields four terms:
\begin{align*}
\sigma^2\bC\bC' ={}& \sigma^2(\bX'\bX)^{-1}\bX'\bX(\bX'\bX)^{-1} && \text{(term A)}\\
& + \sigma^2(\bX'\bX)^{-1}\bX'\bD' && \text{(term B)}\\
& + \sigma^2\bD\bX(\bX'\bX)^{-1} && \text{(term C)}\\
& + \sigma^2\bD\bD'. && \text{(term D)}
\end{align*}
Term A simplifies to $\sigma^2(\bX'\bX)^{-1}$ because $\bX'\bX(\bX'\bX)^{-1} = \bI_{k+1}$. Term C vanishes because $\bD\bX = \bzero$ (Step~2). Term B equals $(\text{term C})'$, so it also vanishes. Hence:
\[
\var(\tilde{\bb}) = \sigma^2(\bX'\bX)^{-1} + \sigma^2\bD\bD' = \var(\hat{\bb}) + \sigma^2\bD\bD'.
\]
Since $\bD\bD'$ is positive semi-definite (for any $\bv \in \Real^{k+1}$, $\bv'\bD\bD'\bv = \|\bD'\bv\|^2 \geqslant 0$):
\[
\var(\tilde{\bb}) - \var(\hat{\bb}) = \sigma^2\bD\bD' \succeq \bzero.
\]
Equality holds if and only if $\bD'\bv = \bzero$ for all $\bv$, i.e., $\bD = \bzero$; equivalently, $\bC = (\bX'\bX)^{-1}\bX'$ and $\tilde{\bb} = \hat{\bb}$.
\end{proof}

\subsection{Estimation of \texorpdfstring{$\sigma^2$}{sigma\texttwosuperior}}

\begin{lemma}[Expected Value of a Quadratic Form]\label{lem:quadratic-form}
Let $\bz$ be an $n \times 1$ random vector with $\expc[\bz] = \boldsymbol\mu$ and $\var(\bz) = \boldsymbol\Sigma$. Let $\bA$ be an $n \times n$ constant matrix. Then:
\begin{equation}\label{eq:quad-expect}
\expc[\bz'\bA\bz] = \tr(\bA\boldsymbol\Sigma) + \boldsymbol\mu'\bA\boldsymbol\mu.
\end{equation}
In particular, if $\expc[\bz] = \bzero$, then $\expc[\bz'\bA\bz] = \tr(\bA\boldsymbol\Sigma)$.
\end{lemma}

\begin{proof}
The quadratic form is a scalar: $\bz'\bA\bz = \sum_{i,j} a_{ij} z_i z_j$. Taking expectations:
\[
\expc[\bz'\bA\bz] = \sum_{i,j} a_{ij} \expc[z_i z_j].
\]
From $\cov(z_i, z_j) = \expc[z_i z_j] - \mu_i\mu_j$, we get $\expc[z_i z_j] = \Sigma_{ij} + \mu_i\mu_j$. Substituting:
\begin{align*}
\expc[\bz'\bA\bz] &= \sum_{i,j} a_{ij}\Sigma_{ij} + \sum_{i,j} a_{ij}\mu_i\mu_j.
\end{align*}
The first sum equals $\tr(\bA\boldsymbol\Sigma)$: indeed, $\tr(\bA\boldsymbol\Sigma) = \sum_i (\bA\boldsymbol\Sigma)_{ii} = \sum_i \sum_j a_{ij}\Sigma_{ji} = \sum_{i,j} a_{ij}\Sigma_{ij}$, where the last step uses the symmetry $\Sigma_{ji} = \Sigma_{ij}$. The second sum equals $\boldsymbol\mu'\bA\boldsymbol\mu$: indeed, $\boldsymbol\mu'\bA\boldsymbol\mu = \sum_i \mu_i (\bA\boldsymbol\mu)_i = \sum_i \mu_i \sum_j a_{ij}\mu_j = \sum_{i,j} a_{ij}\mu_i\mu_j$. Combining yields~\eqref{eq:quad-expect}.
\end{proof}

\begin{theorem}[Unbiasedness of MSE]\label{thm:mse-unbiased}
Under~\ref{A2}--\ref{A4}, the \textbf{mean squared error}
\begin{equation}
s^2 = \MSE = \frac{\SSE}{n - k - 1} = \frac{\mathbf{e}'\mathbf{e}}{n - k - 1}
\end{equation}
is an unbiased estimator of $\sigma^2$: $\expc[\MSE] = \sigma^2$.
\end{theorem}

\begin{proof}
\textbf{Step 1: Express $\SSE$ as a quadratic form in $\be$.} Since $\mathbf{e} = \bM\be$ (Remark~\ref{rem:errors-residuals}):
\[
\SSE = \mathbf{e}'\mathbf{e} = (\bM\be)'(\bM\be) = \be'\bM'\bM\be = \be'\bM\be,
\]
where the last step uses $\bM'\bM = \bM\bM = \bM$ (symmetric idempotent).

\textbf{Step 2: Apply Lemma~\ref{lem:quadratic-form}.} With $\bz = \be$, $\bA = \bM$, $\boldsymbol\mu = \expc[\be] = \bzero$ (by~\ref{A2}), and $\boldsymbol\Sigma = \var(\be) = \sigma^2\bI$ (by~\ref{A3}):
\[
\expc[\SSE] = \expc[\be'\bM\be] = \tr(\bM \cdot \sigma^2\bI) + \bzero'\bM\bzero = \sigma^2\tr(\bM).
\]

\textbf{Step 3: Evaluate the trace.} By Theorem~\ref{thm:HM}\ref{HM:iv}, $\tr(\bM) = n - k - 1$. Hence $\expc[\SSE] = \sigma^2(n - k - 1)$.

\textbf{Step 4: Conclude.}
\[
\expc[\MSE] = \expc\left[\frac{\SSE}{n-k-1}\right] = \frac{\sigma^2(n-k-1)}{n-k-1} = \sigma^2. \qedhere
\]
\end{proof}

\subsection{Variance of Individual Residuals}

\begin{theorem}[Variance of Residuals]\label{thm:var-resid}
Under~\ref{A3} and~\ref{A4}:
\begin{equation}
\var(\mathbf{e}) = \sigma^2\bM, \qquad \text{and in particular, } \var(e_i) = \sigma^2(1 - h_{ii}).
\end{equation}
\end{theorem}

\begin{proof}
Since $\mathbf{e} = \bM\be$ and $\bM$ is symmetric ($\bM' = \bM$) and idempotent ($\bM^2 = \bM$):
\[
\var(\mathbf{e}) = \bM\var(\be)\bM' = \bM(\sigma^2\bI)\bM' = \sigma^2\bM\bM = \sigma^2\bM^2 = \sigma^2\bM.
\]
The $(i, i)$ entry gives $\var(e_i) = \sigma^2 m_{ii}$. Since $\bM = \bI - \bH$, we have $m_{ii} = 1 - h_{ii}$, so $\var(e_i) = \sigma^2(1 - h_{ii})$.
\end{proof}

\begin{remark}\label{rem:var-resid-bound}
By Theorem~\ref{thm:leverage}\ref{lev:i}, $0 \leqslant h_{ii} \leqslant 1$, so $\var(e_i) \leqslant \sigma^2$. High-leverage observations (large $h_{ii}$) produce residuals with smaller variance: the fitted line is pulled toward them, leaving less room for the residual to vary.
\end{remark}

\begin{corollary}[Variance of the Fitted Values]\label{cor:var-fitted}
Under~\ref{A3} and~\ref{A4}, the fitted vector has $\var(\hat{\bY}) = \sigma^2\bH$, and $\cov(\hat{\bY}, \mathbf{e}) = \bzero$. In particular $\var(\hat{Y}_i) = \sigma^2h_{ii}$ and
\[
\var(\hat{Y}_i) + \var(e_i) = \sigma^2 \qquad \text{for every } i.
\]
\end{corollary}

\begin{proof}
Since $\bH\bX = \bX$ (Theorem~\ref{thm:HM}\ref{HM:ii}), $\hat{\bY} = \bH\bY = \bX\bb + \bH\be$, and $\mathbf{e} = \bM\be$ (Remark~\ref{rem:errors-residuals}). By Lemma~\ref{lem:var-linear}\ref{vl:i} and~\ref{A3}, $\var(\hat{\bY}) = \bH(\sigma^2\bI)\bH' = \sigma^2\bH^2 = \sigma^2\bH$ ($\bH$ symmetric idempotent); by Lemma~\ref{lem:var-linear}\ref{vl:ii}, $\cov(\hat{\bY}, \mathbf{e}) = \bH(\sigma^2\bI)\bM' = \sigma^2\bH\bM = \bzero$ (Theorem~\ref{thm:HM}\ref{HM:iii}). The $(i,i)$ entry gives $\var(\hat{Y}_i) = \sigma^2h_{ii}$, and adding $\var(e_i) = \sigma^2(1 - h_{ii})$ (Theorem~\ref{thm:var-resid}) gives $\sigma^2$. Thus $h_{ii}$ is the share of $\sigma^2$ carried by the fitted value of observation $i$, and $1 - h_{ii}$ the share left in its residual.
\end{proof}

\subsection{Standardized Residuals and Unusual Observations}

\begin{definition}[Standardized Residuals, Outliers, and High-Leverage Points]\label{def:std-resid}
Assume $n > k + 1$ and $s > 0$. For an observation with $h_{ii} < 1$, the \textbf{standardized residual} is
\[
\tilde e_i = \frac{e_i}{s\sqrt{1 - h_{ii}}},
\]
the residual divided by the estimate $s\sqrt{1 - h_{ii}}$ of its standard deviation $\sigma\sqrt{1 - h_{ii}}$ (Theorems~\ref{thm:mse-unbiased} and~\ref{thm:var-resid}). By the usual conventions, observation $i$ is flagged as an \textbf{outlier} (unusual in the direction of the response) if $|\tilde e_i| > 2$, and as a \textbf{high-leverage point} (unusual in the regressors) if $h_{ii} > 2(k+1)/n$. An observation can be both.
\end{definition}

\begin{remark}[Rationale for the conventions]\label{rem:std-resid}
(a)~Under~\ref{A2} and~\ref{A3}, $e_i/(\sigma\sqrt{1 - h_{ii}})$ has mean $0$ (as $\expc[\mathbf{e}] = \bM\expc[\be] = \bzero$) and variance $1$ (Theorem~\ref{thm:var-resid}) whatever the leverage, whereas the raw residuals have the unequal variances $\sigma^2(1 - h_{ii})$; $\tilde e_i$ replaces $\sigma$ by $s$. Standardized residuals are dimensionless, so a single cutoff can be used for every observation and every data set. (b)~By Theorem~\ref{thm:leverage}, the average leverage is $(k+1)/n$, and the second rule flags leverages above twice the average. Since the leverages are non-negative and sum to $k + 1$, fewer than $n/2$ observations can be flagged: if $m \geqslant 1$ of them exceed $2(k+1)/n$, then $k + 1 = \sum_i h_{ii} > 2m(k+1)/n$, i.e.\ $m < n/2$. (c)~The restriction $h_{ii} < 1$ is needed. If $h_{ii} = 1$, the identity $h_{ii} = h_{ii}^2 + \sum_{j \neq i}h_{ij}^2$ from the proof of Theorem~\ref{thm:leverage}\ref{lev:i} forces $h_{ij} = 0$ for all $j \neq i$, so $e_i = Y_i - \sum_j h_{ij}Y_j = Y_i - Y_i = 0$ for every $\bY$: the fit passes through observation $i$ whatever its response, and its residual carries no information about it.
\end{remark}

\subsection{The Case \texorpdfstring{$k = 1$}{k = 1}}

\begin{corollary}[Variances in Simple Linear Regression]\label{cor:slr-var}
In the simple linear regression model of \S\ref{sec:k1}, under~\ref{A3} and~\ref{A4},
\[
\var(\hat\beta_1) = \frac{\sigma^2}{S_{xx}}, \qquad \var(\hat\beta_0) = \sigma^2\Bigl(\frac{1}{n} + \frac{\bar x^2}{S_{xx}}\Bigr), \qquad \cov(\hat\beta_0, \hat\beta_1) = -\frac{\sigma^2\bar x}{S_{xx}},
\]
and for every $x_0 \in \Real$, with $\bx_0 = (1, x_0)'$,
\[
\var(\hat\beta_0 + \hat\beta_1x_0) = \sigma^2\bx_0'(\bX'\bX)^{-1}\bx_0 = \sigma^2\Bigl(\frac{1}{n} + \frac{(x_0 - \bar x)^2}{S_{xx}}\Bigr).
\]
\end{corollary}

\begin{proof}
By Theorem~\ref{thm:var-beta}, $\var(\hat{\bb}) = \sigma^2(\bX'\bX)^{-1}$, whose entries are read off Corollary~\ref{cor:slr}\ref{slr:ii}. For the last statement, $\hat\beta_0 + \hat\beta_1x_0 = \bx_0'\hat{\bb}$, so Lemma~\ref{lem:var-linear}\ref{vl:i} gives
\[
\var(\bx_0'\hat{\bb}) = \sigma^2\bx_0'(\bX'\bX)^{-1}\bx_0 = \sigma^2\Bigl(\frac{1}{n} + \frac{\bar x^2 - 2x_0\bar x + x_0^2}{S_{xx}}\Bigr) = \sigma^2\Bigl(\frac{1}{n} + \frac{(x_0 - \bar x)^2}{S_{xx}}\Bigr).
\]
(The slope's variance also follows from Proposition~\ref{prop:weights}: $\hat\beta_1 = \sum_i w_iY_i$ with $\var(\bY) = \sigma^2\bI$, so $\var(\hat\beta_1) = \sigma^2\sum_i w_i^2 = \sigma^2/S_{xx}$ by Lemma~\ref{lem:var-linear}.)
\end{proof}

%=============================================================================
\section{Multivariate Normal Tools}\label{sec:mvn}
%=============================================================================

This section collects the distributional results needed for exact inference in Section~\ref{sec:exact}.

\subsection{Prerequisites: the Multivariate Normal Distribution}\label{sec:prereq-prob}

For a random vector $\bz \in \Real^n$, the \textbf{moment-generating function} (MGF) is $M_{\bz}(\mathbf{t}) = \expc[\exp(\mathbf{t}'\bz)]$, $\mathbf{t} \in \Real^n$, whenever the expectation is finite.

\begin{definition}[Multivariate normal distribution]\label{def:mvn}
Let $\boldsymbol\mu \in \Real^n$ and let $\boldsymbol\Sigma$ be a symmetric positive semi-definite $n \times n$ matrix. A random vector $\bz$ has the \textbf{multivariate normal distribution} $N(\boldsymbol\mu, \boldsymbol\Sigma)$ if
\begin{equation}\label{eq:mvn-mgf}
M_{\bz}(\mathbf{t}) = \exp\bigl(\mathbf{t}'\boldsymbol\mu + \tfrac{1}{2}\mathbf{t}'\boldsymbol\Sigma\mathbf{t}\bigr) \quad \text{for all } \mathbf{t} \in \Real^n.
\end{equation}
Its mean vector is $\boldsymbol\mu$ and its covariance matrix is $\boldsymbol\Sigma$: the gradient of the right-hand side of~\eqref{eq:mvn-mgf} is $(\boldsymbol\mu + \boldsymbol\Sigma\mathbf{t})\,M_{\bz}(\mathbf{t})$ and its Hessian is $\bigl[\boldsymbol\Sigma + (\boldsymbol\mu + \boldsymbol\Sigma\mathbf{t})(\boldsymbol\mu + \boldsymbol\Sigma\mathbf{t})'\bigr]M_{\bz}(\mathbf{t})$, so evaluating at $\mathbf{t} = \bzero$ and using fact~\ref{N6} below gives $\expc[\bz] = \boldsymbol\mu$ and $\expc[\bz\bz'] = \boldsymbol\Sigma + \boldsymbol\mu\boldsymbol\mu'$, whence $\var(\bz) = \expc[\bz\bz'] - \boldsymbol\mu\boldsymbol\mu' = \boldsymbol\Sigma$. For $n = 1$ this is the usual $N(\mu, \sigma^2)$, $\sigma^2 \geqslant 0$.
\end{definition}

The definition is stated through the MGF rather than through a density because $\boldsymbol\Sigma$ is allowed to be singular: several normal vectors below, such as the residual vector $\mathbf{e} = \bM\be$ with $\var(\mathbf{e}) = \sigma^2\bM$ (Theorems~\ref{thm:var-resid} and~\ref{thm:sse-dist}), have singular covariance matrices and no density. The following facts about MGFs and distributions are used without proof; together with \ref{L1}--\ref{L4} they are the only results imported from outside these notes.
\begin{enumerate}[label=\textup{(N\arabic*)}]
\item\label{N1} \emph{Uniqueness.} If two random vectors have MGFs that are finite and equal on an open neighborhood of $\bzero$, they have the same distribution. (Used in Theorem~\ref{thm:mvn-linear}.)
\item\label{N2} \emph{Independence and factorization.} Random vectors $\mathbf{U}_1, \ldots, \mathbf{U}_m$ (with MGFs finite near $\bzero$) are mutually independent if and only if their joint MGF factorizes: $M_{(\mathbf{U}_1, \ldots, \mathbf{U}_m)}(\mathbf{t}_1, \ldots, \mathbf{t}_m) = \prod_{i=1}^m M_{\mathbf{U}_i}(\mathbf{t}_i)$ near $\bzero$. (Used in Theorems~\ref{thm:mvn-uncorr-indep}, \ref{thm:quadratic-chisq} and~\ref{thm:pred}.)
\item\label{N3} \emph{Functions of independent vectors.} If $\mathbf{U}$ and $\mathbf{V}$ are independent (written $\mathbf{U} \indep \mathbf{V}$; the symbol $\perp$ is reserved for orthogonality of vectors and subspaces), then $f(\mathbf{U})$ and $g(\mathbf{V})$ are independent for any (measurable) functions $f$, $g$. (Used whenever ``a function of $\hat{\bb}$ is independent of $\SSE$'' is asserted, e.g.\ in Theorems~\ref{thm:cochran}, \ref{thm:sse-dist}, \ref{thm:t-test}, \ref{thm:general-F}, \ref{thm:ci} and~\ref{thm:pred}.)
\item\label{N4} \emph{Density in the nonsingular case.} If $\boldsymbol\Sigma$ is positive definite, $N(\boldsymbol\mu, \boldsymbol\Sigma)$ has the density
\[
f(\bz) = (2\pi)^{-n/2}\,|\boldsymbol\Sigma|^{-1/2}\exp\bigl(-\tfrac{1}{2}(\bz - \boldsymbol\mu)'\boldsymbol\Sigma^{-1}(\bz - \boldsymbol\mu)\bigr)
\]
on $\Real^n$. (Used in Theorem~\ref{thm:mle}.)
\item\label{N5} \emph{Distributions derived from the normal.} The \textbf{chi-square distribution} $\chi^2_r$ ($r \geqslant 1$) is the distribution of $\sum_{i=1}^r W_i^2$ for i.i.d.\ $W_i \sim N(0,1)$. If $Z \sim N(0,1)$ and $V \sim \chi^2_\nu$ are independent, $Z/\sqrt{V/\nu}$ has the \textbf{$t$ distribution} $t_\nu$. If $U \sim \chi^2_a$ and $V \sim \chi^2_b$ are independent, $(U/a)/(V/b)$ has the \textbf{$F$ distribution} $F_{a,b}$. For $\alpha \in (0,1)$, the \textbf{upper $\alpha$ quantiles} $t_{\alpha,\nu}$ and $F_{\alpha,a,b}$ are the numbers with $\prb(T > t_{\alpha,\nu}) = \alpha$ for $T \sim t_\nu$ and $\prb(F > F_{\alpha,a,b}) = \alpha$ for $F \sim F_{a,b}$; all three have continuous distribution functions that are strictly increasing on their supports (the support of $t_\nu$ is $\Real$, and that of $\chi^2_r$ and of $F_{a,b}$ is $[0, \infty)$), so these quantiles exist and are unique. (Used from Theorem~\ref{thm:quadratic-chisq} onward; the $t$ and $F$ quantiles enter in Theorems~\ref{thm:ci}, \ref{thm:ellipsoid}, \ref{thm:ci-mean} and~\ref{thm:pred}.)
\item\label{N6} \emph{Moments from the MGF.} If $M_{\bz}$ is finite on an open neighborhood of $\bzero$, then its first- and second-order partial derivatives exist there and, at $\mathbf{t} = \bzero$, $\partial M_{\bz}/\partial t_i = \expc[z_i]$ and $\partial^2 M_{\bz}/\partial t_i\,\partial t_j = \expc[z_iz_j]$; that is, the gradient of $M_{\bz}$ at $\bzero$ is $\expc[\bz]$ and its Hessian (the matrix of second partial derivatives) is $\expc[\bz\bz']$. (Used in Definition~\ref{def:mvn}.)
\item\label{N7} \emph{Independence implies uncorrelatedness.} If $\mathbf{U}$ and $\mathbf{V}$ are independent random vectors with finite second moments, then $\expc[\mathbf{U}\mathbf{V}'] = \expc[\mathbf{U}]\,\expc[\mathbf{V}]'$, i.e.\ $\cov(\mathbf{U}, \mathbf{V}) = \bzero$ with $\cov$ as defined in Lemma~\ref{lem:var-linear}. Consequently, for independent scalar random variables $U$ and $V$, $\var(U \pm V) = \var(U) + \var(V)$: apply Lemma~\ref{lem:var-linear}\ref{vl:i} with $\bz = (U, V)'$ and $\bA = (1, \pm 1)$, the off-diagonal entries of $\var(\bz)$ being $\cov(U, V) = 0$. (Used in Theorems~\ref{thm:mvn-uncorr-indep} and~\ref{thm:pred} and in the solution to Exercise~\ref{ex:B6}.)
\end{enumerate}

The following consequence of \ref{N3} and \ref{N5} is used in Theorems~\ref{thm:ci} and~\ref{thm:pred}: \emph{the $t$ distribution is symmetric about $0$}. Let $T = Z/\sqrt{V/\nu}$ with $Z \sim N(0,1)$ and $V \sim \chi^2_\nu$ independent, as in \ref{N5}. Then $M_{-Z}(t) = \expc[e^{t(-Z)}] = M_Z(-t) = e^{t^2/2} = M_Z(t)$ for all $t$, so $-Z \sim N(0,1)$ by Definition~\ref{def:mvn}; $-Z$ is still independent of $V$ (\ref{N3}); hence $-T = (-Z)/\sqrt{V/\nu} \sim t_\nu$ by \ref{N5}, and $\prb(T < -c) = \prb(-T > c) = \prb(T > c)$ for every $c \in \Real$. Taking $c = t_{\alpha,\nu}$ gives the first identity below; the second follows from the first with $\alpha/2$ in place of $\alpha$, because $\prb(-t_{\alpha/2,\nu} \leqslant T \leqslant t_{\alpha/2,\nu}) = 1 - \prb(T < -t_{\alpha/2,\nu}) - \prb(T > t_{\alpha/2,\nu})$:
\begin{equation}\label{eq:t-symm}
\prb(T < -t_{\alpha,\nu}) = \alpha \qquad\text{and}\qquad \prb(-t_{\alpha/2,\nu} \leqslant T \leqslant t_{\alpha/2,\nu}) = 1 - \alpha \qquad (T \sim t_\nu,\; \alpha \in (0,1)).
\end{equation}
A second consequence, used in Proposition~\ref{prop:t-pvalue} and Corollary~\ref{cor:drop-one}: \emph{the square of a $t_\nu$ variable has the $F_{1,\nu}$ distribution}. With $T = Z/\sqrt{V/\nu}$ as above, $T^2 = (Z^2/1)/(V/\nu)$, where $Z^2 \sim \chi^2_1$ by the definition in~\ref{N5} (with $r = 1$) and $Z^2$ is independent of $V$ by~\ref{N3}; hence $T^2 \sim F_{1,\nu}$ by~\ref{N5}. In particular $\prb(T^2 > t_{\alpha/2,\nu}^2) = \prb(|T| > t_{\alpha/2,\nu}) = \alpha$ by~\eqref{eq:t-symm}, so $t_{\alpha/2,\nu}^2 = F_{\alpha,1,\nu}$.

\subsection{Linear Transformations, Independence, and Quadratic Forms}

\begin{theorem}[Closure Under Linear Transformations]\label{thm:mvn-linear}
If $\bz \sim N(\boldsymbol\mu, \boldsymbol\Sigma)$ is an $n \times 1$ normal random vector, $\bA$ is an $m \times n$ constant matrix, and $\mathbf{b} \in \Real^m$ is a constant vector, then
\begin{equation}
\bA\bz + \mathbf{b} \sim N(\bA\boldsymbol\mu + \mathbf{b},\; \bA\boldsymbol\Sigma\bA').
\end{equation}
\end{theorem}

\begin{proof}
By Definition~\ref{def:mvn}, the MGF of $\bz$ is $M_{\bz}(\mathbf{t}) = \exp(\mathbf{t}'\boldsymbol\mu + \frac{1}{2}\mathbf{t}'\boldsymbol\Sigma\mathbf{t})$. The MGF of $\bw = \bA\bz + \mathbf{b}$ is:
\begin{align*}
M_{\bw}(\mathbf{t}) &= \expc\!\left[e^{\mathbf{t}'(\bA\bz + \mathbf{b})}\right] = e^{\mathbf{t}'\mathbf{b}} \cdot \expc\!\left[e^{(\bA'\mathbf{t})' \bz}\right] = e^{\mathbf{t}'\mathbf{b}} \cdot M_{\bz}(\bA'\mathbf{t})\\
&= e^{\mathbf{t}'\mathbf{b}} \cdot \exp\bigl((\bA'\mathbf{t})'\boldsymbol\mu + \tfrac{1}{2}(\bA'\mathbf{t})'\boldsymbol\Sigma(\bA'\mathbf{t})\bigr)\\
&= \exp\bigl(\mathbf{t}'(\bA\boldsymbol\mu + \mathbf{b}) + \tfrac{1}{2}\mathbf{t}'\bA\boldsymbol\Sigma\bA'\mathbf{t}\bigr).
\end{align*}
This is the MGF~\eqref{eq:mvn-mgf} of $N(\bA\boldsymbol\mu + \mathbf{b},\; \bA\boldsymbol\Sigma\bA')$ (note that $\bA\boldsymbol\Sigma\bA'$ is symmetric positive semi-definite). By uniqueness of the MGF, fact~\ref{N1}, the result follows.
\end{proof}

\begin{theorem}[Uncorrelatedness $\Leftrightarrow$ Independence Under Joint Normality]\label{thm:mvn-uncorr-indep}
Let $(\mathbf{U}', \mathbf{V}')'$ be jointly normal:
\[
\begin{pmatrix} \mathbf{U} \\ \mathbf{V} \end{pmatrix} \sim N\!\left(\begin{pmatrix} \boldsymbol\mu_U \\ \boldsymbol\mu_V \end{pmatrix}, \begin{pmatrix} \boldsymbol\Sigma_{UU} & \boldsymbol\Sigma_{UV} \\ \boldsymbol\Sigma_{VU} & \boldsymbol\Sigma_{VV} \end{pmatrix}\right).
\]
Then $\boldsymbol\Sigma_{UV} = \bzero$ if and only if $\mathbf{U}$ and $\mathbf{V}$ are independent.
\end{theorem}

\begin{proof}
$(\Leftarrow)$: Independence implies uncorrelatedness for any distribution (fact~\ref{N7}), and $\boldsymbol\Sigma_{UV} = \cov(\mathbf{U}, \mathbf{V})$ because $\boldsymbol\Sigma$ is the covariance matrix of $(\mathbf{U}', \mathbf{V}')'$ (Definition~\ref{def:mvn}); hence $\boldsymbol\Sigma_{UV} = \bzero$.

$(\Rightarrow)$: With $\boldsymbol\Sigma_{UV} = \bzero$, the covariance is block-diagonal. The joint MGF is
\[
M_{(\mathbf{U},\mathbf{V})}(\mathbf{s}, \mathbf{t}) = \exp\bigl((\mathbf{s}', \mathbf{t}')\boldsymbol\mu + \tfrac{1}{2}(\mathbf{s}', \mathbf{t}')\boldsymbol\Sigma(\mathbf{s}', \mathbf{t}')'\bigr).
\]
Block expansion gives
\[
(\mathbf{s}', \mathbf{t}')\boldsymbol\mu = \mathbf{s}'\boldsymbol\mu_U + \mathbf{t}'\boldsymbol\mu_V, \qquad (\mathbf{s}', \mathbf{t}')\boldsymbol\Sigma\begin{pmatrix}\mathbf{s}\\\mathbf{t}\end{pmatrix} = \mathbf{s}'\boldsymbol\Sigma_{UU}\mathbf{s} + \mathbf{t}'\boldsymbol\Sigma_{VV}\mathbf{t},
\]
using $\boldsymbol\Sigma_{UV} = \boldsymbol\Sigma_{VU}' = \bzero$. The joint MGF factors:
\[
M_{(\mathbf{U},\mathbf{V})}(\mathbf{s}, \mathbf{t}) = \exp\bigl(\mathbf{s}'\boldsymbol\mu_U + \tfrac{1}{2}\mathbf{s}'\boldsymbol\Sigma_{UU}\mathbf{s}\bigr) \cdot \exp\bigl(\mathbf{t}'\boldsymbol\mu_V + \tfrac{1}{2}\mathbf{t}'\boldsymbol\Sigma_{VV}\mathbf{t}\bigr) = M_{\mathbf{U}}(\mathbf{s}) M_{\mathbf{V}}(\mathbf{t}),
\]
where the last step identifies the two factors as the marginal MGFs by setting $\mathbf{t} = \bzero$, respectively $\mathbf{s} = \bzero$, in the joint MGF. By fact~\ref{N2}, this factorization implies independence.
\end{proof}

\begin{remark}\label{rem:uncorr-indep}
For general (non-normal) distributions, uncorrelatedness does \emph{not} imply independence. The equivalence is a special property of the multivariate normal.
\end{remark}

\begin{theorem}[Chi-Square Distribution of Quadratic Forms]\label{thm:quadratic-chisq}
Let $\bz \sim N(\bzero, \bI_n)$ and $\bA$ be an $n \times n$ symmetric idempotent matrix with $\bA \neq \bzero$ (so that $r = \rank(\bA) \geqslant 1$, as the definition of $\chi^2_r$ in~\ref{N5} requires). Then
\begin{equation}
\bz'\bA\bz \sim \chi^2_r, \quad \text{where } r = \rank(\bA) = \tr(\bA).
\end{equation}
\end{theorem}

\begin{proof}
\textbf{Step 1: Spectral decomposition.} Since $\bA$ is symmetric, there exists an orthogonal matrix $\bP$ ($\bP'\bP = \bP\bP' = \bI$) such that $\bA = \bP\boldsymbol\Lambda\bP'$ (\ref{L3}), where $\boldsymbol\Lambda = \diag(\lambda_1, \ldots, \lambda_n)$. By idempotency, each $\lambda_i \in \{0, 1\}$ (Step~1 of the proof of Lemma~\ref{lem:idempotent}). Without loss of generality, let the first $r$ eigenvalues be $1$ and the rest be $0$:
\[
\boldsymbol\Lambda = \diag(\underbrace{1, \ldots, 1}_{r}, \underbrace{0, \ldots, 0}_{n-r}).
\]

\textbf{Step 2: Apply orthogonal transformation.} Let $\bw = \bP'\bz$. Since $\bP$ is orthogonal:
\[
\expc[\bw] = \bP'\expc[\bz] = \bzero, \qquad \var(\bw) = \bP'\bI\bP = \bI_n.
\]
By Theorem~\ref{thm:mvn-linear}, $\bw \sim N(\bzero, \bI_n)$. Its MGF $\exp(\tfrac{1}{2}\mathbf{t}'\mathbf{t}) = \prod_{i=1}^n \exp(\tfrac{1}{2}t_i^2)$ is the product of $n$ standard normal MGFs, so by fact~\ref{N2} (the $n$-block form of the argument in Theorem~\ref{thm:mvn-uncorr-indep}) $w_1, \ldots, w_n$ are i.i.d.\ $N(0, 1)$.

\textbf{Step 3: Evaluate the quadratic form.}
\[
\bz'\bA\bz = \bz'\bP\boldsymbol\Lambda\bP'\bz = (\bP'\bz)'\boldsymbol\Lambda(\bP'\bz) = \bw'\boldsymbol\Lambda\bw = \sum_{i=1}^{n}\lambda_i w_i^2 = \sum_{i=1}^{r} w_i^2.
\]

\textbf{Step 4: Conclude.} Since $w_1, \ldots, w_r$ are independent $N(0, 1)$, by definition $\sum_{i=1}^{r} w_i^2 \sim \chi^2_r$.
\end{proof}

\begin{theorem}[Independence of Quadratic Forms --- a Special Case of Craig's Theorem]\label{thm:cochran}
Let $\bz \sim N(\bzero, \bI_n)$, and let $\bA_1, \bA_2$ be $n \times n$ symmetric idempotent matrices with $\bA_1\bA_2 = \bzero$. Then $\bz'\bA_1\bz$ and $\bz'\bA_2\bz$ are independent.
\end{theorem}

\begin{proof}
Define $\mathbf{U} = \bA_1\bz$ and $\mathbf{V} = \bA_2\bz$. Stack them into a single vector obtained from $\bz$ by one linear map:
\[
\begin{pmatrix}\mathbf{U}\\\mathbf{V}\end{pmatrix} = \begin{pmatrix}\bA_1 \\ \bA_2\end{pmatrix}\bz.
\]
Since $\bz \sim N(\bzero, \bI_n)$, Theorem~\ref{thm:mvn-linear} implies $(\mathbf{U}', \mathbf{V}')'$ is \emph{jointly} normal. Their cross-covariance is, by Lemma~\ref{lem:var-linear}\ref{vl:ii},
\[
\cov(\mathbf{U}, \mathbf{V}) = \bA_1\var(\bz)\bA_2' = \bA_1\bI\bA_2 = \bA_1\bA_2 = \bzero,
\]
where $\bA_2' = \bA_2$ by symmetry. By Theorem~\ref{thm:mvn-uncorr-indep}, $\mathbf{U}$ and $\mathbf{V}$ are independent. By symmetric idempotency of $\bA_1$, $\bz'\bA_1\bz = \bz'\bA_1'\bA_1\bz = \|\mathbf{U}\|^2$ is a function of $\mathbf{U}$ alone; similarly $\bz'\bA_2\bz = \|\mathbf{V}\|^2$ depends only on $\mathbf{V}$. The two quadratic forms are therefore independent (fact~\ref{N3}).
\end{proof}

\begin{remark}\label{rem:cochran-use}
In Lemma~\ref{lem:SSR-chisq} (Section~\ref{sec:test}), Theorem~\ref{thm:cochran} is applied with $\bA_1 = \bH - n^{-1}\bone\bone'$ and $\bA_2 = \bM$ to prove $\SSR \indep \SSE$.
\end{remark}

%=============================================================================
\section{Exact Inference Under Normality}\label{sec:exact}
%=============================================================================

In this section and in Sections~\ref{sec:test} and~\ref{sec:pred}, all five assumptions~\ref{A1}--\ref{A5} are in force, with exceptions noted where they occur: the identities of Corollaries~\ref{cor:nested-F} and~\ref{cor:t2F} are algebraic and, like Theorem~\ref{thm:cls}, use only~\ref{A4} and $\rank(\bR) = q$, and the results of \S\ref{sec:test-more} that are algebraic or concern only moments say so. Sections~\ref{sec:R2}, \ref{sec:fwl} and~\ref{sec:special}, by contrast, consist of algebraic identities that hold for every realization of $\bY$ and use only the standing assumption~\ref{A4} (Theorem~\ref{thm:ovb} also uses~\ref{A2}, and the statements about expectations and variances at the end of Section~\ref{sec:fwl} use~\ref{A1}--\ref{A3} as stated there); this is restated at the start of each of those sections. Section~\ref{sec:qual} is likewise algebraic except where a distribution is asserted.

\subsection{Sampling Distribution of \texorpdfstring{$\hat{\bb}$}{the OLS estimator}}

\begin{theorem}[Normality of $\hat{\bb}$]\label{thm:beta-dist}
Under~\ref{A1}--\ref{A5}:
\begin{equation}
\boxed{\hat{\bb} \sim N\!\left(\bb,\; \sigma^2(\bX'\bX)^{-1}\right).}
\end{equation}
\end{theorem}

\begin{proof}
From~\eqref{eq:beta-decomp}, $\hat{\bb} = \bb + (\bX'\bX)^{-1}\bX'\be$. This is an affine transformation of $\be$. Under~\ref{A5}, $\be \sim N(\bzero, \sigma^2\bI)$. By Theorem~\ref{thm:mvn-linear} with $\bA = (\bX'\bX)^{-1}\bX'$, $\bz = \be$, $\mathbf{b} = \bb$:
\begin{align*}
\hat{\bb} &\sim N\!\bigl(\bb + (\bX'\bX)^{-1}\bX' \cdot \bzero,\; (\bX'\bX)^{-1}\bX'(\sigma^2\bI)\bX(\bX'\bX)^{-1}\bigr)\\
&= N\!\bigl(\bb,\; \sigma^2(\bX'\bX)^{-1}\bigr). \qedhere
\end{align*}
\end{proof}

\subsection{Distribution of \texorpdfstring{$\SSE/\sigma^2$}{SSE/sigma\texttwosuperior} and Independence from \texorpdfstring{$\hat{\bb}$}{the OLS estimator}}

\begin{theorem}[Chi-Square Distribution of $\SSE/\sigma^2$]\label{thm:sse-dist}
Under~\ref{A1}--\ref{A5}:
\begin{equation}
\frac{\SSE}{\sigma^2} \sim \chi^2_{n-k-1},
\end{equation}
and $\SSE$ is independent of $\hat{\bb}$.
\end{theorem}

\begin{proof}
\textbf{Step 1: Standardize.} Let $\bz = \be/\sigma$. Under~\ref{A5}, $\bz \sim N(\bzero, \bI_n)$. Since $\mathbf{e} = \bM\be$:
\[
\frac{\SSE}{\sigma^2} = \frac{\be'\bM\be}{\sigma^2} = \bz'\bM\bz.
\]

\textbf{Step 2: Apply Theorem~\ref{thm:quadratic-chisq}.} The matrix $\bM$ is symmetric and idempotent (Theorem~\ref{thm:HM}\ref{HM:i}), with $\tr(\bM) = n - k - 1$ (Theorem~\ref{thm:HM}\ref{HM:iv}). Hence $\bz'\bM\bz \sim \chi^2_{n-k-1}$.

\textbf{Step 3: Prove independence.} Stack $\hat{\bb} - \bb = (\bX'\bX)^{-1}\bX'\be$ and $\mathbf{e} = \bM\be$ into a single vector obtained from $\be$ by one linear map:
\[
\begin{pmatrix}\hat{\bb} - \bb \\ \mathbf{e}\end{pmatrix} = \begin{pmatrix}(\bX'\bX)^{-1}\bX' \\ \bM\end{pmatrix}\be.
\]
Since $\be \sim N(\bzero, \sigma^2\bI)$, Theorem~\ref{thm:mvn-linear} makes this stacked vector \emph{jointly} normal (individual normality alone would not suffice; joint normality requires that both pieces come from a single normal vector). Their cross-covariance is, by Lemma~\ref{lem:var-linear}\ref{vl:ii},
\begin{align*}
\cov\!\bigl((\bX'\bX)^{-1}\bX'\be,\; \bM\be\bigr) &= (\bX'\bX)^{-1}\bX'\var(\be)\bM'\\
&= (\bX'\bX)^{-1}\bX'(\sigma^2\bI)\bM = \sigma^2(\bX'\bX)^{-1}\bX'\bM = \bzero,
\end{align*}
using $\bX'\bM = (\bM\bX)' = \bzero$ (Theorem~\ref{thm:HM}\ref{HM:ii}). By Theorem~\ref{thm:mvn-uncorr-indep}, $\hat{\bb} - \bb$ and $\mathbf{e}$ are independent. Since $\hat{\bb}$ is $(\hat{\bb} - \bb)$ plus a constant and $\SSE = \mathbf{e}'\mathbf{e}$ is a function of $\mathbf{e}$, $\hat{\bb}$ is independent of $\mathbf{e}$ (hence of $\SSE$).
\end{proof}

\subsection{OLS as Maximum Likelihood Under Normality}

\begin{theorem}[OLS Equals MLE Under Normality]\label{thm:mle}
Under~\ref{A1}--\ref{A5}, the OLS estimator $\hat{\bb}$ is also the maximum likelihood estimator of $\bb$. The MLE of $\sigma^2$ is $\hat{\sigma}^2_{\mathrm{MLE}} = \SSE/n$, which is biased.
\end{theorem}

\begin{proof}
Under~\ref{A5}, $\be \sim N(\bzero, \sigma^2\bI)$, so by Theorem~\ref{thm:mvn-linear} applied to the affine transformation $\bY = \bX\bb + \be$, $\bY \sim N(\bX\bb, \sigma^2\bI)$. Since $\sigma^2\bI$ is positive definite, $\bY$ has the density of fact~\ref{N4} with $|\sigma^2\bI| = (\sigma^2)^n$, and the log-likelihood is:
\begin{align*}
\ell(\bb, \sigma^2) &= -\frac{n}{2}\ln(2\pi) - \frac{n}{2}\ln\sigma^2 - \frac{1}{2\sigma^2}(\bY - \bX\bb)'(\bY - \bX\bb).
\end{align*}

\textbf{Maximizing over $\bb$:} In $\ell(\bb, \sigma^2)$, the first two terms do not involve $\bb$, and the last term is $-(\bY - \bX\bb)'(\bY - \bX\bb)/(2\sigma^2)$ with $\sigma^2 > 0$ treated as fixed. Since $-1/(2\sigma^2) < 0$, maximizing $\ell$ over $\bb$ is equivalent to minimizing $(\bY - \bX\bb)'(\bY - \bX\bb) = Q(\bb)$. This is exactly the OLS problem, so $\hat{\bb}_{\mathrm{MLE}} = \hat{\bb}_{\mathrm{OLS}} = (\bX'\bX)^{-1}\bX'\bY$.

\textbf{Maximizing over $\sigma^2$:} Substituting $\hat{\bb}$ and noting that $(\bY - \bX\hat{\bb})'(\bY - \bX\hat{\bb}) = \SSE$, the profile log-likelihood becomes
\[
\ell(\hat{\bb}, \sigma^2) = -\frac{n}{2}\ln(2\pi) - \frac{n}{2}\ln\sigma^2 - \frac{\SSE}{2\sigma^2}.
\]
Differentiating with respect to $\sigma^2$:
\[
\frac{\partial \ell}{\partial \sigma^2} = -\frac{n}{2\sigma^2} + \frac{\SSE}{2(\sigma^2)^2} = 0 \implies \hat{\sigma}^2_{\mathrm{MLE}} = \frac{\SSE}{n}.
\]
\textbf{Second-order condition:} Writing $u = \sigma^2$ for brevity,
\[
\frac{\partial^2 \ell}{\partial u^2} = \frac{n}{2u^2} - \frac{\SSE}{u^3}.
\]
At the critical point $\hat{u} = \SSE/n$, so $u^2 = \SSE^2/n^2$ and $u^3 = \SSE^3/n^3$:
\[
\frac{\partial^2 \ell}{\partial u^2}\bigg|_{\hat{u}} = \frac{n \cdot n^2}{2\SSE^2} - \frac{\SSE \cdot n^3}{\SSE^3} = \frac{n^3}{2\SSE^2} - \frac{n^3}{\SSE^2} = -\frac{n^3}{2\SSE^2} < 0,
\]
confirming a local maximum. It is the global maximum over $u \in (0, \infty)$: here $\SSE > 0$ (this holds with probability one, since $\SSE/\sigma^2 \sim \chi^2_{n-k-1}$ by Theorem~\ref{thm:sse-dist} is a continuous distribution; if $\SSE = 0$ the likelihood is unbounded and no MLE of $\sigma^2$ exists), so $\ell(\hat{\bb}, u) \to -\infty$ both as $u \to 0^+$ (the term $-\SSE/(2u)$ dominates) and as $u \to \infty$ (the term $-\tfrac{n}{2}\ln u$ dominates). A continuous function on $(0,\infty)$ that tends to $-\infty$ at both ends attains its supremum at a critical point, and $\hat{u}$ is the only one. Finally, since for every fixed $u$ the maximizer over $\bb$ is $\hat{\bb}$ (the unique global minimizer of $Q$, Theorem~\ref{thm:ols}), $(\hat{\bb}, \hat{u})$ maximizes $\ell$ jointly.

\textbf{Bias of $\hat{\sigma}^2_{\mathrm{MLE}}$.} Using Theorem~\ref{thm:mse-unbiased} ($\expc[\SSE] = (n-k-1)\sigma^2$):
\[
\expc[\hat{\sigma}^2_{\mathrm{MLE}}] = \expc\!\left[\frac{\SSE}{n}\right] = \frac{(n-k-1)\sigma^2}{n} = \frac{n-k-1}{n}\sigma^2 \;<\; \sigma^2,
\]
so $\hat{\sigma}^2_{\mathrm{MLE}}$ underestimates $\sigma^2$ on average by a factor of $(n-k-1)/n$. The unbiased estimator is $\MSE = \SSE/(n-k-1)$, which rescales by this factor.
\end{proof}

\subsection{Distribution of the Residuals}

\begin{corollary}[Normality of the Residuals]\label{cor:resid-normal}
Under~\ref{A1}--\ref{A5}, $\mathbf{e} \sim N(\bzero, \sigma^2\bM)$, and for every $i$ with $h_{ii} < 1$,
\[
\frac{e_i}{\sigma\sqrt{1 - h_{ii}}} \sim N(0, 1).
\]
\end{corollary}

\begin{proof}
By Remark~\ref{rem:errors-residuals}, $\mathbf{e} = \bM\be$, so Theorem~\ref{thm:mvn-linear} with $\bA = \bM$ gives $\mathbf{e} \sim N(\bzero, \bM(\sigma^2\bI)\bM') = N(\bzero, \sigma^2\bM)$. Applying Theorem~\ref{thm:mvn-linear} again with $\bA = \boldsymbol\iota_i'/(\sigma\sqrt{1 - h_{ii}})$, the variance of $e_i/(\sigma\sqrt{1 - h_{ii}})$ is $\sigma^2(1 - h_{ii})/(\sigma^2(1 - h_{ii})) = 1$ and its mean is $0$.
\end{proof}

\begin{remark}\label{rem:std-resid-law}
The standardized residual of Definition~\ref{def:std-resid} is the $N(0,1)$ variable of Corollary~\ref{cor:resid-normal} multiplied by $\sigma/s$. It is neither exactly normal nor exactly $t$-distributed, because it is bounded: $e_i = \boldsymbol\iota_i'\mathbf{e} = \boldsymbol\iota_i'\bM\mathbf{e} = (\bM\boldsymbol\iota_i)'\mathbf{e}$ (as $\bM\mathbf{e} = \bM^2\bY = \mathbf{e}$), so Lemma~\ref{lem:cs} gives $e_i^2 \leqslant \|\bM\boldsymbol\iota_i\|^2\|\mathbf{e}\|^2 = \boldsymbol\iota_i'\bM\boldsymbol\iota_i \cdot \SSE = (1 - h_{ii})\SSE$, hence $\tilde e_i^2 \leqslant \SSE/s^2 = n - k - 1$. A bounded variable is not $N(0,1)$, since $\prb(Z > c) > 0$ for $Z \sim N(0,1)$ and every $c > 0$; nor is it $t_\nu$, since for $T = Z/\sqrt{V/\nu}$ as in~\ref{N5} we have $T \geqslant Z$ on $\{Z > 0,\ V \leqslant \nu\}$, so $\prb(T > c) \geqslant \prb(Z > c)\,\prb(V \leqslant \nu) > 0$, where $\prb(V \leqslant \nu) \geqslant \prb(|W_1| \leqslant 1, \ldots, |W_\nu| \leqslant 1) > 0$. The cutoff $|\tilde e_i| > 2$ is a rule of thumb, which we do not prove: when $n - k - 1$ is large, $s$ is close to $\sigma$ with high probability, $\tilde e_i$ then behaves approximately like the $N(0,1)$ variable of Corollary~\ref{cor:resid-normal}, and the cutoff flags, in a correctly specified model, roughly the $5\%$ of observations in the two tails of the standard normal distribution.
\end{remark}

%=============================================================================
\section{Sum of Squares and the Coefficient of Determination}\label{sec:R2}
%=============================================================================

The results of this section are algebraic identities about the fitted values and residuals of a given data set. They use only the standing assumption~\ref{A4} (applied to every design matrix involved) and, where stated, the presence of an intercept column; none of \ref{A1}--\ref{A3} or~\ref{A5} is used. (Theorem~\ref{thm:adj-R2} refers to the $t$ statistic of Section~\ref{sec:test} only as a name for an algebraic ratio.)

\subsection{Variance Decomposition}

\begin{definition}[Centering Matrix]\label{def:centering}
Let $\bJ = \frac{1}{n}\bone\bone' = \bone(\bone'\bone)^{-1}\bone'$, the hat matrix (Definition~\ref{def:hat-matrix}) of the intercept-only design $\bX = \bone$; by Remark~\ref{rem:hat-projection} it is the orthogonal projection onto $\calC(\bone)$, and $\bJ\bv = \bar{v}\bone$ for any vector $\bv$, where $\bar{v} = \frac{1}{n}\bone'\bv$. The \textbf{centering matrix} is the corresponding residual maker $\bM_0 = \bI_n - \bJ = \bI_n - \frac{1}{n}\bone\bone'$. For any vector $\bv$, $\bM_0\bv = \bv - \bar{v}\bone$.
\end{definition}

\begin{remark}[Notation: $\bM_0$ versus $\bM$]\label{rem:M0-vs-M}
The subscript ``$0$'' in $\bM_0$ signals that it is the residual maker with respect to the \emph{intercept only} (the single column $\bone$), whereas $\bM = \bI - \bX(\bX'\bX)^{-1}\bX'$ from Definition~\ref{def:hat-matrix} is the residual maker with respect to the \emph{full} design matrix $\bX$. Conceptually,
\[
\bM_0 \text{ centers a vector},\qquad \bM \text{ regresses out all of } \bX.
\]
Both are symmetric idempotent, but they project onto different subspaces:
\[
\bM_0 : \Real^n \to \bone^{\perp},\qquad \bM : \Real^n \to \calC(\bX)^{\perp}.
\]
If the model includes an intercept then $\bone \in \calC(\bX)$, hence $\calC(\bX)^\perp \subseteq \bone^\perp$, and so $\bM\bM_0 = \bM_0\bM = \bM$ (this follows from the projection identity $\bH\bJ = \bJ\bH = \bJ$ proved in Lemma~\ref{lem:SSR-chisq}, Step~2, by taking complements). In particular, $\bM$ and $\bM_0$ are not equal, nor is one a simple scalar multiple of the other; they coincide only in the trivial case $\bX = \bone$.
\end{remark}

\begin{remark}\label{rem:M0-props}
$\bM_0$ is symmetric and idempotent with $\tr(\bM_0) = n - 1$: it is the residual maker $\bM$ of the intercept-only design $\bX = \bone$, so Theorem~\ref{thm:HM}\ref{HM:i} and~\ref{HM:iv} apply with $k = 0$.
\end{remark}

\begin{theorem}[Variance Decomposition]\label{thm:anova}
In the model with an intercept, the sums of squares introduced in Definition~\ref{def:sumsq} satisfy:
\begin{equation}
\boxed{\SST = \SSR + \SSE,}
\end{equation}
with matrix expressions:
\begin{itemize}
\item $\SST = \sum(Y_i - \bar{Y})^2 = \bY'\bM_0\bY$ \quad (Total Sum of Squares),
\item $\SSR = \sum(\hat{Y}_i - \bar{Y})^2 = \hat{\bY}'\bM_0\hat{\bY}$ \quad (Regression Sum of Squares),
\item $\SSE = \sum(Y_i - \hat{Y}_i)^2 = \mathbf{e}'\mathbf{e}$ \quad (Error/Residual Sum of Squares).
\end{itemize}
\end{theorem}

\begin{proof}
\textbf{Step 1: Verify the matrix expressions for $\SST$ and $\SSR$.} For any vector $\bv \in \Real^n$, Definition~\ref{def:centering} gives $\bM_0\bv = \bv - \bar{v}\bone$, so
\[
\bv'\bM_0\bv = (\bM_0\bv)'(\bM_0\bv) = \|\bv - \bar{v}\bone\|^2 = \sum_{i=1}^n (v_i - \bar{v})^2,
\]
using $\bM_0'\bM_0 = \bM_0^2 = \bM_0$ (symmetric idempotent). Applying this with $\bv = \bY$ gives $\SST = \bY'\bM_0\bY$; with $\bv = \hat{\bY}$ gives $\SSR = \hat{\bY}'\bM_0\hat{\bY}$, provided $\overline{\hat{Y}} = \bar{Y}$—which follows from $\bone'\mathbf{e} = 0$ (Theorem~\ref{thm:normal-eq}\ref{ne:ii}) since $\bone'\hat{\bY} = \bone'(\bY - \mathbf{e}) = \bone'\bY$.

\textbf{Step 2: Decompose $\bM_0\bY$ into orthogonal components.} Write
\begin{equation}\label{eq:M0-decomp}
\bM_0\bY = \bM_0(\hat{\bY} + \mathbf{e}) = \bM_0\hat{\bY} + \bM_0\mathbf{e} = \bM_0\hat{\bY} + \mathbf{e},
\end{equation}
where the last equality uses $\bM_0\mathbf{e} = \mathbf{e} - \bar{e}\bone = \mathbf{e}$ because $\bar{e} = n^{-1}\bone'\mathbf{e} = 0$ (Theorem~\ref{thm:normal-eq}\ref{ne:ii}).

\textbf{Step 3: Expand $\SST$.} Using Step~2 and $\bM_0^2 = \bM_0$:
\begin{align*}
\SST = \bY'\bM_0\bY = (\bM_0\bY)'(\bM_0\bY) &= (\bM_0\hat{\bY} + \mathbf{e})'(\bM_0\hat{\bY} + \mathbf{e})\\
&= \SSR + 2\mathbf{e}'\bM_0\hat{\bY} + \SSE.
\end{align*}

\textbf{Step 4: The cross term vanishes.} Using $\mathbf{e}'\hat{\bY} = 0$ and $\mathbf{e}'\bone = 0$ (Theorem~\ref{thm:normal-eq}\ref{ne:ii}, \ref{ne:iii}):
\[
\mathbf{e}'\bM_0\hat{\bY} = \mathbf{e}'\hat{\bY} - \bar{Y}\,\mathbf{e}'\bone = 0.
\]
Hence $\SST = \SSR + \SSE$.
\end{proof}

\begin{remark}[Geometric view]\label{rem:geometric}
Equation~\eqref{eq:M0-decomp} is an orthogonal decomposition of $\bM_0\bY$ into two pieces: the centered fitted values $\bM_0\hat{\bY}$, which lie in the subspace $\calC(\bX) \cap \bone^\perp$ (the non-intercept directions of the column space), and the residuals $\mathbf{e}$, which lie in $\calC(\bX)^\perp$. Pythagoras applied to these two orthogonal pieces is precisely $\SST = \SSR + \SSE$.
\end{remark}

\subsection{Coefficient of Determination}

\begin{definition}[$R^2$ and Adjusted $R^2$]\label{def:R2}
Assume the model has an intercept and $\SST > 0$, i.e.\ $\bY$ is not a constant vector ($\SST = \|\bM_0\bY\|^2 = 0$ iff $\bY = \bar Y\bone$). The \textbf{coefficient of determination} is
\begin{equation}
R^2 = \frac{\SSR}{\SST} = 1 - \frac{\SSE}{\SST}.
\end{equation}
It measures the proportion of total variation in $\bY$ explained by the regression. The \textbf{adjusted $R^2$} is
\begin{equation}
\bar{R}^2 = 1 - \frac{\SSE/(n-k-1)}{\SST/(n-1)} = 1 - \frac{n-1}{n-k-1}(1 - R^2) = 1 - \frac{\MSE}{s_Y^2},
\end{equation}
where $s_Y^2 = \SST/(n-1)$. The adjusted $R^2$ can be negative.
\end{definition}

\begin{proposition}[Bounds on $R^2$]\label{prop:R2-bounds}
In any model with an intercept and $\SST > 0$, $0 \leqslant R^2 \leqslant 1$.
\end{proposition}

\begin{proof}
$R^2 = \SSR/\SST \geqslant 0$ because $\SSR = \sum(\hat{Y}_i - \bar{Y})^2 \geqslant 0$ and $\SST = \sum(Y_i - \bar{Y})^2 > 0$ by hypothesis. For the upper bound, Theorem~\ref{thm:anova} gives $\SST = \SSR + \SSE$ with $\SSE \geqslant 0$, so $\SSR \leqslant \SST$, i.e., $R^2 \leqslant 1$.
\end{proof}

\begin{theorem}[Monotonicity of $R^2$]\label{thm:R2-monotone}
Adding a regressor to the model cannot decrease $R^2$. That is, if $R^2_k$ is the coefficient of determination with $k$ regressors and $R^2_{k+1}$ is that with an additional regressor, then $R^2_{k+1} \geqslant R^2_k$.
\end{theorem}

\begin{proof}
Let $\calC_k \subseteq \calC_{k+1}$ denote the column spaces of the smaller and larger design matrices (the inclusion holds because every column of $\bX_k$ is also a column of $\bX_{k+1}$). Write $\hat{\bY}_k, \hat{\bY}_{k+1}$ for the corresponding OLS fits.

By Theorem~\ref{thm:ols-proj}, $\hat{\bY}_{k+1}$ is the point in $\calC_{k+1}$ closest to $\bY$, so $\|\bY - \hat{\bY}_{k+1}\| \leqslant \|\bY - \bv\|$ for every $\bv \in \calC_{k+1}$. In particular, $\hat{\bY}_k \in \calC_k \subseteq \calC_{k+1}$, so
\[
\SSE_{k+1} = \|\bY - \hat{\bY}_{k+1}\|^2 \;\leqslant\; \|\bY - \hat{\bY}_k\|^2 = \SSE_k.
\]
Since $\SST$ does not depend on the regressors,
\[
R^2_{k+1} = 1 - \frac{\SSE_{k+1}}{\SST} \;\geqslant\; 1 - \frac{\SSE_k}{\SST} = R^2_k. \qedhere
\]
\end{proof}

\begin{lemma}[Reduction in $\SSE$ from adding one regressor]\label{lem:sse-reduction}
Write the enlarged design as $\bX_{k+1} = [\bX_k \;\; \bx]$, where $\bx$ is the new regressor column, and assume $\bx'\bM_k\bx > 0$ (where $\bM_k$ is the residual maker for $\bX_k$), which holds whenever $\bx$ is not in $\calC(\bX_k)$. Then the OLS coefficient on $\bx$ in the enlarged model is
\begin{equation}\label{eq:beta-new}
\hat{\beta}_{k+1} = \frac{\bx'\bM_k\bY}{\bx'\bM_k\bx},
\end{equation}
and the reduction in residual sum of squares is
\begin{equation}\label{eq:sse-reduction}
\SSE_k - \SSE_{k+1} = \hat{\beta}_{k+1}^2 \cdot \bx'\bM_k\bx = \frac{(\bx'\bM_k\bY)^2}{\bx'\bM_k\bx} \;\geqslant\; 0.
\end{equation}
This identity uses only normal equations (no forward dependence on Section~\ref{sec:fwl}).
\end{lemma}

\begin{proof}
The partitioned normal equations $\bX_{k+1}'\bX_{k+1}\hat{\bb}_{k+1} = \bX_{k+1}'\bY$ read
\begin{align*}
\bX_k'\bX_k\hat{\bb}_k^* + \bX_k'\bx\hat{\beta}_{k+1} &= \bX_k'\bY, \tag{a}\\
\bx'\bX_k\hat{\bb}_k^* + \bx'\bx\,\hat{\beta}_{k+1} &= \bx'\bY, \tag{b}
\end{align*}
where $\hat{\bb}_k^*$ denotes the $\bX_k$-block of the enlarged estimator (its first $k+1$ coordinates, the coefficients on the columns of $\bX_k$; generally different from the short-regression estimator $\hat{\bb}_k$). From (a), $\hat{\bb}_k^* = (\bX_k'\bX_k)^{-1}\bX_k'(\bY - \bx\hat{\beta}_{k+1})$. Substituting into (b):
\[
\bx'\bX_k(\bX_k'\bX_k)^{-1}\bX_k'(\bY - \bx\hat{\beta}_{k+1}) + \bx'\bx\,\hat{\beta}_{k+1} = \bx'\bY.
\]
Writing $\bH_k = \bX_k(\bX_k'\bX_k)^{-1}\bX_k' = \bI - \bM_k$ and rearranging:
\[
\bx'\bH_k\bY - \bx'\bH_k\bx\,\hat{\beta}_{k+1} + \bx'\bx\,\hat{\beta}_{k+1} = \bx'\bY
\;\ifff\; \bx'(\bI - \bH_k)\bx\,\hat{\beta}_{k+1} = \bx'(\bI - \bH_k)\bY,
\]
which gives~\eqref{eq:beta-new}.

For~\eqref{eq:sse-reduction}, relate $\hat{\bb}_k^*$ to the short-regression estimator $\hat{\bb}_k = (\bX_k'\bX_k)^{-1}\bX_k'\bY$: from (a),
\[
\hat{\bb}_k^* = \hat{\bb}_k - (\bX_k'\bX_k)^{-1}\bX_k'\bx\,\hat{\beta}_{k+1}.
\]
Multiplying by $\bX_k$ gives $\bX_k\hat{\bb}_k^* = \hat{\bY}_k - \bH_k\bx\,\hat{\beta}_{k+1}$, so
\[
\hat{\bY}_{k+1} = \bX_k\hat{\bb}_k^* + \bx\hat{\beta}_{k+1} = \hat{\bY}_k + (\bI - \bH_k)\bx\,\hat{\beta}_{k+1} = \hat{\bY}_k + \bM_k\bx\,\hat{\beta}_{k+1},
\]
and the residuals satisfy $\mathbf{e}_{k+1} = \mathbf{e}_k - \bM_k\bx\,\hat{\beta}_{k+1}$. Taking squared norms, using $\mathbf{e}_k = \bM_k\bY$ and $\bM_k^2 = \bM_k$:
\[
\SSE_{k+1} = \SSE_k - 2\hat{\beta}_{k+1}\cdot\bY'\bM_k\bx + \hat{\beta}_{k+1}^2\cdot\bx'\bM_k\bx.
\]
Substituting $\hat{\beta}_{k+1} = (\bx'\bM_k\bY)/(\bx'\bM_k\bx)$ makes the middle term $-2\hat{\beta}_{k+1}^2\cdot\bx'\bM_k\bx$, so
\[
\SSE_{k+1} = \SSE_k - \hat{\beta}_{k+1}^2\cdot\bx'\bM_k\bx,
\]
which is~\eqref{eq:sse-reduction}. The non-negativity reflects Theorem~\ref{thm:R2-monotone}.
\end{proof}

\begin{theorem}[$\bar{R}^2$ Increases iff $|t| > 1$]\label{thm:adj-R2}
Let $\bX_{k+1} = [\bX_k \;\; \bx]$ as in Lemma~\ref{lem:sse-reduction}, with $\bx \notin \calC(\bX_k)$ (so that $\bx'\bM_k\bx > 0$ and $\bX_{k+1}$ has full column rank), and assume $n > k + 2$, $\SST > 0$ and $\SSE_{k+1} > 0$, so that both adjusted $R^2$'s and the $t$ statistic below are defined. Then adding the regressor $\bx$ increases $\bar{R}^2$ if and only if its $t$-statistic in the enlarged model exceeds $1$ in absolute value.
\end{theorem}

\begin{proof}
Let the model with $k$ regressors have $\MSE_k = \SSE_k/(n-k-1)$, and the model with an additional regressor have $\MSE_{k+1} = \SSE_{k+1}/(n-k-2)$. Since $\bar{R}^2 = 1 - \MSE/s_Y^2$ with the same $s_Y^2 = \SST/(n-1) > 0$ in both models, $\bar{R}^2_{k+1} > \bar{R}^2_k$ if and only if $\MSE_{k+1} < \MSE_k$, i.e.,
\begin{equation}\label{eq:mse-compare}
\frac{\SSE_{k+1}}{n-k-2} < \frac{\SSE_k}{n-k-1}.
\end{equation}

\textbf{Step 1: Apply Lemma~\ref{lem:sse-reduction}.}
By~\eqref{eq:sse-reduction}, $\SSE_k = \SSE_{k+1} + \hat{\beta}_{k+1}^2\cdot\bx'\bM_k\bx$. Substitute this into~\eqref{eq:mse-compare}:
\begin{align*}
\frac{\SSE_{k+1}}{n-k-2} &< \frac{\SSE_{k+1} + \hat{\beta}_{k+1}^2\cdot\bx'\bM_k\bx}{n-k-1}\\
(n-k-1)\SSE_{k+1} &< (n-k-2)\bigl[\SSE_{k+1} + \hat{\beta}_{k+1}^2\cdot\bx'\bM_k\bx\bigr]\\
\SSE_{k+1} &< (n-k-2)\cdot\hat{\beta}_{k+1}^2\cdot\bx'\bM_k\bx.
\end{align*}
Dividing by $\SSE_{k+1} = (n-k-2)\MSE_{k+1} > 0$:
\begin{equation}\label{eq:mse-simplified}
1 < \frac{\hat{\beta}_{k+1}^2\cdot\bx'\bM_k\bx}{\MSE_{k+1}}.
\end{equation}

\textbf{Step 2: Identify the right-hand side as $t^2$.}
The $t$-statistic for $\hat{\beta}_{k+1}$ in the enlarged model is $t_{k+1} = \hat{\beta}_{k+1}/\SE(\hat{\beta}_{k+1})$ with $\SE(\hat{\beta}_{k+1})^2 = \MSE_{k+1}\cdot [(\bX_{k+1}'\bX_{k+1})^{-1}]_{k+1,k+1}$ (defined formally in Theorem~\ref{thm:t-test}; here $[\cdot]_{k+1,k+1}$ is the diagonal entry belonging to $\beta_{k+1}$, the last one). We claim
\begin{equation}\label{eq:last-diag}
[(\bX_{k+1}'\bX_{k+1})^{-1}]_{k+1,k+1} = (\bx'\bM_k\bx)^{-1}.
\end{equation}
Indeed, let $\mathbf{G} = (\bX_{k+1}'\bX_{k+1})^{-1}\bX_{k+1}'$, so that $\hat{\bb}_{k+1} = \mathbf{G}\bY$ for every $\bY \in \Real^n$. Lemma~\ref{lem:sse-reduction} says that the last coordinate of $\mathbf{G}\bY$ is $\bx'\bM_k\bY/(\bx'\bM_k\bx)$ for every $\bY$, so the last row of $\mathbf{G}$ is $\mathbf{g}' := \bx'\bM_k/(\bx'\bM_k\bx)$. On the other hand $\mathbf{G}\mathbf{G}' = (\bX_{k+1}'\bX_{k+1})^{-1}\bX_{k+1}'\bX_{k+1}(\bX_{k+1}'\bX_{k+1})^{-1} = (\bX_{k+1}'\bX_{k+1})^{-1}$, whose last diagonal entry is the squared norm of the last row of $\mathbf{G}$:
\[
[(\bX_{k+1}'\bX_{k+1})^{-1}]_{k+1,k+1} = \mathbf{g}'\mathbf{g} = \frac{\bx'\bM_k\bM_k\bx}{(\bx'\bM_k\bx)^2} = \frac{1}{\bx'\bM_k\bx},
\]
using $\bM_k' = \bM_k = \bM_k^2$. (Equivalently: $\var(\hat{\beta}_{k+1})$ computed from~\eqref{eq:beta-new} via Lemma~\ref{lem:var-linear} is $\sigma^2/(\bx'\bM_k\bx)$, and computed from Theorem~\ref{thm:var-beta} it is $\sigma^2[(\bX_{k+1}'\bX_{k+1})^{-1}]_{k+1,k+1}$; this is the block-inverse, or Schur complement, formula for the last diagonal entry, obtained here from Lemma~\ref{lem:sse-reduction} alone.) Hence $\SE(\hat{\beta}_{k+1})^2 = \MSE_{k+1}/(\bx'\bM_k\bx)$, so
\[
t_{k+1}^2 = \frac{\hat{\beta}_{k+1}^2}{\SE(\hat{\beta}_{k+1})^2} = \frac{\hat{\beta}_{k+1}^2\cdot\bx'\bM_k\bx}{\MSE_{k+1}}.
\]
Substituting into~\eqref{eq:mse-simplified} gives $t_{k+1}^2 > 1$, i.e., $|t_{k+1}| > 1$.
\end{proof}

\begin{remark}[Logical status of the forward reference]\label{rem:forward-adjR2}
The proof refers to one later statement: the form of $\SE(\hat{\beta}_{k+1})$ in Theorem~\ref{thm:t-test}, which is a \emph{definition} of the $t$ statistic rather than a result whose proof is needed here. It is used only to \emph{interpret} the ratio $\hat{\beta}_{k+1}^2\cdot\bx'\bM_k\bx/\MSE_{k+1}$ as $t_{k+1}^2$; equation~\eqref{eq:mse-simplified} and the identity~\eqref{eq:last-diag} are self-contained algebraic consequences of Lemma~\ref{lem:sse-reduction}. A reader seeking strict linearity can defer Theorem~\ref{thm:adj-R2} until after Section~\ref{sec:test}.
\end{remark}

\subsection{\texorpdfstring{$R^2$}{R2} as a Squared Correlation}

\begin{proposition}[$R^2$ as a Squared Correlation]\label{prop:R2-corr}
In a model with an intercept and $\SST > 0$:
\begin{enumerate}[label=\textup{(\roman*)}]
\item if $\SSR > 0$, then $\hat{\bY}$ is not constant and $R^2 = \corr(\bY, \hat{\bY})^2$;
\item $\corr(\bY, \bv)^2 \leqslant R^2$ for every non-constant $\bv \in \calC(\bX)$; by (i), equality holds for $\bv = \hat{\bY}$ when $\SSR > 0$.
\end{enumerate}
Thus $\hat{\bY}$ is the linear combination of the regressors that is most highly correlated with $\bY$, and $R^2$ is the square of that correlation.
\end{proposition}

\begin{proof}
\textbf{Step 1: Centred cross-products.} For $\bu, \bv \in \Real^n$, Definition~\ref{def:centering} gives $\bM_0\bu = \bu - \bar u\bone$, so $S_{uv} = (\bM_0\bu)'(\bM_0\bv)$ in the notation of Definition~\ref{def:corr}. Let $\bv \in \calC(\bX)$. Then $\bM_0\bv = \bv - \bar v\bone \in \calC(\bX)$ because $\bone \in \calC(\bX)$, so $\mathbf{e}'\bM_0\bv = 0$ (Theorem~\ref{thm:normal-eq}\ref{ne:i}); together with $\bM_0\mathbf{e} = \mathbf{e}$ (equation~\eqref{eq:M0-decomp}) this gives
\[
S_{Yv} = (\bM_0\hat{\bY} + \mathbf{e})'\bM_0\bv = (\bM_0\hat{\bY})'(\bM_0\bv) = S_{\hat Y v}.
\]

\textbf{Step 2: The correlation with $\hat{\bY}$.} By Theorem~\ref{thm:anova}, $S_{\hat Y\hat Y} = \|\bM_0\hat{\bY}\|^2 = \SSR > 0$ and $S_{YY} = \SST$; in particular $\hat{\bY}$ is not constant. Step~1 with $\bv = \hat{\bY}$ gives $S_{Y\hat Y} = \SSR$, so
\[
\corr(\bY, \hat{\bY})^2 = \frac{S_{Y\hat Y}^2}{S_{YY}S_{\hat Y\hat Y}} = \frac{\SSR^2}{\SST\cdot\SSR} = \frac{\SSR}{\SST} = R^2.
\]

\textbf{Step 3: The maximum.} Let $\bv \in \calC(\bX)$ be non-constant. By Step~1 and Lemma~\ref{lem:cs} applied to $\bM_0\hat{\bY}$ and $\bM_0\bv$,
\[
S_{Yv}^2 = \bigl((\bM_0\hat{\bY})'(\bM_0\bv)\bigr)^2 \leqslant \|\bM_0\hat{\bY}\|^2\,\|\bM_0\bv\|^2 = \SSR\cdot S_{vv},
\]
so $\corr(\bY, \bv)^2 = S_{Yv}^2/(\SST\cdot S_{vv}) \leqslant \SSR/\SST = R^2$.
\end{proof}

\begin{corollary}[$R^2$ in Simple Linear Regression]\label{cor:R2-slr}
In the simple linear regression model of \S\ref{sec:k1}, with $S_{xx} > 0$ and $\SST > 0$,
\[
R^2 = r^2, \qquad \SSR = \hat\beta_1^2S_{xx} = \hat\beta_1^2(n-1)s_x^2,
\]
and, if $n > 2$, $s_Y^2 = \hat\beta_1^2s_x^2 + s^2\,\dfrac{n-2}{n-1}$, where $s^2 = \SSE/(n-2)$ is the mean squared error of Definition~\ref{def:sumsq}.
\end{corollary}

\begin{proof}
By Corollary~\ref{cor:slr}\ref{slr:iv}, $\SSR = \sum_i(\hat Y_i - \bar Y)^2 = \hat\beta_1^2\sum_i(x_i - \bar x)^2 = \hat\beta_1^2S_{xx}$. Since $S_{YY} = \SST > 0$, $\bY$ is not constant, and Corollary~\ref{cor:slr}\ref{slr:iii} gives $\hat\beta_1^2 = r^2S_{YY}/S_{xx}$, so $\SSR = r^2\,\SST$, i.e.\ $R^2 = r^2$. Finally, dividing $\SST = \SSR + \SSE$ (Theorem~\ref{thm:anova}) by $n - 1$, with $\SST/(n-1) = s_Y^2$, $\SSR/(n-1) = \hat\beta_1^2s_x^2$ and $\SSE = (n-2)s^2$, gives the last identity.
\end{proof}

%=============================================================================
\section{Hypothesis Testing and Confidence Regions}\label{sec:test}
%=============================================================================

\subsection{Individual Coefficient \texorpdfstring{$t$}{t} Test}

\begin{theorem}[Individual $t$ Test]\label{thm:t-test}
Under~\ref{A1}--\ref{A5}, for testing $H_0\colon \beta_j = \beta_j^0$:
\begin{equation}
t = \frac{\hat{\beta}_j - \beta_j^0}{\SE(\hat{\beta}_j)} \sim t_{n-k-1} \quad \text{under } H_0,
\end{equation}
where $\SE(\hat{\beta}_j) = s\sqrt{[(\bX'\bX)^{-1}]_{jj}}$ and $s = \sqrt{\MSE}$.
\end{theorem}

\begin{proof}
\textbf{Step 1: Marginal distribution of $\hat\beta_j$.} By Theorem~\ref{thm:beta-dist}, $\hat{\bb} \sim N(\bb, \sigma^2(\bX'\bX)^{-1})$. The $j$-th coordinate of a multivariate normal vector is itself normal, with mean $\beta_j$ and variance equal to the $(j,j)$-entry of the covariance matrix (this is Theorem~\ref{thm:mvn-linear} applied with $\bA = \boldsymbol\iota_j'$, the transpose of the $j$-th standard basis vector):
\[
\hat{\beta}_j = \boldsymbol\iota_j'\hat{\bb} \sim N\!\bigl(\beta_j,\; \sigma^2[(\bX'\bX)^{-1}]_{jj}\bigr).
\]
Standardizing:
\[
Z = \frac{\hat{\beta}_j - \beta_j}{\sigma\sqrt{[(\bX'\bX)^{-1}]_{jj}}} \sim N(0, 1).
\]

\textbf{Step 2: Chi-square quantity independent of $Z$.} By Theorem~\ref{thm:sse-dist}, $V = \SSE/\sigma^2 \sim \chi^2_{n-k-1}$ and is independent of $\hat{\bb}$. Since $Z$ is a (deterministic) function of $\hat{\beta}_j$, which is a coordinate of $\hat{\bb}$, and a function of an independent random variable remains independent, $Z$ is independent of $V$.

\textbf{Step 3.} By definition, if $Z \sim N(0,1)$ and $V \sim \chi^2_\nu$ are independent, then $Z/\sqrt{V/\nu} \sim t_\nu$. Applying this:
\begin{align*}
\frac{Z}{\sqrt{V/(n-k-1)}} &= \frac{\hat{\beta}_j - \beta_j}{\sigma\sqrt{[(\bX'\bX)^{-1}]_{jj}}} \cdot \frac{1}{\sqrt{\SSE/[\sigma^2(n-k-1)]}}\\
&= \frac{\hat{\beta}_j - \beta_j}{\sqrt{\MSE} \cdot \sqrt{[(\bX'\bX)^{-1}]_{jj}}} = \frac{\hat{\beta}_j - \beta_j}{s\sqrt{[(\bX'\bX)^{-1}]_{jj}}} \sim t_{n-k-1}.
\end{align*}
Under $H_0\colon \beta_j = \beta_j^0$, replace $\beta_j$ by $\beta_j^0$.
\end{proof}

\subsection{Overall \texorpdfstring{$F$}{F} Test}

\begin{lemma}[Distribution of $\SSR/\sigma^2$ Under $H_0$]\label{lem:SSR-chisq}
Let $k \geqslant 1$, and let $\bJ = \frac{1}{n}\bone\bone'$ be the projection onto $\calC(\bone)$ of Definition~\ref{def:centering}. Under $H_0\colon \beta_1 = \cdots = \beta_k = 0$ (with~\ref{A1}--\ref{A5}):
\[
\frac{\SSR}{\sigma^2} \sim \chi^2_k, \qquad \text{and } \SSR \text{ is independent of } \SSE.
\]
\end{lemma}

\begin{proof}
\textbf{Step 1: Express $\SSR/\sigma^2$ as a quadratic form in standard normals.} Two matrix observations are useful:
\begin{itemize}
\item $\hat{\bY} = \bH\bY$ by definition of the hat matrix.
\item $\bar{Y}\bone = \bJ\bY$, because $\bJ\bY = n^{-1}\bone(\bone'\bY) = n^{-1}\bone \cdot n\bar{Y} = \bar{Y}\bone$.
\end{itemize}
Therefore
\[
\hat{\bY} - \bar{Y}\bone = \bH\bY - \bJ\bY = (\bH - \bJ)\bY,
\]
and $\SSR = (\hat{\bY} - \bar{Y}\bone)'(\hat{\bY} - \bar{Y}\bone) = \bY'(\bH - \bJ)'(\bH - \bJ)\bY$.

Under $H_0$, $\bb = (\beta_0, 0, \ldots, 0)'$, so $\bX\bb = \beta_0\bone$ and $\bY = \beta_0\bone + \be$. We claim $(\bH - \bJ)\bY = (\bH - \bJ)\be$: indeed, since $\bone \in \calC(\bX)$ we have $\bH\bone = \bone$ (by Theorem~\ref{thm:HM}\ref{HM:ii} applied to the first column of $\bX$), and $\bJ\bone = n^{-1}\bone(\bone'\bone) = \bone$, so
\[
(\bH - \bJ)(\beta_0\bone) = \beta_0(\bH\bone - \bJ\bone) = \beta_0(\bone - \bone) = \bzero,
\]
and hence $(\bH - \bJ)\bY = (\bH - \bJ)(\beta_0\bone + \be) = (\bH - \bJ)\be$.

Let $\bz = \be/\sigma$. Under~\ref{A5}, $\bz \sim N(\bzero, \bI_n)$. Then
\[
\frac{\SSR}{\sigma^2} = \frac{\be'(\bH - \bJ)'(\bH - \bJ)\be}{\sigma^2} = \bz'(\bH - \bJ)'(\bH - \bJ)\bz.
\]
This reduces to $\bz'(\bH - \bJ)\bz$ once we establish in Step~2 that $\bH - \bJ$ is symmetric idempotent, which implies $(\bH - \bJ)'(\bH - \bJ) = (\bH - \bJ)^2 = \bH - \bJ$.

\textbf{Step 2: $\bH - \bJ$ is symmetric idempotent with rank $k$.} Symmetry is immediate. For idempotency, we first show
\begin{equation}\label{eq:HH0-commute}
\bH\bJ = \bJ = \bJ\bH.
\end{equation}
Using $\bH\bone = \bone$: $\bH\bJ = n^{-1}\bH\bone\bone' = n^{-1}\bone\bone' = \bJ$; by symmetry, $\bJ\bH = (\bH\bJ)' = \bJ$. Also $\bJ^2 = n^{-2}\bone(\bone'\bone)\bone' = \bJ$. Then
\[
(\bH - \bJ)^2 = \bH^2 - \bH\bJ - \bJ\bH + \bJ^2 = \bH - \bJ - \bJ + \bJ = \bH - \bJ.
\]
For the rank, Lemma~\ref{lem:idempotent} gives
\[
\rank(\bH - \bJ) = \tr(\bH) - \tr(\bJ) = (k+1) - 1 = k,
\]
with $\tr(\bJ) = n^{-1}\tr(\bone\bone') = 1$.

The Step~1 expression simplifies: $(\bH - \bJ)'(\bH - \bJ) = \bH - \bJ$, so $\SSR/\sigma^2 = \bz'(\bH - \bJ)\bz \sim \chi^2_k$ by Theorem~\ref{thm:quadratic-chisq}.

\textbf{Step 3: Independence of $\SSR$ and $\SSE$.} Apply Theorem~\ref{thm:cochran} to $\bA_1 = \bH - \bJ$ and $\bA_2 = \bM$. The required $\bA_1\bA_2 = \bzero$ follows from $\bH\bM = \bzero$ (Theorem~\ref{thm:HM}\ref{HM:iii}) and $\bJ\bM = \bJ - \bJ\bH = \bzero$ by~\eqref{eq:HH0-commute}:
\[
(\bH - \bJ)\bM = \bzero.
\]
Hence $\SSR/\sigma^2 = \bz'(\bH - \bJ)\bz$ and $\SSE/\sigma^2 = \bz'\bM\bz$ are independent, and so are $\SSR$ and $\SSE$.
\end{proof}

\begin{theorem}[Overall $F$ Test]\label{thm:f-test}
Let $k \geqslant 1$. Under~\ref{A1}--\ref{A5}, for testing $H_0\colon \beta_1 = \beta_2 = \cdots = \beta_k = 0$ (the third expression below presupposes $\SST > 0$, so that $R^2$ is defined; under~\ref{A5} this holds with probability one, see the end of the proof):
\begin{equation}\label{eq:f-test}
\boxed{F = \frac{\MSR}{\MSE} = \frac{\SSR/k}{\SSE/(n-k-1)} = \frac{R^2/k}{(1-R^2)/(n-k-1)} \sim F_{k,\,n-k-1} \quad \text{under } H_0.}
\end{equation}
\end{theorem}

\begin{proof}
By Lemma~\ref{lem:SSR-chisq} and Theorem~\ref{thm:sse-dist}, under $H_0$: $\SSR/\sigma^2 \sim \chi^2_k$, $\SSE/\sigma^2 \sim \chi^2_{n-k-1}$, and they are independent. By the definition of the $F$ distribution (fact~\ref{N5}):
\[
F = \frac{(\SSR/\sigma^2)/k}{(\SSE/\sigma^2)/(n-k-1)} = \frac{\SSR/k}{\SSE/(n-k-1)} = \frac{\MSR}{\MSE} \sim F_{k, n-k-1}.
\]
The equivalence $F = \frac{R^2/k}{(1-R^2)/(n-k-1)}$ follows from $R^2 = \SSR/\SST$ and $1 - R^2 = \SSE/\SST$, which require $\SST > 0$ (Definition~\ref{def:R2}). Under~\ref{A5} this holds with probability one: pick a unit vector $\bv \perp \bone$ (one exists since $n \geqslant 2$), so that $\bM_0\bv = \bv$ and hence $\bv'\bM_0 = \bv'$; if $\SST = \|\bM_0\bY\|^2 = 0$ then $\bv'\bY = \bv'\bM_0\bY = 0$, whereas $\bv'\bY \sim N(\bv'\bX\bb, \sigma^2)$ (Theorem~\ref{thm:mvn-linear} applied to $\bY \sim N(\bX\bb, \sigma^2\bI)$, as in the proof of Theorem~\ref{thm:mle}) with $\sigma^2 > 0$ by~\ref{A3} is a continuous distribution, so $\prb(\bv'\bY = 0) = 0$.
\end{proof}

\subsection{General Linear Hypothesis}

\begin{theorem}[General $F$ Test]\label{thm:general-F}
Under~\ref{A1}--\ref{A5}, consider $H_0\colon \bR\bb = \br$ where $\bR$ is a $q \times (k+1)$ matrix of rank $q$ and $\br \in \Real^q$. The test statistic
\begin{equation}\label{eq:general-F}
\boxed{F = \frac{(\bR\hat{\bb} - \br)'[\bR(\bX'\bX)^{-1}\bR']^{-1}(\bR\hat{\bb} - \br)}{q \cdot s^2} \sim F_{q,\, n-k-1} \quad \text{under } H_0.}
\end{equation}
\end{theorem}

\begin{proof}
\textbf{Step 1.} By Theorem~\ref{thm:beta-dist}, $\hat{\bb} \sim N(\bb, \sigma^2(\bX'\bX)^{-1})$. By Theorem~\ref{thm:mvn-linear}:
\[
\bR\hat{\bb} \sim N\!\bigl(\bR\bb,\; \sigma^2\bR(\bX'\bX)^{-1}\bR'\bigr).
\]
Under $H_0$, $\bR\bb = \br$, so $\bR\hat{\bb} - \br \sim N(\bzero, \sigma^2\bR(\bX'\bX)^{-1}\bR')$.

\textbf{Step 2: A quadratic form in $\bR\hat{\bb} - \br$ is chi-square.} Let $\boldsymbol\Sigma_R = \bR(\bX'\bX)^{-1}\bR'$. $\boldsymbol\Sigma_R$ is positive definite: for nonzero $\bv \in \Real^q$,
\[
\bv'\boldsymbol\Sigma_R\bv = (\bR'\bv)'(\bX'\bX)^{-1}(\bR'\bv) > 0,
\]
since $\rank(\bR) = q$ makes $\bR'\bv \neq \bzero$, and $(\bX'\bX)^{-1}$ is positive definite (Theorem~\ref{thm:ols}, Step~4). Let $\boldsymbol\Sigma_R^{1/2}$ be its positive definite symmetric square root (\ref{L4}).

Set $\bw = \sigma^{-1}\boldsymbol\Sigma_R^{-1/2}(\bR\hat{\bb} - \br)$. Theorem~\ref{thm:mvn-linear} gives $\bw \sim N(\bzero, \bI_q)$, so $\bw'\bw \sim \chi^2_q$ by Theorem~\ref{thm:quadratic-chisq}:
\[
W \;:=\; \bw'\bw = \frac{(\bR\hat{\bb} - \br)'\boldsymbol\Sigma_R^{-1}(\bR\hat{\bb} - \br)}{\sigma^2} \sim \chi^2_q.
\]

\textbf{Step 3.} By Theorem~\ref{thm:sse-dist}, $V = \SSE/\sigma^2 \sim \chi^2_{n-k-1}$, independent of $\hat{\bb}$; since $W$ is a function of $\hat{\bb}$, $W$ is independent of $V$. Then:
\[
F = \frac{W/q}{V/(n-k-1)} = \frac{(\bR\hat{\bb} - \br)'\boldsymbol\Sigma_R^{-1}(\bR\hat{\bb} - \br)/(q\sigma^2)}{\SSE/[\sigma^2(n-k-1)]} = \frac{(\bR\hat{\bb} - \br)'\boldsymbol\Sigma_R^{-1}(\bR\hat{\bb} - \br)}{q \cdot s^2} \sim F_{q, n-k-1}. \qedhere
\]
\end{proof}

\begin{remark}[Special cases]\label{rem:special-cases}
The individual $t$ test (Theorem~\ref{thm:t-test}) is the case $q = 1$, $\bR = \boldsymbol\iota_j'$ (the transpose of the $j$-th standard basis vector), $\br = \beta_j^0$, where $t^2 = F \sim F_{1, n-k-1}$. The overall $F$ test (Theorem~\ref{thm:f-test}) is the case $\bR = [\bzero_{k \times 1} \mid \bI_k]$ (of rank $q = k$), $\br = \bzero_k$: by Corollary~\ref{cor:nested-F} below, the statistic~\eqref{eq:general-F} then equals $\frac{(\SSE_R - \SSE_U)/k}{\SSE_U/(n-k-1)}$, where the reduced model minimizes $\|\bY - \bX\bb\|^2$ over $\bb = (\beta_0, 0, \ldots, 0)'$, i.e.\ $\|\bY - \beta_0\bone\|^2$ over $\beta_0$; its normal equation $\bone'\bone\,\beta_0 = \bone'\bY$ (Theorem~\ref{thm:ols} with $\bX = \bone$) gives $\hat{\beta}_0 = \bar{Y}$, so the reduced fit is $\bar{Y}\bone$, $\SSE_R = \|\bY - \bar{Y}\bone\|^2 = \SST$, and $\SSE_R - \SSE_U = \SST - \SSE = \SSR$ by Theorem~\ref{thm:anova}. The statistic is therefore $(\SSR/k)/(\SSE/(n-k-1)) = \MSR/\MSE$, as in~\eqref{eq:f-test}.
\end{remark}

\begin{corollary}[Nested $F$ Test via Residual Sum of Squares]\label{cor:nested-F}
Let $\SSE_U$ and $\SSE_R$ be the residual sums of squares from the full and reduced models under $H_0\colon \bR\bb = \br$ ($q$ constraints, $\rank(\bR) = q$). Then the $F$ statistic in Theorem~\ref{thm:general-F} is equivalently expressed as
\begin{equation}
F = \frac{(\SSE_R - \SSE_U)/q}{\SSE_U/(n-k-1)}.
\end{equation}
(This is an algebraic identity between the two expressions for $F$: like Theorem~\ref{thm:cls}, on which its proof rests, it uses only~\ref{A4} and $\rank(\bR) = q$. Assumptions~\ref{A1}--\ref{A5} enter only through the distribution $F \sim F_{q,\, n-k-1}$ under $H_0$ in Theorem~\ref{thm:general-F}.)
\end{corollary}

\begin{proof}
By Theorem~\ref{thm:cls}, the constrained estimator satisfies
\begin{equation}\label{eq:delta-beta}
\hat{\bb} - \hat{\bb}_R = (\bX'\bX)^{-1}\bR'\bigl[\bR(\bX'\bX)^{-1}\bR'\bigr]^{-1}(\bR\hat{\bb} - \br).
\end{equation}

\textbf{Step 1: Decompose $\SSE_R$.} Since $\bY - \bX\hat{\bb}_R = (\bY - \bX\hat{\bb}) + \bX(\hat{\bb} - \hat{\bb}_R) = \mathbf{e} + \bX(\hat{\bb} - \hat{\bb}_R)$,
\begin{align*}
\SSE_R = \|\bY - \bX\hat{\bb}_R\|^2 &= \|\mathbf{e}\|^2 + \|\bX(\hat{\bb} - \hat{\bb}_R)\|^2 + 2\mathbf{e}'\bX(\hat{\bb} - \hat{\bb}_R)\\
&= \SSE_U + \|\bX(\hat{\bb} - \hat{\bb}_R)\|^2,
\end{align*}
since the cross term vanishes: $\mathbf{e}'\bX = \bzero'$ by Theorem~\ref{thm:normal-eq}\ref{ne:i}.

\textbf{Step 2: Simplify the squared norm.} Let $\ba = [\bR(\bX'\bX)^{-1}\bR']^{-1}(\bR\hat{\bb} - \br)$ denote the trailing factor (a $q \times 1$ vector) in~\eqref{eq:delta-beta}, so that $\hat{\bb} - \hat{\bb}_R = (\bX'\bX)^{-1}\bR'\ba$. Then
\begin{align*}
\|\bX(\hat{\bb} - \hat{\bb}_R)\|^2 &= (\hat{\bb} - \hat{\bb}_R)'\bX'\bX(\hat{\bb} - \hat{\bb}_R)\\
&= \bigl[(\bX'\bX)^{-1}\bR'\ba\bigr]'\bX'\bX\bigl[(\bX'\bX)^{-1}\bR'\ba\bigr]\\
&= \ba'\bR\underbrace{(\bX'\bX)^{-1}(\bX'\bX)(\bX'\bX)^{-1}}_{=\,(\bX'\bX)^{-1}}\bR'\ba\\
&= \ba'\bR(\bX'\bX)^{-1}\bR'\ba.
\end{align*}
Expanding $\ba$:
\[
\|\bX(\hat{\bb} - \hat{\bb}_R)\|^2 = (\bR\hat{\bb} - \br)'\bigl[\bR(\bX'\bX)^{-1}\bR'\bigr]^{-1}\underbrace{\bR(\bX'\bX)^{-1}\bR'\bigl[\bR(\bX'\bX)^{-1}\bR'\bigr]^{-1}}_{=\,\bI_q}(\bR\hat{\bb} - \br).
\]
The bracketed product collapses to $\bI_q$, leaving
\begin{equation}\label{eq:sseR-sseU}
\SSE_R - \SSE_U = \|\bX(\hat{\bb} - \hat{\bb}_R)\|^2 = (\bR\hat{\bb} - \br)'\bigl[\bR(\bX'\bX)^{-1}\bR'\bigr]^{-1}(\bR\hat{\bb} - \br).
\end{equation}

\textbf{Step 3: Compare with Theorem~\ref{thm:general-F}.} The numerator in Theorem~\ref{thm:general-F}'s $F$ statistic equals the right-hand side of~\eqref{eq:sseR-sseU} divided by $q$. Since the denominator $s^2 = \SSE_U/(n-k-1)$ is the same, substituting yields
\[
F = \frac{(\SSE_R - \SSE_U)/q}{\SSE_U/(n-k-1)}. \qedhere
\]
\end{proof}

\begin{remark}[Logical status of the forward reference]\label{rem:forward-nestedF}
The proof relies on the constrained estimator formula~\eqref{eq:delta-beta} from Theorem~\ref{thm:cls} in Section~\ref{sec:special}; Theorem~\ref{thm:cls} is purely algebraic and does not depend on anything in this section, so no circularity arises. A reader seeking strict linearity can defer this corollary until after Section~\ref{sec:special}.

In the special case of \emph{homogeneous} constraints, $\br = \bzero$, there is also a derivation without constrained OLS: the feasible set $\{\bb : \bR\bb = \bzero\}$ is then a subspace, $\bX\{\bb : \bR\bb = \bzero\}$ is a subspace of $\calC(\bX)$ of dimension $k + 1 - q$ (by~\ref{L1} and~\ref{A4}), and $\SSE_R$ is the squared distance from $\bY$ to it. Writing $\bH_R$ for the orthogonal projection onto this subspace (Remark~\ref{rem:hat-projection}; it is symmetric idempotent with column space the subspace), $\bH\bH_R = \bH_R$ (the subspace lies in $\calC(\bX)$, on which $\bH$ acts as the identity), so $\bH - \bH_R$ is symmetric idempotent of rank $q$, $\bY = \bH_R\bY + (\bH - \bH_R)\bY + \bM\bY$ is an orthogonal decomposition, and $\SSE_R - \SSE_U = \bY'(\bH - \bH_R)\bY$ by Pythagoras; Theorem~\ref{thm:cls} with $\br = \bzero$ gives the explicit form $\bH_R = \bH - \bX(\bX'\bX)^{-1}\bR'[\bR(\bX'\bX)^{-1}\bR']^{-1}\bR(\bX'\bX)^{-1}\bX'$. For $\br \neq \bzero$ the feasible set is an affine set, not a subspace, and this identity is \emph{false}: $\bY'(\bH - \bH_R)\bY$ with the same $\bH_R$ equals $(\bR\hat{\bb})'[\bR(\bX'\bX)^{-1}\bR']^{-1}(\bR\hat{\bb})$, whereas the correct value~\eqref{eq:sseR-sseU} is centered at $\br$. The constrained-OLS route of the proof above covers general $\br$.
\end{remark}

\begin{corollary}[$t^2 = F$ for a Single Linear Constraint]\label{cor:t2F}
For testing $H_0\colon \bc'\bb = r_0$ (a single constraint, $q = 1$, $\bR = \bc'$ with $\bc \neq \bzero$), the $F$ statistic from Theorem~\ref{thm:general-F} equals the square of the $t$ statistic:
\[
t = \frac{\bc'\hat{\bb} - r_0}{s\sqrt{\bc'(\bX'\bX)^{-1}\bc}}, \qquad F = t^2,
\]
and $F = t^2 \sim F_{1, n-k-1}$ under $H_0$ by Theorem~\ref{thm:general-F}. The identity $F = t^2$ itself is algebraic: like Corollary~\ref{cor:nested-F}, it uses only~\ref{A4} (through $\bc'(\bX'\bX)^{-1}\bc > 0$ for $\bc \neq \bzero$, Theorem~\ref{thm:ols}, Step~4), and \ref{A1}--\ref{A5} enter only through the distribution.
\end{corollary}

\begin{proof}
With $q = 1$, $\bR = \bc'$, $\br = r_0$, the general $F$ statistic becomes:
\[
F = \frac{(\bc'\hat{\bb} - r_0)[\bc'(\bX'\bX)^{-1}\bc]^{-1}(\bc'\hat{\bb} - r_0)}{1 \cdot s^2} = \left(\frac{\bc'\hat{\bb} - r_0}{s\sqrt{\bc'(\bX'\bX)^{-1}\bc}}\right)^2 = t^2. \qedhere
\]
\end{proof}

\subsection{Confidence Intervals and Regions}

\begin{theorem}[Confidence Interval for $\bc'\bb$]\label{thm:ci}
Under~\ref{A1}--\ref{A5}, for $\bc \neq \bzero$ (so that $\bc'(\bX'\bX)^{-1}\bc > 0$ by Theorem~\ref{thm:ols}, Step~4), a $(1-\alpha) \times 100\%$ confidence interval for $\bc'\bb$ is
\begin{equation}
\bc'\hat{\bb} \pm t_{\alpha/2,\, n-k-1} \cdot s\sqrt{\bc'(\bX'\bX)^{-1}\bc},
\end{equation}
where $t_{\alpha/2, n-k-1}$ is the upper $\alpha/2$ quantile of $t_{n-k-1}$.
\end{theorem}

\begin{proof}
By Theorem~\ref{thm:mvn-linear} applied to $\hat{\bb} \sim N(\bb, \sigma^2(\bX'\bX)^{-1})$ with $\bA = \bc'$ (a $1 \times (k+1)$ row vector), $\bc'\hat{\bb} \sim N(\bc'\bb,\; \sigma^2\bc'(\bX'\bX)^{-1}\bc)$. Standardizing:
\[
Z = \frac{\bc'\hat{\bb} - \bc'\bb}{\sigma\sqrt{\bc'(\bX'\bX)^{-1}\bc}} \sim N(0, 1).
\]
By Theorem~\ref{thm:sse-dist}, $V = \SSE/\sigma^2 \sim \chi^2_{n-k-1}$ and is independent of $\hat{\bb}$; since $Z$ is a function of $\hat{\bb}$ alone, $Z$ is independent of $V$. By definition of the $t$ distribution:
\[
\frac{Z}{\sqrt{V/(n-k-1)}} = \frac{\bc'\hat{\bb} - \bc'\bb}{\sigma\sqrt{\bc'(\bX'\bX)^{-1}\bc}} \cdot \frac{1}{\sqrt{\SSE/[\sigma^2(n-k-1)]}} = \frac{\bc'\hat{\bb} - \bc'\bb}{s\sqrt{\bc'(\bX'\bX)^{-1}\bc}} \sim t_{n-k-1},
\]
where the last step uses $s^2 = \SSE/(n-k-1)$. The pivot gives, by the symmetry of the $t$ distribution about $0$, equation~\eqref{eq:t-symm},
\[
\prb\!\left(-t_{\alpha/2} \leqslant \frac{\bc'\hat{\bb} - \bc'\bb}{s\sqrt{\bc'(\bX'\bX)^{-1}\bc}} \leqslant t_{\alpha/2}\right) = 1 - \alpha,
\]
and rearranging for $\bc'\bb$ yields the interval.
\end{proof}

\begin{theorem}[Joint Confidence Ellipsoid]\label{thm:ellipsoid}
Under~\ref{A1}--\ref{A5}, a $(1-\alpha) \times 100\%$ joint confidence region for $\bb$ is
\begin{equation}
\bigl\{\bb \in \Real^{k+1} : (\hat{\bb} - \bb)'(\bX'\bX)(\hat{\bb} - \bb) \leqslant (k+1) s^2 F_{\alpha, k+1, n-k-1}\bigr\}.
\end{equation}
\end{theorem}

\begin{proof}
\textbf{Step 1: A chi-square quadratic form in $\hat{\bb} - \bb$.} By Theorem~\ref{thm:beta-dist},
\[
\hat{\bb} - \bb \sim N(\bzero, \sigma^2(\bX'\bX)^{-1}).
\]
Let $(\bX'\bX)^{1/2}$ be the positive definite symmetric square root (\ref{L4}) and define $\bw = \sigma^{-1}(\bX'\bX)^{1/2}(\hat{\bb} - \bb)$. Theorem~\ref{thm:mvn-linear} gives $\bw \sim N(\bzero, \bI_{k+1})$, so $\bw'\bw \sim \chi^2_{k+1}$ by Theorem~\ref{thm:quadratic-chisq}:
\[
\bw'\bw = \frac{(\hat{\bb} - \bb)'(\bX'\bX)(\hat{\bb} - \bb)}{\sigma^2} \sim \chi^2_{k+1}.
\]

\textbf{Step 2: Form the $F$ ratio.} This is a function of $\hat{\bb}$, hence independent of $\SSE/\sigma^2 \sim \chi^2_{n-k-1}$ (Theorem~\ref{thm:sse-dist}). By definition of $F$,
\[
\frac{(\hat{\bb} - \bb)'(\bX'\bX)(\hat{\bb} - \bb)/[(k+1)\sigma^2]}{\SSE/[\sigma^2(n-k-1)]} = \frac{(\hat{\bb} - \bb)'(\bX'\bX)(\hat{\bb} - \bb)}{(k+1)s^2} \sim F_{k+1, n-k-1}.
\]
The confidence region collects the $\bb$ for which this statistic is at most $F_{\alpha, k+1, n-k-1}$.
\end{proof}

\subsection{Testing Conventions and Further Consequences}\label{sec:test-more}

Throughout this subsection $\nu = n - k - 1$, $\bc \neq \bzero$, $s > 0$ (under~\ref{A5} this holds with probability one, since $\SSE/\sigma^2 \sim \chi^2_\nu$ has a continuous distribution function by~\ref{N5}), and $t$ denotes the statistic of Corollary~\ref{cor:t2F} for $H_0\colon \bc'\bb = r_0$,
\[
t = \frac{\bc'\hat{\bb} - r_0}{s\sqrt{\bc'(\bX'\bX)^{-1}\bc}},
\]
evaluated at the observed data; for $\bc = \boldsymbol\iota_j$ and $r_0 = \beta_j^0$ it is the statistic of Theorem~\ref{thm:t-test}. The proof of Theorem~\ref{thm:ci} shows that, under~\ref{A1}--\ref{A5}, $(\bc'\hat{\bb} - \bc'\bb)/(s\sqrt{\bc'(\bX'\bX)^{-1}\bc}) \sim t_\nu$; hence $t \sim t_\nu$ under $H_0$.

\begin{definition}[One- and Two-Sided $t$ Tests and $p$-Values]\label{def:t-tests}
Let $\alpha \in (0, 1)$ and $T \sim t_\nu$. The level-$\alpha$ $t$ tests of $H_0\colon \bc'\bb = r_0$ against the three alternatives, and their \textbf{$p$-values}, are:
\begin{center}
\begin{tabular}{@{}lll@{}}
\toprule
Alternative $H_1$ & Reject $H_0$ if & $p$-value\\
\midrule
$\bc'\bb > r_0$ & $t > t_{\alpha,\nu}$ & $\prb(T \geqslant t)$\\
$\bc'\bb < r_0$ & $t < -t_{\alpha,\nu}$ & $\prb(T \leqslant t)$\\
$\bc'\bb \neq r_0$ & $|t| > t_{\alpha/2,\nu}$ & $\prb(|T| \geqslant |t|)$\\
\bottomrule
\end{tabular}
\end{center}
In each $p$-value, $t$ is the observed value and $T$ is an independent $t_\nu$ variable: the $p$-value is the probability, under the null distribution of $t$, of a value at least as extreme, in the direction of $H_1$, as the one observed. The most common case is $H_0\colon \beta_j = 0$ against $\beta_j \neq 0$, with $t = \hat\beta_j/\SE(\hat\beta_j)$.
\end{definition}

\begin{proposition}[Exact Level and $p$-Values]\label{prop:t-pvalue}
Under~\ref{A1}--\ref{A5}:
\begin{enumerate}[label=\textup{(\roman*)}]
\item if $H_0\colon \bc'\bb = r_0$ holds, each of the three tests of Definition~\ref{def:t-tests} rejects $H_0$ with probability exactly $\alpha$; and, whether or not $H_0$ holds, the interval of Theorem~\ref{thm:ci} contains $\bc'\bb$ with probability exactly $1 - \alpha$;
\item each test rejects $H_0$ if and only if its $p$-value is smaller than $\alpha$;
\item the two-sided $p$-value equals $\prb(F \geqslant t^2)$ for $F \sim F_{1,\nu}$, the $p$-value of the $F$ test of Theorem~\ref{thm:general-F} with $\bR = \bc'$ and $\br = r_0$ (Corollary~\ref{cor:t2F}), where the $p$-value of an $F$ test with $q$ restrictions and observed statistic $F_{\mathrm{obs}}$ is $\prb(F \geqslant F_{\mathrm{obs}})$ for $F \sim F_{q,\nu}$, its null distribution.
\end{enumerate}
\end{proposition}

\begin{proof}
\textbf{(i)} Under $H_0$, $t \sim t_\nu$. By the definition of the upper quantile in~\ref{N5}, $\prb(t > t_{\alpha,\nu}) = \alpha$; by~\eqref{eq:t-symm}, $\prb(t < -t_{\alpha,\nu}) = \alpha$ and $\prb(|t| > t_{\alpha/2,\nu}) = 1 - \prb(-t_{\alpha/2,\nu} \leqslant t \leqslant t_{\alpha/2,\nu}) = \alpha$. The coverage probability is the display at the end of the proof of Theorem~\ref{thm:ci}, which holds for the true $\bc'\bb$.

\textbf{(ii)} Let $G(a) = \prb(T \geqslant a)$ for $a \in \Real$. By~\ref{N5}, the distribution function of $t_\nu$ is continuous and strictly increasing on $\Real$, so $G$ is continuous and strictly decreasing, with $G(t_{\alpha,\nu}) = \alpha$. \emph{Upper alternative:} $p = G(t) < \alpha = G(t_{\alpha,\nu})$ iff $t > t_{\alpha,\nu}$. \emph{Lower alternative:} since $-T \sim t_\nu$ (\S\ref{sec:prereq-prob}), $p = \prb(T \leqslant t) = \prb(-T \geqslant -t) = G(-t)$, which is $< \alpha$ iff $-t > t_{\alpha,\nu}$. \emph{Two-sided:} for $a \geqslant 0$, $\prb(|T| \geqslant a) = \prb(T \geqslant a) + \prb(T \leqslant -a) = 2G(a)$ (the two events are disjoint for $a > 0$, and for $a = 0$ both sides equal $1$ because $G(0) = \frac{1}{2}$ by symmetry and continuity), so $p = 2G(|t|) < \alpha$ iff $G(|t|) < \alpha/2 = G(t_{\alpha/2,\nu})$ iff $|t| > t_{\alpha/2,\nu}$.

\textbf{(iii)} $\prb(|T| \geqslant |t|) = \prb(T^2 \geqslant t^2)$, and $T^2 \sim F_{1,\nu}$ (\S\ref{sec:prereq-prob}). By Corollary~\ref{cor:t2F}, $t^2$ is the observed value of the statistic of Theorem~\ref{thm:general-F} with $\bR = \bc'$, $\br = r_0$, whose null distribution is $F_{1,\nu}$.
\end{proof}

\begin{proposition}[Duality of Tests and Confidence Intervals]\label{prop:duality}
Assume $s > 0$ (under~\ref{A5} this holds with probability one, since $\SSE/\sigma^2 \sim \chi^2_\nu$ has a continuous distribution function by~\ref{N5}). For every $r_0 \in \Real$, the two-sided level-$\alpha$ test of $H_0\colon \bc'\bb = r_0$ rejects if and only if $r_0$ lies outside the $(1-\alpha) \times 100\%$ confidence interval of Theorem~\ref{thm:ci}. Hence that interval is the set of hypothesized values $r_0$ that the test does not reject; in particular, $H_0\colon \beta_j = 0$ is rejected at level $\alpha$ iff the interval for $\beta_j$ excludes $0$.
\end{proposition}

\begin{proof}
Write $c = t_{\alpha/2,\nu}$ and $\SE = s\sqrt{\bc'(\bX'\bX)^{-1}\bc}$, which is positive because $s > 0$ and $\bc'(\bX'\bX)^{-1}\bc > 0$ (Theorem~\ref{thm:ols}, Step~4). The test rejects iff $|\bc'\hat{\bb} - r_0| > c\,\SE$, i.e.\ iff $r_0 < \bc'\hat{\bb} - c\,\SE$ or $r_0 > \bc'\hat{\bb} + c\,\SE$, i.e.\ iff $r_0 \notin [\bc'\hat{\bb} - c\,\SE,\; \bc'\hat{\bb} + c\,\SE]$, which is the interval of Theorem~\ref{thm:ci}. The last statement is the case $\bc = \boldsymbol\iota_j$, $r_0 = 0$.
\end{proof}

\begin{proposition}[Expected Regression Mean Square]\label{prop:E-SSR}
Write $\bX = [\bone \;\; \bX_*]$ and $\bb = (\beta_0, \bb_*')'$, where $\bb_* = (\beta_1, \ldots, \beta_k)'$ collects the slopes. Under~\ref{A1}--\ref{A3} (normality is not needed),
\[
\expc[\SSR] = k\sigma^2 + \|\bM_0\bX_*\bb_*\|^2, \qquad\text{so}\qquad \expc[\MSR] = \sigma^2 + \frac{\|\bM_0\bX_*\bb_*\|^2}{k} \;\geqslant\; \sigma^2 = \expc[\MSE],
\]
with equality if and only if $\bb_* = \bzero$, i.e.\ iff the null hypothesis of the overall $F$ test (Theorem~\ref{thm:f-test}) holds.
\end{proposition}

\begin{proof}
\textbf{Step 1: $\SSR$ as a quadratic form.} Step~1 of the proof of Lemma~\ref{lem:SSR-chisq} shows, without using $H_0$ or~\ref{A5}, that $\SSR = \bY'(\bH - \bJ)'(\bH - \bJ)\bY$ and $(\bH - \bJ)\bone = \bzero$; Step~2 shows, again without them, that $\bH - \bJ$ is symmetric idempotent of trace $k$. Hence $\SSR = \bY'(\bH - \bJ)\bY$.

\textbf{Step 2: Expectation.} By~\ref{A1} and Lemma~\ref{lem:var-linear} (with~\ref{A3}), $\expc[\bY] = \bX\bb$ and $\var(\bY) = \sigma^2\bI$, so Lemma~\ref{lem:quadratic-form} gives
\[
\expc[\SSR] = \tr\bigl((\bH - \bJ)\sigma^2\bI\bigr) + \bb'\bX'(\bH - \bJ)\bX\bb = k\sigma^2 + \|(\bH - \bJ)\bX\bb\|^2,
\]
using $\bH - \bJ = (\bH - \bJ)'(\bH - \bJ)$.

\textbf{Step 3: The non-centrality term.} Since $(\bH - \bJ)\bone = \bzero$ and $\bH\bX_* = \bX_*$ (Theorem~\ref{thm:HM}\ref{HM:ii}), $(\bH - \bJ)\bX\bb = (\bH - \bJ)\bX_*\bb_* = \bX_*\bb_* - \bJ\bX_*\bb_* = \bM_0\bX_*\bb_*$. This vanishes iff $\bX_*\bb_* = a\bone$ for some $a$, i.e.\ iff $\bX(-a, \bb_*')' = \bzero$, which by~\ref{A4} forces $a = 0$ and $\bb_* = \bzero$ (and conversely $\bb_* = \bzero$ gives $\bM_0\bX_*\bb_* = \bzero$). Dividing by $k$ and recalling $\expc[\MSE] = \sigma^2$ (Theorem~\ref{thm:mse-unbiased}) completes the proof.
\end{proof}

The overall $F$ statistic is thus a ratio of two estimators of $\sigma^2$ when $H_0$ holds, and its numerator has the larger expectation otherwise; values of $F$ well above $1$ are evidence against $H_0$, and how far above $1$ is decided by the $F_{k,\nu}$ quantile.

\begin{corollary}[$R^2$ from the Overall $F$ Statistic]\label{cor:R2-F}
If $\SST > 0$ and $\SSE > 0$, the overall $F$ statistic of Theorem~\ref{thm:f-test} determines $R^2$:
\[
R^2 = \frac{kF}{kF + n - k - 1}.
\]
Under~\ref{A5} both conditions hold with probability one.
\end{corollary}

\begin{proof}
$\SSE > 0$ gives $R^2 < 1$, so by~\eqref{eq:f-test}, $F = (n-k-1)R^2/\{k(1 - R^2)\}$, i.e.\ $kF(1 - R^2) = (n-k-1)R^2$, i.e.\ $R^2(kF + n - k - 1) = kF$; the bracket is positive since $F \geqslant 0$ and $n > k + 1$. Under~\ref{A5}, $\SST > 0$ with probability one by the proof of Theorem~\ref{thm:f-test}, and $\prb(\SSE = 0) = 0$ because $\SSE/\sigma^2 \sim \chi^2_\nu$ (Theorem~\ref{thm:sse-dist}) has a continuous distribution function (\ref{N5}).
\end{proof}

\begin{corollary}[Nested $F$ Test via $R^2$]\label{cor:nested-R2}
In the setting of Corollary~\ref{cor:nested-F}, assume $\SST > 0$ and $\SSE_U > 0$, and put $R^2_U = 1 - \SSE_U/\SST$ and $R^2_R = 1 - \SSE_R/\SST$. Then
\[
F = \frac{(R^2_U - R^2_R)/q}{(1 - R^2_U)/(n-k-1)}.
\]
If the constraints set $q$ of the coefficients to zero ($\bR$ consists of the corresponding rows of $\bI_{k+1}$, $\br = \bzero$), then $\SSE_R$ is the residual sum of squares of the OLS regression of $\bY$ on the remaining columns of $\bX$; if these include the intercept column (as when only slopes are set to zero), $R^2_R$ is the coefficient of determination (Definition~\ref{def:R2}) of that regression.
\end{corollary}

\begin{proof}
Both fits have the same response, hence the same $\SST$; so $\SSE_R - \SSE_U = \SST(R^2_U - R^2_R)$ and $\SSE_U = \SST(1 - R^2_U)$, and substituting into Corollary~\ref{cor:nested-F} and cancelling $\SST$ gives the formula. For the second statement, $\bX\bb$ with $\beta_j = 0$ for the $q$ constrained indices is $\bX_R\bb_R$, where $\bX_R$ consists of the remaining columns and $\bb_R$ of the remaining coefficients; minimizing $\|\bY - \bX\bb\|^2$ under the constraints is therefore minimizing $\|\bY - \bX_R\bb_R\|^2$ over $\bb_R$, the OLS problem for the design $\bX_R$, which has full column rank by~\ref{A4}; when $\bX_R$ contains the intercept column, $1 - \SSE_R/\SST$ is its $R^2$.
\end{proof}

For the next result and for Section~\ref{sec:fwl}, write $\bX_{\cdot j}$ for the column of $\bX$ that multiplies $\beta_j$ ($j = 0, \ldots, k$; $\bX_{\cdot 0} = \bone$), $\bX_{(j)}$ for the $n \times k$ matrix obtained from $\bX$ by deleting this column (here $k \geqslant 1$), which has full column rank by~\ref{A4}, $\bM_{(j)}$ for the residual maker of $\bX_{(j)}$, and $\SSE_{(j)}$ for the residual sum of squares of the regression of $\bY$ on $\bX_{(j)}$.

\begin{corollary}[Dropping One Regressor: $F = t^2$]\label{cor:drop-one}
Let $k \geqslant 1$, fix $j \in \{0, \ldots, k\}$, assume $\SSE > 0$, and let $t_j = \hat\beta_j/\SE(\hat\beta_j)$ be the statistic of Theorem~\ref{thm:t-test} for $\beta_j^0 = 0$. Then
\[
\SSE_{(j)} - \SSE = \hat\beta_j^2\,\bX_{\cdot j}'\bM_{(j)}\bX_{\cdot j}, \qquad [(\bX'\bX)^{-1}]_{jj} = \frac{1}{\bX_{\cdot j}'\bM_{(j)}\bX_{\cdot j}}, \qquad \frac{(\SSE_{(j)} - \SSE)/1}{\SSE/(n-k-1)} = t_j^2.
\]
Under~\ref{A1}--\ref{A5} and $H_0\colon \beta_j = 0$, this statistic has the $F_{1,n-k-1}$ distribution. For $k = 1$ and $j = 1$, it is the overall $F$ statistic of Theorem~\ref{thm:f-test}, so that $F = t_1^2$.
\end{corollary}

\begin{proof}
\textbf{Step 1: Reordering the columns.} Let $\bP$ be the permutation matrix that moves column $j$ to the last position, so that $\tilde{\bX} = \bX\bP = [\bX_{(j)} \;\; \bX_{\cdot j}]$. Since $\bP'\bP = \bI$, $(\tilde{\bX}'\tilde{\bX})^{-1} = \bP'(\bX'\bX)^{-1}\bP$, and the OLS estimator for $\tilde{\bX}$ is $\bP'(\bX'\bX)^{-1}\bP\bP'\bX'\bY = \bP'\hat{\bb}$, the same coefficients in the new order; the fitted values $\tilde{\bX}\bP'\hat{\bb} = \bX\hat{\bb}$, the residuals and $\SSE$ are unchanged; and the last diagonal entry of $(\tilde{\bX}'\tilde{\bX})^{-1}$ is $[(\bX'\bX)^{-1}]_{jj}$.

\textbf{Step 2: The identities.} By~\ref{A4}, $\bX_{\cdot j} \notin \calC(\bX_{(j)})$, so Lemma~\ref{lem:sse-reduction} applies with $\bX_{(j)}$ in place of $\bX_k$ and $\bX_{\cdot j}$ in place of $\bx$; it gives the first identity, and identity~\eqref{eq:last-diag}, whose proof uses Lemma~\ref{lem:sse-reduction} alone, gives the second. Hence $\SE(\hat\beta_j)^2 = s^2/(\bX_{\cdot j}'\bM_{(j)}\bX_{\cdot j})$, and
\[
\frac{\SSE_{(j)} - \SSE}{s^2} = \frac{\hat\beta_j^2\,\bX_{\cdot j}'\bM_{(j)}\bX_{\cdot j}}{s^2} = \frac{\hat\beta_j^2}{\SE(\hat\beta_j)^2} = t_j^2.
\]
(Equivalently, by Corollary~\ref{cor:nested-R2} the regression on $\bX_{(j)}$ is the constrained fit for $\bR = \boldsymbol\iota_j'$, $\br = 0$, so the identity also follows from Corollaries~\ref{cor:nested-F} and~\ref{cor:t2F}.)

\textbf{Step 3: Distribution.} Under $H_0$, $t_j \sim t_{n-k-1}$ by Theorem~\ref{thm:t-test}, so $t_j^2 \sim F_{1,n-k-1}$ (\S\ref{sec:prereq-prob}).

\textbf{Step 4: The case $k = 1$.} Then $\bX_{(1)} = \bone$, whose fit is $\bar Y\bone$ with residual sum of squares $\SST$ (Remark~\ref{rem:special-cases}), so the statistic is $(\SST - \SSE)/(\SSE/(n-2)) = \SSR/(\SSE/(n-2))$ by Theorem~\ref{thm:anova}, which is $F$ of~\eqref{eq:f-test} with $k = 1$.
\end{proof}

%=============================================================================
\section{Partitioned Regression}\label{sec:fwl}
%=============================================================================

Theorem~\ref{thm:fwl} is an algebraic identity that uses only the standing assumption~\ref{A4}; Theorem~\ref{thm:ovb} uses~\ref{A2} in addition. The results that follow Theorem~\ref{thm:ovb} are algebraic except for their statements about expectations and variances, which use~\ref{A1}--\ref{A3} as stated in each. No distributional assumption is used in this section.

\subsection{The Frisch--Waugh--Lovell Theorem}

\begin{theorem}[Frisch--Waugh--Lovell]\label{thm:fwl}
Partition $\bX = [\bX_1 \;\; \bX_2]$ and $\bb = (\bb_1', \bb_2')'$, so the model is $\bY = \bX_1\bb_1 + \bX_2\bb_2 + \be$. Let $\bM_2 = \bI - \bX_2(\bX_2'\bX_2)^{-1}\bX_2'$ be the residual maker for $\bX_2$. Then:
\begin{enumerate}[label=\textup{(\roman*)}]
\item\label{fwl:i} $\hat{\bb}_1 = (\bX_1'\bM_2\bX_1)^{-1}\bX_1'\bM_2\bY$.
\item The OLS coefficient vector from regressing $\bM_2\bY$ on $\bM_2\bX_1$ is the $\hat{\bb}_1$ of~\ref{fwl:i}, and the residual vector of that regression equals the residual vector $\mathbf{e}$ from the full regression.
\end{enumerate}
\end{theorem}

\begin{proof}
\textbf{(i).} The normal equations for the full model are:
\[
\begin{pmatrix} \bX_1'\bX_1 & \bX_1'\bX_2 \\ \bX_2'\bX_1 & \bX_2'\bX_2 \end{pmatrix} \begin{pmatrix} \hat{\bb}_1 \\ \hat{\bb}_2 \end{pmatrix} = \begin{pmatrix} \bX_1'\bY \\ \bX_2'\bY \end{pmatrix}.
\]
The second block equation reads $\bX_2'\bX_1\hat{\bb}_1 + \bX_2'\bX_2\hat{\bb}_2 = \bX_2'\bY$, which rearranges to
\[
\hat{\bb}_2 = (\bX_2'\bX_2)^{-1}\bX_2'(\bY - \bX_1\hat{\bb}_1).
\]
Substituting this into the first block equation $\bX_1'\bX_1\hat{\bb}_1 + \bX_1'\bX_2\hat{\bb}_2 = \bX_1'\bY$:
\begin{equation}\label{eq:fwl-sub}
\bX_1'\bX_1\hat{\bb}_1 + \underbrace{\bX_1'\bX_2(\bX_2'\bX_2)^{-1}\bX_2'}_{=\,\bX_1'\bH_2}(\bY - \bX_1\hat{\bb}_1) = \bX_1'\bY,
\end{equation}
where we recognize the projection onto $\calC(\bX_2)$:
\[
\bH_2 := \bX_2(\bX_2'\bX_2)^{-1}\bX_2' = \bI - \bM_2.
\]
Expanding~\eqref{eq:fwl-sub} using $\bX_1'\bH_2 = \bX_1' - \bX_1'\bM_2$:
\[
\bX_1'\bX_1\hat{\bb}_1 + \bX_1'(\bY - \bX_1\hat{\bb}_1) - \bX_1'\bM_2(\bY - \bX_1\hat{\bb}_1) = \bX_1'\bY.
\]
The first two terms combine to $\bX_1'\bY$, cancelling the right-hand side, leaving $\bX_1'\bM_2(\bY - \bX_1\hat{\bb}_1) = \bzero$, or
\[
\bX_1'\bM_2\bX_1\hat{\bb}_1 = \bX_1'\bM_2\bY.
\]
Since $\bX_1'\bM_2\bX_1 = (\bM_2\bX_1)'(\bM_2\bX_1)$ is positive definite (the columns of $\bM_2\bX_1$ are linearly independent: if $\bM_2\bX_1\bv = \bzero$ then $\bX_1\bv = \bH_2\bX_1\bv \in \calC(\bX_2)$, say $\bX_1\bv = \bX_2\bw$, so $[\bX_1 \;\; \bX_2](\bv', -\bw')' = \bzero$ and~\ref{A4} forces $\bv = \bzero$):
\[
\hat{\bb}_1 = (\bX_1'\bM_2\bX_1)^{-1}\bX_1'\bM_2\bY.
\]

\textbf{(ii).} Let $\tilde{\bY} = \bM_2\bY$ and $\tilde{\bX}_1 = \bM_2\bX_1$. The columns of $\tilde{\bX}_1$ are linearly independent (shown in (i)), so the OLS coefficient vector from regressing $\tilde{\bY}$ on $\tilde{\bX}_1$ is defined, and, using $\bM_2' = \bM_2 = \bM_2^2$ (Theorem~\ref{thm:HM}\ref{HM:i} for the design $\bX_2$, which has full column rank by~\ref{A4}) and then (i),
\[
(\tilde{\bX}_1'\tilde{\bX}_1)^{-1}\tilde{\bX}_1'\tilde{\bY} = (\bX_1'\bM_2'\bM_2\bX_1)^{-1}\bX_1'\bM_2'\bM_2\bY = (\bX_1'\bM_2\bX_1)^{-1}\bX_1'\bM_2\bY = \hat{\bb}_1.
\]
The residual from this regression is therefore
\[
\tilde{\mathbf{e}} = \bM_2\bY - \bM_2\bX_1\hat{\bb}_1 = \bM_2(\bY - \bX_1\hat{\bb}_1).
\]
Substituting $\bY - \bX_1\hat{\bb}_1 = \bX_2\hat{\bb}_2 + \mathbf{e}$ from the full regression and using $\bM_2\bX_2 = \bzero$:
\begin{equation}
\tilde{\mathbf{e}} = \bM_2\bX_2\hat{\bb}_2 + \bM_2\mathbf{e} = \bM_2\mathbf{e}.
\end{equation}
Finally, $\bM_2\mathbf{e} = \mathbf{e}$: since $\bX'\mathbf{e} = \bzero$ (Theorem~\ref{thm:normal-eq}\ref{ne:i}) includes $\bX_2'\mathbf{e} = \bzero$,
\[
\bM_2\mathbf{e} = \mathbf{e} - \bX_2(\bX_2'\bX_2)^{-1}\bX_2'\mathbf{e} = \mathbf{e}.
\]
Hence $\tilde{\mathbf{e}} = \mathbf{e}$.
\end{proof}

\begin{remark}[The key identity $\bM_2\bM = \bM$]\label{rem:M2M}
The computation above rests on the identity $\bM_2\mathbf{e} = \mathbf{e}$, which can be written equivalently as $\bM_2\bM = \bM$ (applying both sides to $\bY$). This holds because $\calC(\bX_2) \subseteq \calC(\bX)$, so the orthogonal complement $\calC(\bX)^\perp \subseteq \calC(\bX_2)^\perp$: anything already orthogonal to $\bX$ is automatically orthogonal to $\bX_2$, and $\bM_2$ leaves it unchanged. This identity is the abstract reason why ``residualizing with respect to a subset of regressors'' preserves the full-model residual.
\end{remark}

\subsection{Omitted Variable Bias}

\begin{theorem}[Omitted Variable Bias]\label{thm:ovb}
Suppose the true model is $\bY = \bX_1\bb_1 + \bX_2\bb_2 + \be$ with $\expc[\be] = \bzero$ (assumption~\ref{A2}), but the researcher estimates the ``short'' regression $\bY = \bX_1\boldsymbol\gamma + \mathbf{u}$, obtaining $\tilde{\boldsymbol\gamma} = (\bX_1'\bX_1)^{-1}\bX_1'\bY$. Then:
\begin{equation}\label{eq:ovb}
\expc[\tilde{\boldsymbol\gamma}] = \bb_1 + \underbrace{(\bX_1'\bX_1)^{-1}\bX_1'\bX_2}_{\text{auxiliary regression coefficients}}\,\bb_2.
\end{equation}
The estimator $\tilde{\boldsymbol\gamma}$ is unbiased for $\bb_1$ if and only if $\bX_1'\bX_2\bb_2 = \bzero$, which holds when either $\bX_1'\bX_2 = \bzero$ ($\bX_1$ and $\bX_2$ are orthogonal) or $\bb_2 = \bzero$ (the omitted variables are irrelevant).
\end{theorem}

\begin{proof}
Substituting the true model $\bY = \bX_1\bb_1 + \bX_2\bb_2 + \be$ into the short-regression estimator:
\begin{align*}
\tilde{\boldsymbol\gamma} &= (\bX_1'\bX_1)^{-1}\bX_1'\bY = (\bX_1'\bX_1)^{-1}\bX_1'(\bX_1\bb_1 + \bX_2\bb_2 + \be)\\
&= \bb_1 + (\bX_1'\bX_1)^{-1}\bX_1'\bX_2\bb_2 + (\bX_1'\bX_1)^{-1}\bX_1'\be.
\end{align*}
Taking expectations (using $\expc[\be] = \bzero$):
\[
\expc[\tilde{\boldsymbol\gamma}] = \bb_1 + (\bX_1'\bX_1)^{-1}\bX_1'\bX_2\bb_2. \qedhere
\]
\end{proof}

\begin{remark}[Two-factor structure of the bias]\label{rem:ovb-structure}
The bias term $(\bX_1'\bX_1)^{-1}\bX_1'\bX_2\,\bb_2$ is the product of two conceptually distinct quantities:
\begin{itemize}
\item $(\bX_1'\bX_1)^{-1}\bX_1'\bX_2$ is the matrix of OLS coefficients from regressing each column of $\bX_2$ on $\bX_1$---an \emph{auxiliary regression} describing how the omitted variables relate to the included ones.
\item $\bb_2$ is the effect of the omitted variables on $\bY$ in the true model.
\end{itemize}
The bias vanishes if either factor is zero: if $\bX_1$ and $\bX_2$ are orthogonal ($\bX_1'\bX_2 = \bzero$), omission is harmless; if the omitted variables have no effect ($\bb_2 = \bzero$), there is nothing to omit. Both ingredients must be nonzero to generate bias.
\end{remark}

\begin{remark}[Connection to FWL]\label{rem:ovb-fwl}
Theorem~\ref{thm:fwl} shows that the correct (full-model) coefficient $\hat{\bb}_1$ is obtained by regressing $\bM_2\bY$ on $\bM_2\bX_1$---that is, after partialling $\bX_2$ out of both $\bY$ and $\bX_1$. The short regression instead regresses $\bY$ directly on $\bX_1$, effectively replacing $\bM_2$ with $\bI$. The two coefficient vectors differ by an exact sample identity: the first block of the normal equations of the full model, $\bX_1'\bX_1\hat{\bb}_1 + \bX_1'\bX_2\hat{\bb}_2 = \bX_1'\bY = \bX_1'\bX_1\tilde{\boldsymbol\gamma}$, gives
\[
\tilde{\boldsymbol\gamma} - \hat{\bb}_1 = (\bX_1'\bX_1)^{-1}\bX_1'\bX_2\,\hat{\bb}_2,
\]
which has the same form as the bias term in~\eqref{eq:ovb} with $\hat{\bb}_2$ in place of $\bb_2$; taking expectations, with $\expc[\hat{\bb}_1] = \bb_1$ and $\expc[\hat{\bb}_2] = \bb_2$ (Theorem~\ref{thm:unbiased}, under~\ref{A2}), recovers~\eqref{eq:ovb}. OVB measures the cost of not partialling out the omitted variables.
\end{remark}

\begin{theorem}[Omitted Variables: Fitted Values and Error Variance]\label{thm:ovb-fit}
In the setting of Theorem~\ref{thm:ovb}, assume~\ref{A1}--\ref{A3} for the true model $\bY = \bX_1\bb_1 + \bX_2\bb_2 + \be$, whose design $\bX = [\bX_1 \;\; \bX_2]$ satisfies~\ref{A4}, and let $\bX_1$ have $p < n$ columns. Let $\bH_1 = \bX_1(\bX_1'\bX_1)^{-1}\bX_1'$ and $\bM_1 = \bI - \bH_1$ be the hat matrix and residual maker of the short regression, $\hat{\bY}_1 = \bH_1\bY$ its fitted values and $\tilde s^2 = \bY'\bM_1\bY/(n-p)$ its mean squared error. Then:
\begin{enumerate}[label=\textup{(\roman*)}]
\item the fitted values are biased: $\expc[\hat{\bY}_1] - \expc[\bY] = -\bM_1\bX_2\bb_2$;
\item\label{ovbfit:ii} $\ds\expc[\tilde s^2] = \sigma^2 + \frac{\|\bM_1\bX_2\bb_2\|^2}{n-p} = \sigma^2 + \frac{1}{n-p}\sum_{i=1}^n\bigl(\expc[\hat Y_{1i}] - \expc[Y_i]\bigr)^2 \;\geqslant\; \sigma^2$, with equality if and only if $\bb_2 = \bzero$.
\end{enumerate}
\end{theorem}

\begin{proof}
The columns of $\bX_1$ are among those of $\bX$, so $\bX_1$ has full column rank and Theorem~\ref{thm:HM} applies to $\bH_1$ and $\bM_1$ (with $p$ in place of $k + 1$).

\textbf{(i)} By~\ref{A1}, $\expc[\bY] = \bX_1\bb_1 + \bX_2\bb_2$, and $\bH_1\bX_1 = \bX_1$, so $\expc[\bH_1\bY] = \bX_1\bb_1 + \bH_1\bX_2\bb_2$; subtracting $\expc[\bY]$ leaves $-(\bI - \bH_1)\bX_2\bb_2 = -\bM_1\bX_2\bb_2$.

\textbf{(ii)} By Lemma~\ref{lem:quadratic-form} with $\expc[\bY] = \bX\bb$ and $\var(\bY) = \sigma^2\bI$ (Lemma~\ref{lem:var-linear}, \ref{A3}),
\[
\expc[\bY'\bM_1\bY] = \sigma^2\tr(\bM_1) + (\bX_1\bb_1 + \bX_2\bb_2)'\bM_1(\bX_1\bb_1 + \bX_2\bb_2) = \sigma^2(n - p) + \|\bM_1\bX_2\bb_2\|^2,
\]
using $\tr(\bM_1) = n - p$, $\bM_1\bX_1 = \bzero$ and $\bM_1 = \bM_1'\bM_1$. Dividing by $n - p$ gives the first expression, and by (i) $\|\bM_1\bX_2\bb_2\|^2$ is the sum of the squared biases of the fitted values. It vanishes iff $\bX_2\bb_2 = \bH_1\bX_2\bb_2$, i.e.\ iff $\bX_2\bb_2 = \bX_1\ba$ for some $\ba$, i.e.\ iff $\bX(-\ba', \bb_2')' = \bzero$, which by~\ref{A4} holds iff $\ba = \bzero$ and $\bb_2 = \bzero$.
\end{proof}

\begin{proposition}[Extraneous Variables]\label{prop:extraneous}
Suppose the true model is $\bY = \bX_1\bb_1 + \be$ with~\ref{A1}--\ref{A3}, where $\bX_1$ has $p$ columns, but the long regression on $\bX = [\bX_1 \;\; \bX_2]$ is fitted, where $\bX$ has full column rank and $n > k + 1$, giving $\hat{\bb} = (\hat{\bb}_1', \hat{\bb}_2')'$ and $s^2$. Then $\expc[\hat{\bb}_1] = \bb_1$, $\expc[\hat{\bb}_2] = \bzero$, $\expc[\hat{\bY}] = \expc[\bY]$ and $\expc[s^2] = \sigma^2$; and $\var(\hat{\bb}_1) - \var(\tilde{\boldsymbol\gamma})$ is positive semi-definite, where $\tilde{\boldsymbol\gamma} = (\bX_1'\bX_1)^{-1}\bX_1'\bY$ is the OLS estimator of the true model.
\end{proposition}

\begin{proof}
The true model is the model with design $\bX$ and coefficient vector $\bb = (\bb_1', \bzero')'$, and~\ref{A1}--\ref{A4} hold for it. Theorem~\ref{thm:unbiased} gives $\expc[\hat{\bb}] = (\bb_1', \bzero')'$, hence $\expc[\hat{\bY}] = \bX\expc[\hat{\bb}] = \bX_1\bb_1 = \expc[\bY]$, and Theorem~\ref{thm:mse-unbiased} gives $\expc[s^2] = \sigma^2$. For the variances, $\hat{\bb}_1 = \bC\bY$, where $\bC$ consists of the first $p$ rows of $(\bX'\bX)^{-1}\bX'$; since $(\bX'\bX)^{-1}\bX'\bX = \bI_{k+1}$, $\bC\bX_1 = \bI_p$, so $\expc[\hat{\bb}_1] = \bC\bX_1\bb_1 = \bb_1$ for every $\bb_1$: $\hat{\bb}_1$ is a linear unbiased estimator in the true model, whose design $\bX_1$ has full column rank (its columns are among those of $\bX$). The Gauss--Markov theorem (Theorem~\ref{thm:gm}) applied to the true model, with $\bX_1$ in place of $\bX$, gives $\var(\hat{\bb}_1) - \var(\tilde{\boldsymbol\gamma}) \succeq \bzero$.
\end{proof}

\begin{remark}[Omitting versus including]\label{rem:omit-include}
Theorems~\ref{thm:ovb} and~\ref{thm:ovb-fit} and Proposition~\ref{prop:extraneous} compare the two errors of specification. Omitting a relevant regressor biases the coefficients (unless $\bX_1'\bX_2\bb_2 = \bzero$) and the fitted values, and inflates $\expc[\tilde s^2]$ by the sum of the squared biases of the fitted values divided by $n - p$ (Theorem~\ref{thm:ovb-fit}\ref{ovbfit:ii}), which tends to widen intervals computed from $\tilde s$. Including an irrelevant regressor keeps every estimator unbiased and costs only precision.
\end{remark}

\subsection{Added-Variable Plots, Partial Correlation, and Variance Inflation}

In this subsection $j \in \{1, \ldots, k\}$, and $\bX_{\cdot j}$, $\bX_{(j)}$, $\bM_{(j)}$ and $\SSE_{(j)}$ are as defined before Corollary~\ref{cor:drop-one}; $\bX_{(j)}$ contains the intercept column $\bone$.

\begin{definition}[Added-Variable Plot and Partial Correlation]\label{def:avp}
Let $\mathbf{e}_{(j)} = \bM_{(j)}\bY$ be the residuals of the regression of $\bY$ on the other columns, and $\tilde{\bx}_j = \bM_{(j)}\bX_{\cdot j}$ the residuals of the regression of the $j$-th regressor on them. The \textbf{added-variable plot} (or partial regression plot) for regressor $j$ is the scatter plot of the points $(\tilde x_{ji}, e_{(j)i})$, $i = 1, \ldots, n$. When $\mathbf{e}_{(j)}$ is not constant, the \textbf{partial correlation} of $\bY$ and regressor $j$ given the other regressors is $r_{Yj\cdot} = \corr(\tilde{\bx}_j, \mathbf{e}_{(j)})$ (Definition~\ref{def:corr}; $\tilde{\bx}_j$ is never constant, by Corollary~\ref{cor:avp}\ref{avp:i} below).
\end{definition}

\begin{corollary}[What the Added-Variable Plot Shows]\label{cor:avp}
\leavevmode
\begin{enumerate}[label=\textup{(\roman*)}]
\item\label{avp:i} $\mathbf{e}_{(j)}$ and $\tilde{\bx}_j$ have mean zero, and $\tilde{\bx}_j$ is not constant.
\item\label{avp:ii} The OLS line through the added-variable plot (the regression of $\mathbf{e}_{(j)}$ on $[\bone \;\; \tilde{\bx}_j]$) has intercept $0$ and slope
\[
\frac{\tilde{\bx}_j'\mathbf{e}_{(j)}}{\tilde{\bx}_j'\tilde{\bx}_j} = \frac{\bX_{\cdot j}'\bM_{(j)}\bY}{\bX_{\cdot j}'\bM_{(j)}\bX_{\cdot j}} = \hat\beta_j,
\]
the coefficient of regressor $j$ in the full regression.
\item The residuals of that line are the residuals $\mathbf{e}$ of the full regression; in particular its residual sum of squares is $\SSE$.
\item\label{avp:iv} Under~\ref{A3}, $\var(\hat\beta_j) = \sigma^2/(\bX_{\cdot j}'\bM_{(j)}\bX_{\cdot j}) = \sigma^2/\|\tilde{\bx}_j\|^2$.
\end{enumerate}
\end{corollary}

\begin{proof}
\textbf{(i)} Since $\bone \in \calC(\bX_{(j)})$, $\bM_{(j)}\bone = \bzero$ (Theorem~\ref{thm:HM}\ref{HM:ii} for the design $\bX_{(j)}$), and by symmetry $\bone'\bM_{(j)} = \bzero'$; hence $\bone'\mathbf{e}_{(j)} = 0$ and $\bone'\tilde{\bx}_j = 0$. If $\tilde{\bx}_j = \bzero$, then $\bX_{\cdot j} = (\bI - \bM_{(j)})\bX_{\cdot j} \in \calC(\bX_{(j)})$, contradicting~\ref{A4}; a nonzero vector with mean zero is not constant.

\textbf{(ii)} By (i), Corollary~\ref{cor:slr} applies to the regression of $\mathbf{e}_{(j)}$ on $[\bone \;\; \tilde{\bx}_j]$. Because both vectors have mean zero, its slope is $S_{\tilde x e}/S_{\tilde x\tilde x} = \tilde{\bx}_j'\mathbf{e}_{(j)}/\tilde{\bx}_j'\tilde{\bx}_j$ and its intercept is $0 - \text{slope}\cdot 0 = 0$. Since $\bM_{(j)}$ is symmetric idempotent, $\tilde{\bx}_j'\mathbf{e}_{(j)} = \bX_{\cdot j}'\bM_{(j)}\bY$ and $\tilde{\bx}_j'\tilde{\bx}_j = \bX_{\cdot j}'\bM_{(j)}\bX_{\cdot j}$. By Theorem~\ref{thm:fwl}\ref{fwl:i} with $\bX_1 = \bX_{\cdot j}$ and $\bX_2 = \bX_{(j)}$ (after reordering the columns, as in Step~1 of the proof of Corollary~\ref{cor:drop-one}), this ratio is $\hat\beta_j$.

\textbf{(iii)} As the intercept is $0$, the line's fitted values are $\hat\beta_j\tilde{\bx}_j$ and its residuals are $\bM_{(j)}\bY - \bM_{(j)}\bX_{\cdot j}\hat\beta_j$, the residuals of the regression of $\bM_{(j)}\bY$ on $\bM_{(j)}\bX_{\cdot j}$; by Theorem~\ref{thm:fwl}(ii) they equal $\mathbf{e}$.

\textbf{(iv)} By (ii), $\hat\beta_j = \ba'\bY$ with the constant vector $\ba = \tilde{\bx}_j/\|\tilde{\bx}_j\|^2$ (because $\tilde{\bx}_j'\bY = \bX_{\cdot j}'\bM_{(j)}\bY$), so Lemma~\ref{lem:var-linear} and~\ref{A3} give $\var(\hat\beta_j) = \sigma^2\ba'\ba = \sigma^2/\|\tilde{\bx}_j\|^2$. (This agrees with $\sigma^2[(\bX'\bX)^{-1}]_{jj}$ by Corollary~\ref{cor:drop-one}.)
\end{proof}

\begin{proposition}[Partial Correlation and the $t$ Statistic]\label{prop:partial-t}
Assume $\SSE > 0$ and let $t_j$ be as in Corollary~\ref{cor:drop-one}. Then $\mathbf{e}_{(j)}$ is not constant, and
\[
r_{Yj\cdot}^2 = \frac{\SSE_{(j)} - \SSE}{\SSE_{(j)}}, \qquad r_{Yj\cdot} = \frac{t_j}{\sqrt{t_j^2 + n - k - 1}}.
\]
\end{proposition}

\begin{proof}
\textbf{Step 1: The correlation.} By Corollary~\ref{cor:drop-one}, $\SSE_{(j)} \geqslant \SSE > 0$, and $\SSE_{(j)} = \|\mathbf{e}_{(j)}\|^2$; so $\mathbf{e}_{(j)} \neq \bzero$, and having mean zero (Corollary~\ref{cor:avp}\ref{avp:i}) it is not constant. Since both vectors have mean zero, $r_{Yj\cdot} = \tilde{\bx}_j'\mathbf{e}_{(j)}/(\|\tilde{\bx}_j\|\,\|\mathbf{e}_{(j)}\|)$.

\textbf{Step 2: The first identity.} By Corollaries~\ref{cor:drop-one} and~\ref{cor:avp}\ref{avp:ii}, $\SSE_{(j)} - \SSE = \hat\beta_j^2\|\tilde{\bx}_j\|^2$, which equals $(\tilde{\bx}_j'\mathbf{e}_{(j)})^2/\|\tilde{\bx}_j\|^2$. Dividing by $\SSE_{(j)} = \|\mathbf{e}_{(j)}\|^2$ gives $r_{Yj\cdot}^2$.

\textbf{Step 3: The second identity.} Write $\Delta = \SSE_{(j)} - \SSE$. By Corollary~\ref{cor:drop-one}, $\Delta/\SSE = t_j^2/(n-k-1)$, so
\[
r_{Yj\cdot}^2 = \frac{\Delta}{\SSE + \Delta} = \frac{\Delta/\SSE}{1 + \Delta/\SSE} = \frac{t_j^2}{t_j^2 + n - k - 1}.
\]
The sign of $r_{Yj\cdot}$ is that of $\tilde{\bx}_j'\mathbf{e}_{(j)}$, which by Corollary~\ref{cor:avp}\ref{avp:ii} is the sign of $\hat\beta_j$ and hence of $t_j$ (as $\SE(\hat\beta_j) > 0$). Taking the square root with this sign gives the second identity.
\end{proof}

\begin{corollary}[Variance Inflation Factor]\label{cor:vif}
Let $R_j^2$ be the coefficient of determination of the regression of $\bX_{\cdot j}$ on $\bX_{(j)}$, and $S_{jj} = \sum_i(X_{ij} - \bar X_j)^2 = (n-1)s_j^2$, where $\bar X_j$ and $s_j^2$ are the mean and sample variance of the $j$-th regressor. Then $S_{jj} > 0$ and $R_j^2 < 1$, and under~\ref{A3}
\[
\var(\hat\beta_j) = \frac{\sigma^2}{S_{jj}(1 - R_j^2)} = \frac{\sigma^2\,\VIF_j}{(n-1)s_j^2}, \qquad \VIF_j = \frac{1}{1 - R_j^2},
\]
where $\VIF_j$ is the \textbf{variance inflation factor} of regressor $j$; correspondingly,
\[
\SE(\hat\beta_j) = \frac{s\sqrt{\VIF_j}}{s_j\sqrt{n-1}}.
\]
For $k = 1$, $R_1^2 = 0$ and the formula reduces to $\var(\hat\beta_1) = \sigma^2/S_{xx}$ of Corollary~\ref{cor:slr-var}.
\end{corollary}

\begin{proof}
$\bX_{\cdot j}$ is not constant, for otherwise $\bX_{\cdot j} \in \calC(\bone) \subseteq \calC(\bX_{(j)})$, contradicting~\ref{A4}; so $S_{jj} > 0$, and $S_{jj}$ is the total sum of squares of the auxiliary regression of $\bX_{\cdot j}$ on $\bX_{(j)}$, which has an intercept. Its residual sum of squares is $\|\bM_{(j)}\bX_{\cdot j}\|^2 = \|\tilde{\bx}_j\|^2 > 0$ (Corollary~\ref{cor:avp}\ref{avp:i}), so $R_j^2 = 1 - \|\tilde{\bx}_j\|^2/S_{jj} < 1$ (Definition~\ref{def:R2}) and $\|\tilde{\bx}_j\|^2 = S_{jj}(1 - R_j^2)$. Substituting into Corollary~\ref{cor:avp}\ref{avp:iv} gives $\var(\hat\beta_j)$; replacing $\sigma^2[(\bX'\bX)^{-1}]_{jj} = \sigma^2/\|\tilde{\bx}_j\|^2$ by $s^2/\|\tilde{\bx}_j\|^2$ gives $\SE(\hat\beta_j)^2$ (Theorem~\ref{thm:t-test} and Corollary~\ref{cor:drop-one}). For $k = 1$, $\bX_{(1)} = \bone$, whose fitted values are $\bar x\bone$, so the residual sum of squares is $S_{xx}$ and $R_1^2 = 0$.
\end{proof}

\begin{remark}[Sources of imprecision]\label{rem:imprecision}
By Corollary~\ref{cor:vif}, $\var(\hat\beta_j)$ is large when the error variance $\sigma^2$ is large, when the $j$-th regressor varies little in the sample ($S_{jj} = (n-1)s_j^2$ small, through a small $n$ or a small spread $s_j$), or when it is nearly a linear combination of the other regressors ($R_j^2$ close to $1$, i.e.\ a large $\VIF_j$: collinearity). $\VIF_j \geqslant 1$, with equality iff $R_j^2 = 0$, i.e.\ iff the fitted values of the auxiliary regression are constant.
\end{remark}

%=============================================================================
\section{Prediction}\label{sec:pred}
%=============================================================================

Let $\bx_0 \in \Real^{k+1}$ be a new observation's regressor vector. The predicted value is $\hat{Y}_0 = \bx_0'\hat{\bb}$.

\begin{theorem}[Confidence Interval for the Conditional Mean]\label{thm:ci-mean}
Under~\ref{A1}--\ref{A5}, a $(1-\alpha) \times 100\%$ confidence interval for $\expc[Y_0 \mid \bx_0] = \bx_0'\bb$ is
\begin{equation}
\bx_0'\hat{\bb} \pm t_{\alpha/2,\,n-k-1} \cdot s\sqrt{\bx_0'(\bX'\bX)^{-1}\bx_0}.
\end{equation}
\end{theorem}

\begin{proof}
This is the special case $\bc = \bx_0$ of Theorem~\ref{thm:ci}; its hypothesis $\bc \neq \bzero$ holds because the first component of $\bx_0$ is the intercept's $1$ (Definition~\ref{def:model}).
\end{proof}

\begin{theorem}[Prediction Interval for a New Observation]\label{thm:pred}
Under~\ref{A1}--\ref{A5}, let $Y_0 = \bx_0'\bb + \varepsilon_0$ be a new observation whose error $\varepsilon_0 \sim N(0, \sigma^2)$ is independent of $\be$ (that is, \ref{A3} and~\ref{A5} extend to the new observation). Then a $(1-\alpha) \times 100\%$ prediction interval for $Y_0$ is
\begin{equation}
\bx_0'\hat{\bb} \pm t_{\alpha/2,\,n-k-1} \cdot s\sqrt{1 + \bx_0'(\bX'\bX)^{-1}\bx_0}.
\end{equation}
\end{theorem}

\begin{proof}
\textbf{Step 1: Set up.} The prediction error is
\[
\hat{Y}_0 - Y_0 = \bx_0'(\hat{\bb} - \bb) - \varepsilon_0.
\]
Here $\hat{\bb} - \bb = (\bX'\bX)^{-1}\bX'\be$ depends on the training-sample errors $\be$, whereas $\varepsilon_0$ is the new-observation error, which by hypothesis is independent of $\be$; hence $\varepsilon_0 \indep \bx_0'(\hat{\bb} - \bb)$ (fact~\ref{N3}). Moreover the stacked vector $(\be', \varepsilon_0)'$ is jointly normal: by independence its MGF factorizes as $\expc[e^{\mathbf{t}'\be}]\,\expc[e^{t_0\varepsilon_0}] = \exp(\tfrac{\sigma^2}{2}\mathbf{t}'\mathbf{t})\exp(\tfrac{\sigma^2}{2}t_0^2)$, which is the MGF~\eqref{eq:mvn-mgf} of $N(\bzero, \sigma^2\bI_{n+1})$ (fact~\ref{N2} and Definition~\ref{def:mvn}). Consequently
\begin{equation}\label{eq:pred-stack}
\begin{pmatrix}\hat{\bb} - \bb \\ \mathbf{e} \\ \varepsilon_0\end{pmatrix} = \begin{pmatrix}(\bX'\bX)^{-1}\bX' & \bzero \\ \bM & \bzero \\ \bzero' & 1\end{pmatrix}\begin{pmatrix}\be \\ \varepsilon_0\end{pmatrix}
\end{equation}
is jointly normal (Theorem~\ref{thm:mvn-linear}), with $\cov(\mathbf{e}, \hat{\bb} - \bb) = \sigma^2\bM\bX(\bX'\bX)^{-1} = \bzero$ (as in Theorem~\ref{thm:sse-dist}, Step~3) and $\cov(\mathbf{e}, \varepsilon_0) = [\bM \;\; \bzero]\,\sigma^2\bI_{n+1}\,[\bzero' \;\; 1]' = \bzero$ (Lemma~\ref{lem:var-linear}\ref{vl:ii}).

\textbf{Step 2: Compute the variance.} Using Step~1's independence (so that the variances of the two terms add, fact~\ref{N7}) and Lemma~\ref{lem:var-linear} (with $\bA = \bx_0'$):
\begin{align*}
\var(\hat{Y}_0 - Y_0) &= \var\!\bigl(\bx_0'(\hat{\bb} - \bb)\bigr) + \var(\varepsilon_0)\\
&= \bx_0'\var(\hat{\bb})\bx_0 + \sigma^2\\
&= \sigma^2\bx_0'(\bX'\bX)^{-1}\bx_0 + \sigma^2\\
&= \sigma^2\bigl(1 + \bx_0'(\bX'\bX)^{-1}\bx_0\bigr).
\end{align*}

\textbf{Step 3: Derive the $t$ pivotal quantity.} As a linear function of the jointly normal vector $(\be', \varepsilon_0)'$ of Step~1, namely $\hat{Y}_0 - Y_0 = [\bx_0'(\bX'\bX)^{-1}\bX' \;\; {-1}]\,(\be', \varepsilon_0)'$, the prediction error is normal (Theorem~\ref{thm:mvn-linear}) with mean $0$ and the variance of Step~2:
\[
\hat{Y}_0 - Y_0 \sim N\!\bigl(0,\; \sigma^2(1 + \bx_0'(\bX'\bX)^{-1}\bx_0)\bigr).
\]
For independence from $\SSE$, note that ``$\hat{\bb} \indep \SSE$ and $\varepsilon_0 \indep \SSE$'' would \emph{not} by itself imply that the function $\bx_0'(\hat{\bb} - \bb) - \varepsilon_0$ of the pair is independent of $\SSE$; what is needed is independence of $\mathbf{e}$ from the pair $(\hat{\bb} - \bb, \varepsilon_0)$ jointly. This follows from~\eqref{eq:pred-stack}: the vector there is jointly normal and both cross-covariances of $\mathbf{e}$ with the other two blocks vanish, so Theorem~\ref{thm:mvn-uncorr-indep} (with $\mathbf{U} = \mathbf{e}$ and $\mathbf{V} = ((\hat{\bb} - \bb)', \varepsilon_0)'$) gives $\mathbf{e} \indep (\hat{\bb} - \bb, \varepsilon_0)$. Hence $\SSE = \mathbf{e}'\mathbf{e}$ is independent of $\hat{Y}_0 - Y_0$ (fact~\ref{N3}). Replacing $\sigma$ with $s$ gives a $t_{n-k-1}$ pivot (fact~\ref{N5}, exactly as in Theorem~\ref{thm:ci}):
\[
\frac{\hat{Y}_0 - Y_0}{s\sqrt{1 + \bx_0'(\bX'\bX)^{-1}\bx_0}} \sim t_{n-k-1},
\]
and inverting, using~\eqref{eq:t-symm} exactly as at the end of the proof of Theorem~\ref{thm:ci}, yields the prediction interval.
\end{proof}

\begin{remark}\label{rem:pi-wider}
The prediction interval is always wider than the confidence interval for the mean, because the prediction variance includes the additional $\sigma^2$ term from the irreducible noise $\varepsilon_0$.
\end{remark}

\begin{corollary}[Intervals in Simple Linear Regression]\label{cor:pred-slr}
In the simple linear regression model of \S\ref{sec:k1}, under~\ref{A1}--\ref{A5} and with the hypotheses of Theorem~\ref{thm:pred} for the new observation, at $x_0 \in \Real$ the $(1-\alpha) \times 100\%$ confidence interval for $\beta_0 + \beta_1x_0$ and prediction interval for $Y_0$ are
\[
\hat\beta_0 + \hat\beta_1x_0 \pm t_{\alpha/2,\,n-2}\cdot s\sqrt{\frac{1}{n} + \frac{(x_0 - \bar x)^2}{S_{xx}}}, \qquad \hat\beta_0 + \hat\beta_1x_0 \pm t_{\alpha/2,\,n-2}\cdot s\sqrt{1 + \frac{1}{n} + \frac{(x_0 - \bar x)^2}{(n-1)s_x^2}}.
\]
Both are narrowest at $x_0 = \bar x$ and widen as $x_0$ moves away from $\bar x$; their squared half-widths are quadratic in $x_0 - \bar x$, so the half-widths themselves grow approximately linearly for large $|x_0 - \bar x|$.
\end{corollary}

\begin{proof}
Apply Theorems~\ref{thm:ci-mean} and~\ref{thm:pred} with $k = 1$, $\bx_0 = (1, x_0)'$ and $n - k - 1 = n - 2$; by Corollary~\ref{cor:slr-var}, $\bx_0'(\bX'\bX)^{-1}\bx_0 = 1/n + (x_0 - \bar x)^2/S_{xx}$, and $S_{xx} = (n-1)s_x^2$ (Definition~\ref{def:corr}).
\end{proof}

%=============================================================================
\section{Constrained Least Squares and Modeling Remarks}\label{sec:special}
%=============================================================================

Theorem~\ref{thm:cls} below is a purely algebraic statement about a given $\bY$ and $\bX$; it uses only the standing assumption~\ref{A4} and no distributional assumption.

\subsection{Remarks on Model Specification}

\begin{remark}[No-intercept model]\label{rem:no-intercept}
When the model omits the intercept, $\bX$ has no column of ones. The OLS formula $\hat{\bb} = (\bX'\bX)^{-1}\bX'\bY$ still applies, but Theorem~\ref{thm:normal-eq}\ref{ne:ii} is no longer available (its proof reads $\bone'\mathbf{e} = 0$ off the intercept row of $\bX'\mathbf{e} = \bzero$), so $\sum e_i \neq 0$ in general; since Steps~1, 2 and~4 of the proof of Theorem~\ref{thm:anova} all use $\bone'\mathbf{e} = 0$, the decomposition $\SST = \SSR + \SSE$ may fail. Consequently the two expressions in Definition~\ref{def:R2} no longer agree and Proposition~\ref{prop:R2-bounds} does not apply: $1 - \SSE/\SST$ can be negative, while $\SSR/\SST \geqslant 0$ can exceed $1$. (For $n = 2$ with the single regressor $\bX = (1, 0)'$ and $\bY = (1, 3)'$: $\hat{\bb} = 1$, $\mathbf{e} = (0, 3)'$, $\SST = 2$, $\SSR = 5$, $\SSE = 9$, so $1 - \SSE/\SST = -3.5$ and $\SSR/\SST = 2.5$.) Use the uncentered $R^2 = 1 - \SSE/\sum Y_i^2 = 1 - \SSE/\bY'\bY$ instead (for $\bY \neq \bzero$): it lies in $[0, 1]$ because $\SSE = \|\bM\bY\|^2 = \bY'\bM\bY$ and $\bY'\bY - \bY'\bM\bY = \bY'\bH\bY = \|\bH\bY\|^2 \geqslant 0$ (Theorem~\ref{thm:HM}\ref{HM:i}); in the example it equals $1 - 9/10 = 0.1$.
\end{remark}

\begin{remark}[Dummy variable trap]\label{rem:dummy-trap}
A categorical variable with $g$ levels requires $g-1$ dummy regressors when the model includes an intercept. Including all $g$ dummies creates the linear dependence $\bone = D_1 + \cdots + D_g$, violating~\ref{A4}.
\end{remark}

\subsection{Constrained Least Squares}

\begin{theorem}[Constrained OLS]\label{thm:cls}
Consider $\min_{\bb} \|\bY - \bX\bb\|^2$ subject to $\bR\bb = \br$ ($q$ linear constraints, $\rank(\bR) = q$). The constrained estimator is:
\begin{equation}\label{eq:cls}
\hat{\bb}_R = \hat{\bb} - (\bX'\bX)^{-1}\bR'\bigl[\bR(\bX'\bX)^{-1}\bR'\bigr]^{-1}(\bR\hat{\bb} - \br),
\end{equation}
and $\SSE_R \geqslant \SSE_U$ (where $\SSE_U$ is the unconstrained residual sum of squares).
\end{theorem}

\begin{proof}
The objective $\|\bY - \bX\bb\|^2$ is a convex quadratic function of $\bb$ (its Hessian $2\bX'\bX$ is positive definite by~\ref{A4}), and the constraint $\bR\bb = \br$ is affine. Rather than invoking the theory of constrained optimization (the KKT conditions), we use the Lagrangian only to \emph{produce a candidate} (Steps~1--2) and then verify directly that the candidate is the unique minimizer over the feasible set (Step~3); nothing from outside these notes is needed.

\textbf{Step 1: Lagrangian and first-order conditions.} The Lagrangian is $\mathcal{L}(\bb, \boldsymbol\lambda) = (\bY - \bX\bb)'(\bY - \bX\bb) + 2\boldsymbol\lambda'(\bR\bb - \br)$ (the factor $2$ simplifies later algebra). By Lemma~\ref{lem:matrix-diff}, differentiating $(\bY - \bX\bb)'(\bY - \bX\bb) = \bY'\bY - 2\bY'\bX\bb + \bb'\bX'\bX\bb$ with respect to $\bb'$ gives $-2\bY'\bX + 2\bb'\bX'\bX$. The first-order conditions are
\begin{align*}
\frac{\partial \mathcal{L}}{\partial \bb'} &= -2\bY'\bX + 2\bb'\bX'\bX + 2\boldsymbol\lambda'\bR = \bzero', \tag{i}\\
\frac{\partial \mathcal{L}}{\partial \boldsymbol\lambda'} &= 2(\bR\bb - \br)' = \bzero'. \tag{ii}
\end{align*}
Transposing (i) gives $\bX'\bX\hat{\bb}_R + \bR'\hat{\boldsymbol\lambda} = \bX'\bY$, so
\begin{equation}\label{eq:cls-step1}
\hat{\bb}_R = (\bX'\bX)^{-1}\bX'\bY - (\bX'\bX)^{-1}\bR'\hat{\boldsymbol\lambda} = \hat{\bb} - (\bX'\bX)^{-1}\bR'\hat{\boldsymbol\lambda}.
\end{equation}

\textbf{Step 2: Solve for $\hat{\boldsymbol\lambda}$.} Left-multiply~\eqref{eq:cls-step1} by $\bR$ and impose the constraint (ii), $\bR\hat{\bb}_R = \br$:
\[
\br = \bR\hat{\bb}_R = \bR\hat{\bb} - \bR(\bX'\bX)^{-1}\bR'\hat{\boldsymbol\lambda}.
\]
Since $\rank(\bR) = q$ and $(\bX'\bX)^{-1}$ is positive definite (Theorem~\ref{thm:ols}, Step~4), $\bR(\bX'\bX)^{-1}\bR'$ is positive definite (hence invertible); see Step~2 of the proof of Theorem~\ref{thm:general-F} for the detailed argument. Rearranging:
\[
\hat{\boldsymbol\lambda} = [\bR(\bX'\bX)^{-1}\bR']^{-1}(\bR\hat{\bb} - \br).
\]
Substituting back into~\eqref{eq:cls-step1} gives~\eqref{eq:cls}.

\textbf{Step 3: $\hat{\bb}_R$ is the unique minimizer over the feasible set.} Define $\hat{\bb}_R$ by~\eqref{eq:cls} and $\hat{\boldsymbol\lambda}$ as in Step~2. By construction, $\bR\hat{\bb}_R = \bR\hat{\bb} - \bR(\bX'\bX)^{-1}\bR'\hat{\boldsymbol\lambda} = \bR\hat{\bb} - (\bR\hat{\bb} - \br) = \br$ (feasibility), and, transposing (i),
\begin{equation}\label{eq:cls-foc}
\bX'(\bY - \bX\hat{\bb}_R) = \bX'\bY - \bX'\bX\hat{\bb} + \bR'\hat{\boldsymbol\lambda} = \bR'\hat{\boldsymbol\lambda},
\end{equation}
using $\bX'\bX\hat{\bb} = \bX'\bY$. Now let $\bb$ be any feasible vector, $\bR\bb = \br$. Writing $\bY - \bX\bb = (\bY - \bX\hat{\bb}_R) + \bX(\hat{\bb}_R - \bb)$, the cross term vanishes:
\[
(\bY - \bX\hat{\bb}_R)'\bX(\hat{\bb}_R - \bb) = \hat{\boldsymbol\lambda}'\bR(\hat{\bb}_R - \bb) = \hat{\boldsymbol\lambda}'(\br - \br) = 0
\]
by~\eqref{eq:cls-foc}. Hence
\[
\|\bY - \bX\bb\|^2 = \|\bY - \bX\hat{\bb}_R\|^2 + \|\bX(\hat{\bb}_R - \bb)\|^2 \;\geqslant\; \|\bY - \bX\hat{\bb}_R\|^2,
\]
with equality iff $\bX(\hat{\bb}_R - \bb) = \bzero$, i.e.\ (by~\ref{A4}) iff $\bb = \hat{\bb}_R$. So $\hat{\bb}_R$ is the unique constrained minimizer and $\SSE_R = \|\bY - \bX\hat{\bb}_R\|^2$.

\textbf{Step 4: $\SSE_R \geqslant \SSE_U$.} The constrained problem minimizes $\|\bY - \bX\bb\|^2$ over $\{\bb : \bR\bb = \br\} \subseteq \Real^{k+1}$. Since $\hat{\bb}$ minimizes the same objective over the full space, the minimum over a subset cannot be smaller: $\SSE_R \geqslant \SSE_U$. Equality holds iff $\hat{\bb}$ already satisfies the constraint, i.e., $\bR\hat{\bb} = \br$: if $\bR\hat{\bb} = \br$, then~\eqref{eq:cls} gives $\hat{\bb}_R = \hat{\bb}$ and hence $\SSE_R = \SSE_U$; conversely, if $\SSE_R = \SSE_U$, then $Q(\hat{\bb}_R) = \SSE_R = \SSE_U = Q(\hat{\bb})$, so $\hat{\bb}_R$ minimizes $Q$ over all of $\Real^{k+1}$, hence $\hat{\bb}_R = \hat{\bb}$ by the uniqueness in Theorem~\ref{thm:ols} (by~\eqref{eq:ols-global} with $\bb = \hat{\bb}_R$, $\bX(\hat{\bb}_R - \hat{\bb}) = \bzero$, so $\hat{\bb}_R = \hat{\bb}$ under~\ref{A4}), and therefore $\bR\hat{\bb} = \bR\hat{\bb}_R = \br$ by the feasibility established in Step~3.
\end{proof}

%=============================================================================
\section{Qualitative Regressors, Reparametrization, and Interactions}\label{sec:qual}
%=============================================================================

This section develops the modeling remarks of Section~\ref{sec:special} (Remark~\ref{rem:dummy-trap}). Its results are algebraic identities about a given $\bY$ and design matrix and use only the standing assumption~\ref{A4}, except where a distribution is asserted; such statements use~\ref{A1}--\ref{A5}, as in Section~\ref{sec:test}.

\subsection{Qualitative Regressors}

\begin{definition}[Factors and Indicator Coding]\label{def:factor}
A \textbf{qualitative regressor} (or \textbf{factor}) assigns each observation to exactly one of $g \geqslant 2$ \textbf{levels}. The \textbf{indicator} (or dummy variable) of level $l$ is the vector $D_l \in \Real^n$ with $i$-th entry $1$ if observation $i$ is at level $l$ and $0$ otherwise. \textbf{Baseline coding} enters the intercept and the indicators of all levels but one, the \textbf{baseline}; by Remark~\ref{rem:dummy-trap}, entering all $g$ indicators together with the intercept would violate~\ref{A4}. A \textbf{binary} regressor is a factor with $g = 2$, entered as a single indicator.
\end{definition}

\begin{proposition}[Coefficients of a Factor Are Differences of Level Means]\label{prop:level-means}
Let $\bY$ be regressed on $\bX = [\bone \;\; D_2 \;\; \cdots \;\; D_g]$ (baseline level $1$; here the coefficient of $D_l$ is indexed by $l$, so the coefficients are $\beta_0, \beta_2, \ldots, \beta_g$), and suppose every level is observed: level $l$ has $n_l \geqslant 1$ observations, with mean response $\bar Y_l = D_l'\bY/n_l$. Then $\bX$ has full column rank, and:
\begin{enumerate}[label=\textup{(\roman*)}]
\item the fitted value of each observation is the mean response of its level;
\item $\hat\beta_0 = \bar Y_1$ and $\hat\beta_l = \bar Y_l - \bar Y_1$ for $l = 2, \ldots, g$;
\item $\SSE = \sum_{l=1}^g\sum_{i:\,D_{l,i} = 1}(Y_i - \bar Y_l)^2$ is the within-level sum of squares, and for $n > g$, $s^2 = \SSE/(n-g)$.
\end{enumerate}
\end{proposition}

\begin{proof}
\textbf{Step 1: Column space and rank.} Every observation is at exactly one level, so $D_1 + \cdots + D_g = \bone$, i.e.\ $D_1 = \bone - D_2 - \cdots - D_g$. Hence $\calC(\bX) = \calC([D_1 \;\; \cdots \;\; D_g])$. The indicators have disjoint, non-empty supports, so $D_l'D_m = 0$ for $l \neq m$ and $D_l'D_l = n_l > 0$; if $\sum_l a_lD_l = \bzero$, multiplying by $D_m'$ gives $a_mn_m = 0$, so $D_1, \ldots, D_g$ are linearly independent. Thus $\dim\calC(\bX) = g$, the number of columns of $\bX$, and $\bX$ has full column rank.

\textbf{Step 2: The fitted values.} Let $\bv = \sum_l\bar Y_lD_l \in \calC(\bX)$. For each $m$, $D_m'(\bY - \bv) = D_m'\bY - n_m\bar Y_m = 0$, so $\bY - \bv$ is orthogonal to every $D_m$, hence to $\calC(\bX)$. By Step~1 of the proof of Theorem~\ref{thm:ols-proj}, $\bv$ is then the unique point of $\calC(\bX)$ closest to $\bY$, i.e.\ $\hat{\bY} = \bv$. This is (i), and $\SSE = \|\bY - \bv\|^2$ is the within-level sum of squares in (iii); with $k = g - 1$ regressors, $n - k - 1 = n - g$.

\textbf{Step 3: The coefficients.} At level $1$ all of $D_2, \ldots, D_g$ vanish, so the fitted value is $\hat\beta_0$, and at level $l \geqslant 2$ it is $\hat\beta_0 + \hat\beta_l$. As every level is observed, (i) gives $\hat\beta_0 = \bar Y_1$ and $\hat\beta_0 + \hat\beta_l = \bar Y_l$.
\end{proof}

\begin{corollary}[One Indicator: the Pooled Two-Sample $t$ Test]\label{cor:two-sample}
Let $g = 2$, $\bX = [\bone \;\; D_2]$, $n_1, n_2 \geqslant 1$ and $n = n_1 + n_2 \geqslant 3$. Then $\hat\beta_0 = \bar Y_1$, $\hat\beta_1 = \bar Y_2 - \bar Y_1$,
\[
s^2 = \frac{(n_1 - 1)s_1^2 + (n_2 - 1)s_2^2}{n - 2}, \qquad \SE(\hat\beta_1) = s\sqrt{\frac{1}{n_1} + \frac{1}{n_2}},
\]
where $(n_l - 1)s_l^2$ is the sum of squared deviations of the responses at level $l$ from $\bar Y_l$ (so $s^2$ is the pooled variance), and, when $s > 0$ (under~\ref{A5} this holds with probability one, as in Proposition~\ref{prop:duality}), the statistic of Theorem~\ref{thm:t-test} for $H_0\colon \beta_1 = 0$ is the pooled two-sample $t$ statistic
\[
t = \frac{\bar Y_2 - \bar Y_1}{s\sqrt{1/n_1 + 1/n_2}}.
\]
Under~\ref{A1}--\ref{A5} the responses are normal with common variance $\sigma^2$ and means $\mu_1 = \beta_0$ at level $1$ and $\mu_2 = \beta_0 + \beta_1$ at level $2$, and under $H_0\colon \mu_1 = \mu_2$, $t \sim t_{n-2}$; moreover the overall $F$ statistic equals $t^2$ (Corollary~\ref{cor:drop-one}).
\end{corollary}

\begin{proof}
The coefficients and $s^2$ are Proposition~\ref{prop:level-means} with $g = 2$ (the coefficient of $D_2$, called $\beta_2$ there, is indexed here by its column, as $\beta_1$). For the standard error, apply Corollary~\ref{cor:slr} with $\bx = D_2$: $\bar x = n_2/n$ and $x_i^2 = x_i$, so $S_{xx} = \sum_i x_i^2 - n\bar x^2 = n_2 - n_2^2/n = n_1n_2/n$, and $[(\bX'\bX)^{-1}]_{11} = 1/S_{xx} = n/(n_1n_2) = 1/n_1 + 1/n_2$ (the entry belonging to $\beta_1$). Hence $\SE(\hat\beta_1) = s\sqrt{1/n_1 + 1/n_2}$ (Theorem~\ref{thm:t-test}). Under~\ref{A1}, $\expc[Y_i] = \beta_0 + \beta_1D_{2,i}$, which gives the two means, and $H_0\colon \beta_1 = 0$ is $\mu_1 = \mu_2$; Theorem~\ref{thm:t-test} with $n - k - 1 = n - 2$ gives the distribution, and Corollary~\ref{cor:drop-one} with $k = 1$ gives $F = t^2$.
\end{proof}

\begin{proposition}[Reparametrization]\label{prop:reparam}
Let $\bX$ satisfy~\ref{A4} and let $\tilde{\bX} = \bX\bA$ with $\bA$ an invertible $(k+1) \times (k+1)$ matrix. Then:
\begin{enumerate}[label=\textup{(\roman*)}]
\item $\tilde{\bX}$ has full column rank and $\calC(\tilde{\bX}) = \calC(\bX)$;
\item\label{rp:ii} the regressions of $\bY$ on $\bX$ and on $\tilde{\bX}$ have the same fitted values, residuals, $\SSE$, $\SSR$ and $s^2$, hence the same $R^2$ and overall $F$ statistic;
\item\label{rp:iii} $\hat{\tilde{\bb}} = \bA^{-1}\hat{\bb}$, and under~\ref{A3}, $\var(\hat{\tilde{\bb}}) = \sigma^2\bA^{-1}(\bX'\bX)^{-1}(\bA^{-1})'$;
\item\label{rp:iv} conversely, if $\tilde{\bX}$ is any $n \times (k+1)$ matrix with $\calC(\tilde{\bX}) = \calC(\bX)$, then $\tilde{\bX} = \bX\bA$ for an invertible $\bA$.
\end{enumerate}
\end{proposition}

\begin{proof}
\textbf{(i)} $\tilde{\bX}\bv = \bX(\bA\bv)$ shows $\calC(\tilde{\bX}) \subseteq \calC(\bX)$, and $\bX\bv = \tilde{\bX}(\bA^{-1}\bv)$ the reverse inclusion. If $\tilde{\bX}\bv = \bzero$, then $\bX(\bA\bv) = \bzero$, so $\bA\bv = \bzero$ by~\ref{A4} and $\bv = \bzero$.

\textbf{(ii)} By Theorem~\ref{thm:ols-proj}, each fitted vector is the unique point of the common column space closest to $\bY$, so the fitted values coincide, and with them the residuals, $\SSE$ and $\SSR = \sum_i(\hat Y_i - \bar Y)^2$; both designs have $k + 1$ columns, so $s^2$ agrees as well, and $R^2$ and $F$ are functions of these quantities.

\textbf{(iii)} Using $(\bA'\bX'\bX\bA)^{-1} = \bA^{-1}(\bX'\bX)^{-1}(\bA')^{-1}$,
\[
\hat{\tilde{\bb}} = (\tilde{\bX}'\tilde{\bX})^{-1}\tilde{\bX}'\bY = \bA^{-1}(\bX'\bX)^{-1}(\bA')^{-1}\bA'\bX'\bY = \bA^{-1}\hat{\bb},
\]
and Theorem~\ref{thm:var-beta} for the design $\tilde{\bX}$ gives $\var(\hat{\tilde{\bb}}) = \sigma^2(\tilde{\bX}'\tilde{\bX})^{-1} = \sigma^2\bA^{-1}(\bX'\bX)^{-1}(\bA^{-1})'$.

\textbf{(iv)} Every column of $\tilde{\bX}$ lies in $\calC(\bX)$, so $\tilde{\bX} = \bX\bA$ for some $(k+1) \times (k+1)$ matrix $\bA$; likewise $\bX = \tilde{\bX}\mathbf{B}$ for some $\mathbf{B}$. Then $\bX(\bA\mathbf{B} - \bI) = \bzero$, and~\ref{A4}, applied to each column of $\bA\mathbf{B} - \bI$, gives $\bA\mathbf{B} = \bI$; so the square matrix $\bA$ is invertible.
\end{proof}

\begin{corollary}[Changes of Coding]\label{cor:coding}
Each of the following changes of the design replaces $\bX$ by $\bX\bA$ with $\bA$ invertible, other columns being left unchanged, and therefore does not change the fit (Proposition~\ref{prop:reparam}):
\begin{enumerate}[label=\textup{(\alph*)}]
\item\label{cod:a} changing the baseline level of a factor entered by baseline coding;
\item\label{cod:b} coding a binary regressor as $\pm 1$, i.e.\ replacing its indicator $D$ by $2D - \bone$; with no other regressors the new intercept is $\frac{1}{2}(\bar Y_1 + \bar Y_2)$ and the new slope is $\frac{1}{2}(\bar Y_2 - \bar Y_1)$;
\item centering a regressor, $\bX_{\cdot j} \mapsto \bX_{\cdot j} - \bar X_j\bone$ ($j \geqslant 1$);
\item rescaling a regressor, $\bX_{\cdot j} \mapsto a\bX_{\cdot j}$ with $a \neq 0$, which replaces $\hat\beta_j$ by $\hat\beta_j/a$.
\end{enumerate}
In~\ref{cod:a} and~\ref{cod:b}, the $F$ statistic of Corollary~\ref{cor:nested-F} for the hypothesis that all coefficients of the factor's indicators are zero has the same value under both codings. In~\ref{cod:a} the intercept, the coefficients of the indicators and their $t$ statistics in general change; in~\ref{cod:b} the indicator's coefficient and its standard error are both halved, so its $t$ statistic is unchanged (its square is that $F$ statistic, Corollary~\ref{cor:drop-one}), and only the intercept and its $t$ statistic change in general. In both, the coefficients of the other columns, their standard errors and their $t$ statistics are unchanged.
\end{corollary}

\begin{proof}
\textbf{(a)} Let $\mathbf{Z}$ denote the other columns. The design $[\bone \;\; D_2 \;\; \cdots \;\; D_g \;\; \mathbf{Z}]$ and the design $[\bone \;\; D_1 \;\; D_3 \;\; \cdots \;\; D_g \;\; \mathbf{Z}]$ both have column space $\calC([D_1 \;\; \cdots \;\; D_g \;\; \mathbf{Z}])$, because $\bone = \sum_lD_l$, $D_1 = \bone - \sum_{l \geqslant 2}D_l$ and $D_2 = \bone - \sum_{l \neq 2}D_l$. They have the same number of columns, so Proposition~\ref{prop:reparam}\ref{rp:iv} applies. \textbf{(b)} $[\bone \;\; 2D - \bone] = [\bone \;\; D]\begin{pmatrix} 1 & -1\\ 0 & 2\end{pmatrix}$, an invertible matrix, extended by the identity on the other columns. Without other regressors, $\hat{\bb} = (\bar Y_1, \bar Y_2 - \bar Y_1)'$ by Proposition~\ref{prop:level-means}, and $\bA^{-1} = \begin{pmatrix} 1 & 1/2\\ 0 & 1/2\end{pmatrix}$ gives $\bA^{-1}\hat{\bb} = (\frac{1}{2}(\bar Y_1 + \bar Y_2), \frac{1}{2}(\bar Y_2 - \bar Y_1))'$. \textbf{(c)} The new column is $\bX(\boldsymbol\iota_j - \bar X_j\boldsymbol\iota_0)$, where $\boldsymbol\iota_0$ and $\boldsymbol\iota_j$ select the intercept column and column $j$; so $\bA = \bI - \bar X_j\boldsymbol\iota_0\boldsymbol\iota_j'$, whose inverse is $\bI + \bar X_j\boldsymbol\iota_0\boldsymbol\iota_j'$ because $\boldsymbol\iota_j'\boldsymbol\iota_0 = 0$. \textbf{(d)} $\bA$ is diagonal with $a$ in position $j$ and ones elsewhere, and the $j$-th entry of $\bA^{-1}\hat{\bb}$ is $\hat\beta_j/a$.

For the $F$ statistic: in both~(a) and~(b), setting the coefficients of the factor's indicators to zero leaves the regression on $\bone$ and $\mathbf{Z}$ (Corollary~\ref{cor:nested-R2}), the same under both codings, so $\SSE_R$ agrees; $\SSE_U$ agrees by Proposition~\ref{prop:reparam}\ref{rp:ii}; and $q$ and $n - k - 1$ are the same. The coefficients $\bA^{-1}\hat{\bb}$, however, depend on $\bA$. In~(b), the row of $\bA^{-1}$ belonging to the indicator is $\frac{1}{2}$ times the corresponding unit row, so by Proposition~\ref{prop:reparam}\ref{rp:iii} its new coefficient is $\frac{1}{2}$ times the old one, and since $(\tilde{\bX}'\tilde{\bX})^{-1} = (\bA'\bX'\bX\bA)^{-1} = \bA^{-1}(\bX'\bX)^{-1}(\bA^{-1})'$, the corresponding diagonal entry of $(\tilde{\bX}'\tilde{\bX})^{-1}$ is $\frac{1}{4}$ times the old one; $s$ is unchanged by Proposition~\ref{prop:reparam}\ref{rp:ii}, so the standard error is halved and the $t$ statistic is unchanged. The intercept row of $\bA^{-1}$ is $(1, \frac{1}{2}, 0, \ldots, 0)$, so the new intercept is $\hat\beta_0 + \frac{1}{2}\hat\beta_D$, where $\hat\beta_D$ is the old indicator coefficient. Finally, in both~(a) and~(b) the new columns of the factor are combinations of $\bone$ and the indicators only, so $\bA$ is block diagonal with the identity on the columns of $\mathbf{Z}$, and so is $\bA^{-1}$; the rows of $\bA^{-1}$ belonging to $\mathbf{Z}$ are therefore unit rows, and the same computation shows that the coefficients of $\mathbf{Z}$ and the corresponding diagonal entries of $(\tilde{\bX}'\tilde{\bX})^{-1}$, hence their standard errors and $t$ statistics, are unchanged.
\end{proof}

\subsection{Regressors Defined by Known Functions}

\begin{definition}[Linear in the Parameters]\label{def:linear-params}
The theory of these notes requires linearity in $\bb$, not in the observed variables. A model $\expc[Y_i] = \beta_0 + \beta_1g_1(\mathbf{z}_i) + \cdots + \beta_kg_k(\mathbf{z}_i)$, with known functions $g_1, \ldots, g_k$ of observed variables $\mathbf{z}_i$, is the model of Definition~\ref{def:model} with $X_{ij} = g_j(\mathbf{z}_i)$. Examples are \emph{polynomial regression}, $g_j(z) = z^j$; logarithms, $g(z) = \ln z$; a threshold indicator $g(z) = 1\{z \geqslant \xi\}$; and the hinge $g(z) = 1\{z \geqslant \xi\}(z - \xi)$, with a known knot $\xi$. A model such as $\expc[Y_i] = \beta_0 + \exp(\beta_1z_i)$ is not linear in the parameters.
\end{definition}

\begin{proposition}[Piecewise Linear Regression Functions]\label{prop:piecewise}
Let $\xi$ be known and $d(z) = 1\{z \geqslant \xi\}$.
\begin{enumerate}[label=\textup{(\roman*)}]
\item $f(z) = \beta_0 + \beta_1z + \beta_2d(z)(z - \xi)$ is continuous, with slope $\beta_1$ for $z < \xi$ and $\beta_1 + \beta_2$ for $z > \xi$.
\item $f(z) = \beta_0 + \beta_1z + \beta_2d(z) + \beta_3d(z)z$ has slope $\beta_1$ for $z < \xi$ and $\beta_1 + \beta_3$ for $z > \xi$, and a jump at $\xi$ of size $f(\xi) - \lim_{z \uparrow \xi}f(z) = \beta_2 + \beta_3\xi$.
\end{enumerate}
\end{proposition}

\begin{proof}
For $z < \xi$, $d(z) = 0$ and both functions equal $\beta_0 + \beta_1z$, whose limit at $\xi$ is $\beta_0 + \beta_1\xi$. For $z \geqslant \xi$, $d(z) = 1$: the function in (i) is $(\beta_0 - \beta_2\xi) + (\beta_1 + \beta_2)z$, whose value at $\xi$ is $\beta_0 + \beta_1\xi$, so $f$ is continuous at $\xi$ (and, being linear on each side, everywhere); the function in (ii) is $(\beta_0 + \beta_2) + (\beta_1 + \beta_3)z$, whose value at $\xi$ exceeds $\beta_0 + \beta_1\xi$ by $\beta_2 + \beta_3\xi$.
\end{proof}

\subsection{Interactions}

\begin{definition}[Interaction]\label{def:interaction}
For two regressors $\bX_{\cdot 1}$ and $\bX_{\cdot 2}$, the \textbf{interaction} is their elementwise product $\bX_{\cdot 1} \circ \bX_{\cdot 2}$, with entries $X_{i1}X_{i2}$, entered as a further regressor:
\[
\expc[Y_i] = \beta_0 + \beta_1X_{i1} + \beta_2X_{i2} + \beta_3X_{i1}X_{i2}.
\]
$\bX_{\cdot 1}$ and $\bX_{\cdot 2}$ are the \textbf{main effects}. A model without interaction terms is called \textbf{additive}.
\end{definition}

\begin{proposition}[Effects in an Interaction Model]\label{prop:interaction}
Let $f(x_1, x_2) = \beta_0 + \beta_1x_1 + \beta_2x_2 + \beta_3x_1x_2$. Then $\partial f/\partial x_1 = \beta_1 + \beta_3x_2$, $\partial f/\partial x_2 = \beta_2 + \beta_3x_1$, and $f(x_1 + 1, x_2) - f(x_1, x_2) = \beta_1 + \beta_3x_2$: the effect of $x_1$ depends on the level of $x_2$. If $x_2 = z \in \{0, 1\}$ is an indicator, then $f(x_1, 0) = \beta_0 + \beta_1x_1$ and $f(x_1, 1) = (\beta_0 + \beta_2) + (\beta_1 + \beta_3)x_1$: two lines whose intercepts differ by $\beta_2$ and whose slopes differ by $\beta_3$; they are parallel iff $\beta_3 = 0$.
\end{proposition}

\begin{proof}
Differentiate term by term, using $\partial(\beta_3x_1x_2)/\partial x_1 = \beta_3x_2$ and symmetrically in $x_2$; the difference is $\beta_1 + \beta_3x_2$ by direct expansion; and substituting $z = 0$ and $z = 1$ gives the two lines.
\end{proof}

\begin{proposition}[The Hierarchical Principle]\label{prop:hierarchy}
Let $\bX = [\bone \;\; \bX_{\cdot 1} \;\; \bX_{\cdot 2} \;\; \bX_{\cdot 1} \circ \bX_{\cdot 2} \;\; \mathbf{Z}]$ satisfy~\ref{A4}, where $\mathbf{Z}$ (possibly empty) collects further columns, and for constants $a, b$ let $\tilde{\bX}$ be obtained by replacing $\bX_{\cdot 1}$, $\bX_{\cdot 2}$ and $\bX_{\cdot 1} \circ \bX_{\cdot 2}$ by $\bX_{\cdot 1} + a\bone$, $\bX_{\cdot 2} + b\bone$ and $(\bX_{\cdot 1} + a\bone) \circ (\bX_{\cdot 2} + b\bone)$. Then:
\begin{enumerate}[label=\textup{(\roman*)}]
\item $\tilde{\bX} = \bX\bA$ with $\bA$ invertible, so the two fits coincide (Proposition~\ref{prop:reparam});
\item the coefficients become $\tilde\beta_3 = \hat\beta_3$, $\tilde\beta_1 = \hat\beta_1 - b\hat\beta_3$, $\tilde\beta_2 = \hat\beta_2 - a\hat\beta_3$, $\tilde\beta_0 = \hat\beta_0 - a\hat\beta_1 - b\hat\beta_2 + ab\hat\beta_3$, the coefficients of $\mathbf{Z}$ being unchanged;
\item the standard error and the $t$ statistic (Theorem~\ref{thm:t-test}) of the interaction coefficient are unchanged;
\item\label{hier:iv} if the main effects are omitted, i.e.\ for the design $[\bone \;\; \bX_{\cdot 1} \circ \bX_{\cdot 2} \;\; \mathbf{Z}]$ and its shifted version, the column space changes unless $b\bX_{\cdot 1} + a\bX_{\cdot 2}$ lies in the column space of the unshifted design.
\end{enumerate}
\end{proposition}

\begin{proof}
\textbf{(i)} Since $(\bX_{\cdot 1} + a\bone) \circ (\bX_{\cdot 2} + b\bone) = \bX_{\cdot 1} \circ \bX_{\cdot 2} + b\bX_{\cdot 1} + a\bX_{\cdot 2} + ab\bone$, $\tilde{\bX} = \bX\bA$ with $\bA$ block diagonal, with the identity on the columns of $\mathbf{Z}$ and, on the first four columns,
\[
\bA_4 = \begin{pmatrix} 1 & a & b & ab\\ 0 & 1 & 0 & b\\ 0 & 0 & 1 & a\\ 0 & 0 & 0 & 1\end{pmatrix}.
\]
$\bA_4$ is invertible: $\bA_4\bv = \bzero$ gives $v_4 = 0$ from the last row, then $v_3 = v_2 = 0$ from the third and second rows, and $v_1 = 0$ from the first. Hence $\bA$ is invertible.

\textbf{(ii)} Writing the fitted function in the shifted variables, $\tilde\beta_0 + \tilde\beta_1(x_1 + a) + \tilde\beta_2(x_2 + b) + \tilde\beta_3(x_1 + a)(x_2 + b)$ expands to
\[
(\tilde\beta_0 + a\tilde\beta_1 + b\tilde\beta_2 + ab\tilde\beta_3) + (\tilde\beta_1 + b\tilde\beta_3)x_1 + (\tilde\beta_2 + a\tilde\beta_3)x_2 + \tilde\beta_3x_1x_2.
\]
By (i) the fitted values coincide, i.e.\ $\tilde{\bX}\hat{\tilde{\bb}} = \bX\bA\hat{\tilde{\bb}} = \bX\hat{\bb}$, so $\bA\hat{\tilde{\bb}} = \hat{\bb}$ by~\ref{A4}; reading off the coordinates of $\bA\hat{\tilde{\bb}}$, which are the four brackets above and the coefficients of $\mathbf{Z}$, and solving from the interaction upward gives the stated values.

\textbf{(iii)} $\bA_4^{-1}$ is again upper triangular with unit diagonal (its last row is $(0, 0, 0, 1)$ by the last row of $\bA_4\bA_4^{-1} = \bI$), so for the coordinate vector $\mathbf{u}$ of the interaction, $\mathbf{u}'\bA^{-1} = \mathbf{u}'$. By Proposition~\ref{prop:reparam}\ref{rp:iii}, $(\tilde{\bX}'\tilde{\bX})^{-1} = \bA^{-1}(\bX'\bX)^{-1}(\bA^{-1})'$, whence $\mathbf{u}'(\tilde{\bX}'\tilde{\bX})^{-1}\mathbf{u} = \mathbf{u}'(\bX'\bX)^{-1}\mathbf{u}$. Since $s$ is the same for both fits (Proposition~\ref{prop:reparam}\ref{rp:ii}) and $\tilde\beta_3 = \hat\beta_3$, the standard error and $t$ statistic agree.

\textbf{(iv)} The shifted interaction column equals $\bX_{\cdot 1} \circ \bX_{\cdot 2} + (b\bX_{\cdot 1} + a\bX_{\cdot 2}) + ab\bone$. If the two column spaces coincide, the shifted column lies in the unshifted column space, which contains $\bone$ and $\bX_{\cdot 1} \circ \bX_{\cdot 2}$; hence so does $b\bX_{\cdot 1} + a\bX_{\cdot 2}$.
\end{proof}

\begin{remark}[The meaning of the hierarchical principle]\label{rem:hierarchy}
With both main effects present, the fit and the inference on the interaction do not depend on the origins chosen for $x_1$ and $x_2$, while the main-effect coefficients are the slopes at $x_2 = 0$ and at $x_1 = 0$ (Proposition~\ref{prop:interaction}) and change with the origin; their individual $t$ tests therefore answer origin-dependent questions. Dropping a main effect because its $t$ statistic is small makes, by Proposition~\ref{prop:hierarchy}\ref{hier:iv}, the column space, and in general the fit, depend on an arbitrary origin. (Part~\ref{hier:iv} treats the omission of both main effects; if only $\bX_{\cdot 2}$ is omitted, the same argument, with $\bone$, $\bX_{\cdot 1}$ and $\bX_{\cdot 1} \circ \bX_{\cdot 2}$ in the unshifted design, shows that the column space changes unless $a\bX_{\cdot 2}$ lies in it, and symmetrically with $b\bX_{\cdot 1}$ if only $\bX_{\cdot 1}$ is omitted.) This is the \textbf{hierarchical principle}: if an interaction is included, its main effects are included as well, whatever their $p$-values.
\end{remark}

%=============================================================================
\section{Exercises}
%=============================================================================

\exercisetier{A}
\subsection*{Tier A: Reproduce the Derivation --- Estimation, Distribution Theory and the Overall $F$ Test (Sections~\ref{sec:ols}--\ref{sec:test})}

\begin{exercise}\label{ex:A1}
Establish the matrix differentiation formulas of Lemma~\ref{lem:matrix-diff}: for constant $\ba\in\Real^p$ and constant $\bA\in\Real^{p\times p}$, prove $D(\ba'\bb)=\ba'$ and $D(\bb'\bA\bb)=\bb'(\bA+\bA')$ by computing the partial derivatives $\partial/\partial\beta_j$ coordinate by coordinate. Then apply the formulas to $Q(\bb) = \bY'\bY - 2\bb'\bX'\bY + \bb'\bX'\bX\bb$ to obtain $DQ(\bb) = -2\bY'\bX + 2\bb'\bX'\bX$.
\end{exercise}

\begin{exercise}\label{ex:A2}
Let $\bA\in\Real^{n\times n}$ be symmetric and idempotent ($\bA^2=\bA=\bA'$). (a)~Show every eigenvalue of $\bA$ lies in $\{0,1\}$. (b)~Conclude $\tr(\bA)=\rank(\bA)$. (c)~Verify the formula on $\bH$ and $\bM$: $\tr(\bH)=k+1$ and $\tr(\bM)=n-k-1$.
\end{exercise}

\begin{exercise}\label{ex:A3}
Derive the OLS estimator $\hat{\bb} = (\bX'\bX)^{-1}\bX'\bY$ by minimizing $\|\bY - \bX\bb\|^2$. Verify the second-order condition.
\end{exercise}

\begin{exercise}\label{ex:A4}
For $\bH = \bX(\bX'\bX)^{-1}\bX'$ and $\bM = \bI - \bH$: (a)~prove $\bH$ and $\bM$ are symmetric and idempotent; (b)~prove $\bH\bX = \bX$, $\bM\bX = \bzero$, and $\bH\bM = \bzero$; (c)~conclude $\tr(\bH) = k+1$ and $\tr(\bM) = n-k-1$.
\end{exercise}

\begin{exercise}\label{ex:A5}
Prove that $\expc[\hat{\bb}] = \bb$ under (A2) and (A4), and derive $\var(\hat{\bb}) = \sigma^2(\bX'\bX)^{-1}$ under (A3) and (A4).
\end{exercise}

\begin{exercise}\label{ex:A6}
State and prove the Gauss--Markov theorem.
\end{exercise}

\begin{exercise}\label{ex:A7}
Prove that $\MSE = \SSE/(n-k-1)$ is an unbiased estimator of $\sigma^2$ under (A2)--(A4).
\end{exercise}

\begin{exercise}\label{ex:A8}
Let $\bz \sim N(\boldsymbol\mu, \boldsymbol\Sigma)$. (a)~For any constant matrix $\bA$ and vector $\mathbf{b}$, prove $\bA\bz + \mathbf{b} \sim N(\bA\boldsymbol\mu + \mathbf{b},\;\bA\boldsymbol\Sigma\bA')$ via the moment generating function. (b)~Specialize to $\bz \sim N(\bzero, \bI_n)$ and let $\bA$ be symmetric idempotent of rank $r$; prove $\bz'\bA\bz \sim \chi^2_r$ via the spectral decomposition.
\end{exercise}

\begin{exercise}\label{ex:A9}
\textbf{(Independence of quadratic forms.)} Let $\bz \sim N(\bzero, \bI_n)$ and let $\bA_1, \bA_2$ be symmetric idempotent with $\bA_1\bA_2 = \bzero$. Prove that $\bz'\bA_1\bz$ and $\bz'\bA_2\bz$ are independent. \emph{Hint:} use joint normality of $(\bA_1\bz, \bA_2\bz)$ and the equivalence ``uncorrelated $\Leftrightarrow$ independent'' under joint normality.
\end{exercise}

\begin{exercise}\label{ex:A10}
Under (A1)--(A5), prove $\SSE/\sigma^2 \sim \chi^2_{n-k-1}$ and that $\SSE$ is independent of $\hat{\bb}$.
\end{exercise}

\begin{exercise}\label{ex:A11}
Prove the variance decomposition $\SST = \SSR + \SSE$, and derive the overall $F$ statistic including the proof that, under (A1)--(A5) and $k \geqslant 1$, $\SSR/\sigma^2 \sim \chi^2_k$ under $H_0\colon \beta_1 = \cdots = \beta_k = 0$.
\end{exercise}

\exercisetier{B}
\subsection*{Tier B: Reproduce and Specialize --- Inference and Prediction (Sections~\ref{sec:ols}, \ref{sec:finite}, \ref{sec:test}, \ref{sec:pred}; Exercise~\ref{ex:B5} extends Section~\ref{sec:finite} to non-spherical errors)}

\begin{exercise}\label{ex:B1}
Let $\bH = \bX(\bX'\bX)^{-1}\bX'$. (a)~Show $0 \leqslant h_{ii} \leqslant 1$. (b)~Show $\sum_{i=1}^n h_{ii} = k+1$. (c)~Under (A3), derive $\var(e_i) = \sigma^2(1-h_{ii})$.
\end{exercise}

\begin{exercise}\label{ex:B2}
Consider testing $H_0\colon \bc'\bb = r_0$ for a given nonzero vector $\bc$, under (A1)--(A5). (a)~Derive the distribution of $\bc'\hat{\bb}$. (b)~Construct the $t$ statistic and prove $t^2 \sim F_{1, n-k-1}$ under $H_0$.
\end{exercise}

\begin{exercise}\label{ex:B3}
Consider two nested models: the full model with $k$ regressors and a reduced model with $k_0 < k$ regressors. Let $\SSE_R$ and $\SSE_U$ denote the respective residual sums of squares. (a)~Prove $\SSE_R \geqslant \SSE_U$. (b)~Under (A1)--(A5), show that $\frac{(\SSE_R - \SSE_U)/(k-k_0)}{\SSE_U/(n-k-1)} \sim F_{k-k_0, n-k-1}$ under $H_0\colon \beta_{k_0+1} = \cdots = \beta_k = 0$, where the reduced model retains the intercept and the first $k_0$ regressors.
\end{exercise}

\begin{exercise}\label{ex:B4}
\textbf{(General linear hypothesis.)} For $H_0\colon \bR\bb = \br$ with $\bR$ a $q\times(k+1)$ matrix of rank $q$, derive the test statistic
\[
F = \frac{(\bR\hat{\bb} - \br)'[\bR(\bX'\bX)^{-1}\bR']^{-1}(\bR\hat{\bb} - \br)/q}{\SSE/(n-k-1)}
\]
and prove $F \sim F_{q, n-k-1}$ under $H_0$ and (A1)--(A5).
\end{exercise}

\begin{exercise}\label{ex:B5}
Assume \ref{A1}, \ref{A2} and \ref{A4}, but suppose that in place of \ref{A3} we have $\var(\be) = \sigma^2\boldsymbol\Omega$, where $\boldsymbol\Omega$ is known, symmetric, and positive definite (but $\neq \bI$). (a)~Show that OLS is unbiased but derive its actual (incorrect) variance. (b)~Define $\hat{\bb}_{\mathrm{GLS}} = (\bX'\boldsymbol\Omega^{-1}\bX)^{-1}\bX'\boldsymbol\Omega^{-1}\bY$ and prove it is BLUE.
\end{exercise}

\begin{exercise}\label{ex:B6}
Let $\hat{Y}_0 = \bx_0'\hat{\bb}$ be the predicted value at a new point $\bx_0$, where $Y_0 = \bx_0'\bb + \varepsilon_0$ with $\expc[\varepsilon_0] = 0$, $\var(\varepsilon_0) = \sigma^2$ and $\varepsilon_0$ independent of $\be$. Assume (A1)--(A4). (a)~Derive $\var(\hat{Y}_0 - \expc[Y_0])$. (b)~Derive $\var(\hat{Y}_0 - Y_0)$. (c)~Explain why (b) always exceeds (a).
\end{exercise}

\exercisetier{C}
\subsection*{Tier C: Reproduce the Derivation --- Partitioned Regression, Model Selection and Constraints (Sections~\ref{sec:R2}, \ref{sec:fwl}, \ref{sec:special})}

\begin{exercise}\label{ex:C1}
\textbf{(Frisch--Waugh--Lovell.)} Partition $\bX = [\bX_1 \;\; \bX_2]$. Let $\bM_2 = \bI - \bX_2(\bX_2'\bX_2)^{-1}\bX_2'$. Prove: (a)~$\hat{\bb}_1 = (\bX_1'\bM_2\bX_1)^{-1}\bX_1'\bM_2\bY$. (b)~The full regression residuals equal those from regressing $\bM_2\bY$ on $\bM_2\bX_1$.
\end{exercise}

\begin{exercise}\label{ex:C2}
\textbf{(Omitted variable bias.)} The true model is $\bY = \bX_1\bb_1 + \bX_2\bb_2 + \be$, but only $\bY = \bX_1\boldsymbol\gamma + \mathbf{u}$ is estimated. Show $\expc[\tilde{\boldsymbol\gamma}] = \bb_1 + (\bX_1'\bX_1)^{-1}\bX_1'\bX_2\bb_2$. Under what conditions is $\tilde{\boldsymbol\gamma}$ unbiased?
\end{exercise}

\begin{exercise}\label{ex:C3}
\textbf{(Monotonicity of $R^2$ and the adjusted $R^2$ criterion.)} (a)~Prove that adding a regressor cannot decrease $R^2$. (b)~Show that $\bar{R}^2$ increases if and only if the added variable's $|t|$-statistic exceeds $1$.
\end{exercise}

\begin{exercise}\label{ex:C4}
\textbf{(Constrained least squares.)} Consider $\min_{\bb}\|\bY - \bX\bb\|^2$ subject to $\bR\bb = \br$. (a)~Derive the constrained estimator $\hat{\bb}_R$ via Lagrange multipliers. (b)~Show
\[
\hat{\bb}_R = \hat{\bb} - (\bX'\bX)^{-1}\bR'[\bR(\bX'\bX)^{-1}\bR']^{-1}(\bR\hat{\bb} - \br).
\]
(c)~Prove $\SSE_R \geqslant \SSE_U$.
\end{exercise}

\exercisetier{D}
\subsection*{Tier D: Apply and Extend}

The exercises in this tier apply the theory to objects that do not appear in the text: a data set to be worked by hand, situations in which one has to decide which assumption fails, and two small propositions not proved above (Exercises~\ref{ex:D3}(a) and~\ref{ex:D4}; Exercise~\ref{ex:D3}(b) is the identity in the proof of Corollary~\ref{cor:slr-var} with $x_0 = x_i$, worked as a check on Exercise~\ref{ex:D1}(d)).

\begin{exercise}\label{ex:D1}
\textbf{(A regression by hand.)} Five observations of a regressor $x$ and a response $Y$ are
\[
\begin{array}{c|ccccc}
x_i & 1 & 2 & 3 & 4 & 5\\ \hline
Y_i & 1 & 2 & 2 & 4 & 6
\end{array}
\]
Fit the model $Y_i = \beta_0 + \beta_1 x_i + \varepsilon_i$ under (A1)--(A5). (a)~Compute $\bX'\bX$, $\bX'\bY$, $(\bX'\bX)^{-1}$ and $\hat{\bb}$. (b)~Compute the fitted values $\hat{\bY}$ and residuals $\mathbf{e}$, and verify $\sum_i e_i = 0$ and $\sum_i x_i e_i = 0$ (Theorem~\ref{thm:normal-eq}). (c)~Compute $\SSE$ and $\MSE = s^2$. (d)~Compute the leverages $h_{11}, \ldots, h_{55}$ and verify $0 \leqslant h_{ii} \leqslant 1$ and $\sum_i h_{ii} = 2$ (Theorem~\ref{thm:leverage}). (e)~Compute $\SST$, $\SSR$, verify $\SST = \SSR + \SSE$ (Theorem~\ref{thm:anova}), and compute $R^2$ and $\bar{R}^2$. (f)~Compute the $t$ statistic for $H_0\colon \beta_1 = 0$ and the overall $F$ statistic, and verify $F = t^2$ (Corollary~\ref{cor:drop-one}). (g)~Using $t_{0.025,3} = 3.182$, give a $95\%$ confidence interval for $\beta_1$, and at $x_0 = 6$ a $95\%$ confidence interval for the conditional mean and a $95\%$ prediction interval for a new observation (Theorems~\ref{thm:ci-mean} and~\ref{thm:pred}).
\end{exercise}

\begin{exercise}\label{ex:D2}
\textbf{(Which assumption fails?)} For each situation, state which of (A1)--(A5) is violated, and which of the following still hold: unbiasedness of $\hat{\bb}$ (Theorem~\ref{thm:unbiased}), the variance formula $\var(\hat{\bb}) = \sigma^2(\bX'\bX)^{-1}$ (Theorem~\ref{thm:var-beta}), the Gauss--Markov property (Theorem~\ref{thm:gm}), and the exact $t$ and $F$ tests of Section~\ref{sec:test}.
\begin{enumerate}[label=\textup{(\roman*)}]
\item Household expenditure is regressed on income; the scatter of the observations around the fitted line widens visibly as income grows.
\item Wages are regressed on years of schooling, an intercept, and indicator variables for \emph{all four} regions of a country.
\item Test scores are regressed on class size, but family income---which is correlated with class size and affects test scores---is not observed and is left out.
\item Monthly sales are regressed on monthly advertising expenditure; the errors of consecutive months are positively correlated.
\item The response is the indicator $Y_i \in \{0, 1\}$ of whether individual $i$ is employed, regressed on age and schooling.
\end{enumerate}
\end{exercise}

\begin{exercise}\label{ex:D3}
\textbf{(A lower bound for leverages.)} (a)~Show that in a model with an intercept, $h_{ii} \geqslant 1/n$ for every $i$, so that together with Theorem~\ref{thm:leverage}, $1/n \leqslant h_{ii} \leqslant 1$. \emph{Hint:} apply the argument of Theorem~\ref{thm:leverage}\ref{lev:i} to the symmetric idempotent matrix $\bH - \bJ$ of Lemma~\ref{lem:SSR-chisq}, Step~2. (b)~For the simple regression model $Y_i = \beta_0 + \beta_1 x_i + \varepsilon_i$, show that $h_{ii} = \dfrac{1}{n} + \dfrac{(x_i - \bar{x})^2}{\sum_{l=1}^n (x_l - \bar{x})^2}$, and check this formula against the leverages found in Exercise~\ref{ex:D1}(d).
\end{exercise}

\begin{exercise}\label{ex:D4}
\textbf{(Adding an orthogonal regressor.)} Let $\bX_{k+1} = [\bX_k \;\; \bx]$ with $\bX_k'\bx = \bzero$ and $\bx \neq \bzero$. Show: (a)~$\hat{\beta}_{k+1} = \bx'\bY/(\bx'\bx)$; (b)~the coefficients on $\bX_k$ are the same in the short regression (on $\bX_k$) and in the enlarged regression (on $\bX_{k+1}$); (c)~$\SSE_k - \SSE_{k+1} = (\bx'\bY)^2/(\bx'\bx)$. (d)~Relate (b) to Theorem~\ref{thm:ovb}.
\end{exercise}

\end{document}

% The solutions of clm.tex follow, outside the compiled slide edition.

%=============================================================================
\section{Solutions to Exercises}
%=============================================================================

\subsection*{Solutions to Tier A}

\textbf{Solution to Exercise~\ref{ex:A1}.} $\ba'\bb = \sum_i a_i\beta_i$ gives $\partial(\ba'\bb)/\partial\beta_j = a_j$, i.e.\ $D(\ba'\bb) = \ba'$. In $\bb'\bA\bb = \sum_{i,l}\beta_iA_{il}\beta_l$ the coordinate $\beta_j$ occurs in the terms with $i = j$ and in those with $l = j$, so
\[
\frac{\partial(\bb'\bA\bb)}{\partial\beta_j} = \sum_l A_{jl}\beta_l + \sum_i \beta_iA_{ij} = (\bA\bb)_j + (\bA'\bb)_j, \quad\text{i.e.}\quad D(\bb'\bA\bb) = \bb'(\bA + \bA')
\]
(this is the proof of Lemma~\ref{lem:matrix-diff}). Applying (i) with $\ba = \bX'\bY$ and (ii) with the symmetric $\bA = \bX'\bX$ to $Q(\bb) = \bY'\bY - 2\bb'\bX'\bY + \bb'\bX'\bX\bb$ gives $DQ(\bb) = -2(\bX'\bY)' + 2\bb'\bX'\bX = -2\bY'\bX + 2\bb'\bX'\bX$.

\textbf{Solution to Exercise~\ref{ex:A2}.} (a)~If $\bA\bv = \lambda\bv$ with $\bv \neq \bzero$, then $\lambda\bv = \bA\bv = \bA^2\bv = \lambda^2\bv$, so $\lambda^2 = \lambda$ and $\lambda \in \{0, 1\}$. (b)~$\tr(\bA) = \rank(\bA)$ by Lemma~\ref{lem:idempotent}, which holds for every idempotent matrix. For symmetric $\bA$ the trace can also be read off the spectral decomposition: by~\ref{L3}, $\bA = \bP\boldsymbol\Lambda\bP'$ with $\bP$ orthogonal and $\boldsymbol\Lambda$ diagonal with entries in $\{0,1\}$ by (a), so $\tr(\bA) = \tr(\boldsymbol\Lambda\bP'\bP) = \tr(\boldsymbol\Lambda)$ (\ref{L2}) is the number of unit eigenvalues, and by the lemma this number is the rank. (c)~$\bH$ and $\bM$ are symmetric idempotent (Theorem~\ref{thm:HM}\ref{HM:i}) and $\calC(\bH) = \calC(\bX)$ has dimension $k+1$ (Theorem~\ref{thm:HM}\ref{HM:iv}, Step~1), so $\tr(\bH) = \rank(\bH) = k+1$ and $\tr(\bM) = \tr(\bI_n) - \tr(\bH) = n-k-1$.

\textbf{Solution to Exercise~\ref{ex:A3}.} Expand $Q(\bb) = \bY'\bY - 2\bb'\bX'\bY + \bb'\bX'\bX\bb$ and differentiate with Lemma~\ref{lem:matrix-diff}: $DQ(\bb) = -2\bY'\bX + 2\bb'\bX'\bX$. Setting it to zero gives the normal equation $\bX'\bX\hat{\bb} = \bX'\bY$. Under (A4), $\bX'\bX\bv = \bzero$ implies $\|\bX\bv\|^2 = \bv'\bX'\bX\bv = 0$, so $\bX\bv = \bzero$ and $\bv = \bzero$; hence $\bX'\bX$ is invertible and $\hat{\bb} = (\bX'\bX)^{-1}\bX'\bY$. Second-order condition: the Hessian $2\bX'\bX$ satisfies $\bv'(2\bX'\bX)\bv = 2\|\bX\bv\|^2 > 0$ for $\bv \neq \bzero$; the minimum is global and unique by $Q(\bb) = Q(\hat{\bb}) + \|\bX(\bb - \hat{\bb})\|^2$, equation~\eqref{eq:ols-global}.

\textbf{Solution to Exercise~\ref{ex:A4}.} (a)~$\bH' = \bX[(\bX'\bX)^{-1}]'\bX' = \bX(\bX'\bX)^{-1}\bX' = \bH$ since $(\bX'\bX)^{-1}$ is symmetric; $\bH^2 = \bX(\bX'\bX)^{-1}(\bX'\bX)(\bX'\bX)^{-1}\bX' = \bH$; then $\bM' = \bI - \bH' = \bM$ and $\bM^2 = \bI - 2\bH + \bH^2 = \bM$. (b)~$\bH\bX = \bX(\bX'\bX)^{-1}\bX'\bX = \bX$, $\bM\bX = \bX - \bH\bX = \bzero$, $\bH\bM = \bH - \bH^2 = \bzero$. (c)~By Lemma~\ref{lem:idempotent}, $\tr(\bH) = \rank(\bH)$; every column of $\bH$ is $\bX$ times a vector, so $\calC(\bH) \subseteq \calC(\bX)$, and $\bX = \bH\bX$ gives the reverse inclusion; hence $\rank(\bH) = \dim\calC(\bX) = k+1$ by (A4), and $\tr(\bM) = n - (k+1)$.

\textbf{Solution to Exercise~\ref{ex:A5}.} Substituting $\bY = \bX\bb + \be$: $\hat{\bb} = (\bX'\bX)^{-1}\bX'(\bX\bb + \be) = \bb + (\bX'\bX)^{-1}\bX'\be$, where (A4) is what makes $(\bX'\bX)^{-1}$ exist. Taking expectations and using $\expc[\be] = \bzero$ (A2) gives $\expc[\hat{\bb}] = \bb$. For the variance, $\hat{\bb}$ is the constant $\bb$ plus the linear transformation $(\bX'\bX)^{-1}\bX'\be$, so Lemma~\ref{lem:var-linear}\ref{vl:i} and (A3) give $\var(\hat{\bb}) = (\bX'\bX)^{-1}\bX'(\sigma^2\bI)\bX(\bX'\bX)^{-1} = \sigma^2(\bX'\bX)^{-1}$.

\textbf{Solution to Exercise~\ref{ex:A6}.} \emph{Statement:} under (A1)--(A4), for every linear estimator $\tilde{\bb} = \bC\bY$ that is unbiased for all $\bb$, $\var(\tilde{\bb}) - \var(\hat{\bb})$ is positive semi-definite, with equality iff $\tilde{\bb} = \hat{\bb}$. \emph{Proof:} $\expc[\bC\bY] = \bC\bX\bb = \bb$ for all $\bb$ forces $\bC\bX = \bI_{k+1}$. Write $\bC = (\bX'\bX)^{-1}\bX' + \bD$; then $\bD\bX = \bC\bX - \bI = \bzero$. By Lemma~\ref{lem:var-linear} and (A3), $\var(\tilde{\bb}) = \sigma^2\bC\bC' = \sigma^2(\bX'\bX)^{-1} + \sigma^2(\bX'\bX)^{-1}\bX'\bD' + \sigma^2\bD\bX(\bX'\bX)^{-1} + \sigma^2\bD\bD'$, and the two middle terms vanish because $\bD\bX = \bzero$ and $\bX'\bD' = (\bD\bX)'$. Hence $\var(\tilde{\bb}) - \var(\hat{\bb}) = \sigma^2\bD\bD' \succeq \bzero$, with equality iff $\bD = \bzero$, i.e.\ $\tilde{\bb} = \hat{\bb}$.

\textbf{Solution to Exercise~\ref{ex:A7}.} $\mathbf{e} = \bM\bY = \bM(\bX\bb + \be) = \bM\be$ since $\bM\bX = \bzero$, so $\SSE = \be'\bM'\bM\be = \be'\bM\be$ ($\bM$ symmetric idempotent). By Lemma~\ref{lem:quadratic-form} with $\expc[\be] = \bzero$ (A2) and $\var(\be) = \sigma^2\bI$ (A3), $\expc[\SSE] = \tr(\bM\cdot\sigma^2\bI) = \sigma^2\tr(\bM) = \sigma^2(n-k-1)$ (Theorem~\ref{thm:HM}\ref{HM:iv}). Dividing by $n-k-1$ gives $\expc[\MSE] = \sigma^2$.

\textbf{Solution to Exercise~\ref{ex:A8}.} (a)~$M_{\bA\bz+\mathbf{b}}(\mathbf{t}) = \expc[e^{\mathbf{t}'(\bA\bz+\mathbf{b})}] = e^{\mathbf{t}'\mathbf{b}}M_{\bz}(\bA'\mathbf{t}) = \exp\bigl(\mathbf{t}'(\bA\boldsymbol\mu+\mathbf{b}) + \tfrac12\mathbf{t}'\bA\boldsymbol\Sigma\bA'\mathbf{t}\bigr)$, the MGF of $N(\bA\boldsymbol\mu+\mathbf{b}, \bA\boldsymbol\Sigma\bA')$; conclude by uniqueness~\ref{N1}. (b)~Write $\bA = \bP\boldsymbol\Lambda\bP'$ (\ref{L3}) with $\boldsymbol\Lambda = \diag(1,\ldots,1,0,\ldots,0)$ ($r$ ones: the eigenvalues lie in $\{0,1\}$ by idempotency and $\tr\bA = \rank\bA = r$). Then $\bw = \bP'\bz \sim N(\bzero, \bP'\bP) = N(\bzero, \bI_n)$ by (a), its coordinates are i.i.d.\ $N(0,1)$ (the MGF $e^{\mathbf{t}'\mathbf{t}/2}$ factorizes, \ref{N2}), and $\bz'\bA\bz = \bw'\boldsymbol\Lambda\bw = \sum_{i=1}^r w_i^2 \sim \chi^2_r$ by definition~\ref{N5}.

\textbf{Solution to Exercise~\ref{ex:A9}.} $(\mathbf{U}', \mathbf{V}')' = ((\bA_1\bz)', (\bA_2\bz)')'$ is a linear image of $\bz$, hence jointly normal (Theorem~\ref{thm:mvn-linear}), with $\cov(\mathbf{U}, \mathbf{V}) = \bA_1\bI\bA_2' = \bA_1\bA_2 = \bzero$ (Lemma~\ref{lem:var-linear}\ref{vl:ii}). By Theorem~\ref{thm:mvn-uncorr-indep}, $\mathbf{U} \indep \mathbf{V}$. Since $\bA_i$ is symmetric idempotent, $\bz'\bA_i\bz = \|\bA_i\bz\|^2$ is a function of $\mathbf{U}$ (resp.\ $\mathbf{V}$) alone, so the two quadratic forms are independent (\ref{N3}).

\textbf{Solution to Exercise~\ref{ex:A10}.} Let $\bz = \be/\sigma \sim N(\bzero, \bI_n)$ (Theorem~\ref{thm:mvn-linear}). Then $\SSE/\sigma^2 = \bz'\bM\bz$ with $\bM$ symmetric idempotent of trace $n-k-1$, so $\SSE/\sigma^2 \sim \chi^2_{n-k-1}$ by Theorem~\ref{thm:quadratic-chisq}. For independence, stack $\hat{\bb} - \bb = (\bX'\bX)^{-1}\bX'\be$ and $\mathbf{e} = \bM\be$: the pair is jointly normal by Theorem~\ref{thm:mvn-linear}, with cross-covariance $\sigma^2(\bX'\bX)^{-1}\bX'\bM = \bzero$ because $\bX'\bM = (\bM\bX)' = \bzero$. Theorem~\ref{thm:mvn-uncorr-indep} gives $\hat{\bb} \indep \mathbf{e}$, hence $\hat{\bb} \indep \SSE = \mathbf{e}'\mathbf{e}$ (\ref{N3}).

\textbf{Solution to Exercise~\ref{ex:A11}.} \emph{Decomposition:} $\bM_0\bY = \bM_0\hat{\bY} + \bM_0\mathbf{e} = \bM_0\hat{\bY} + \mathbf{e}$ because $\bone'\mathbf{e} = 0$ (Theorem~\ref{thm:normal-eq}\ref{ne:ii}); the cross term $\mathbf{e}'\bM_0\hat{\bY} = \mathbf{e}'\hat{\bY} - \bar Y\,\mathbf{e}'\bone = 0$ (Theorem~\ref{thm:normal-eq}\ref{ne:ii}, \ref{ne:iii}), so $\SST = \|\bM_0\bY\|^2 = \|\bM_0\hat{\bY}\|^2 + \|\mathbf{e}\|^2 = \SSR + \SSE$ (using $\overline{\hat Y} = \bar Y$). \emph{$F$ statistic:} $\hat{\bY} - \bar Y\bone = (\bH - \bJ)\bY$, and $\bH - \bJ$ is symmetric idempotent of rank $k$ because $\bH\bJ = \bJ = \bJ\bH$ (from $\bH\bone = \bone$). Under $H_0$, $\bX\bb = \beta_0\bone$ and $(\bH - \bJ)\bone = \bzero$, so $\SSR = \be'(\bH - \bJ)\be$ and $\SSR/\sigma^2 \sim \chi^2_k$ (Theorem~\ref{thm:quadratic-chisq}); $(\bH - \bJ)\bM = \bzero$ gives $\SSR \indep \SSE$ (Theorem~\ref{thm:cochran}); with $\SSE/\sigma^2 \sim \chi^2_{n-k-1}$ (Theorem~\ref{thm:sse-dist}), $F = (\SSR/k)/(\SSE/(n-k-1)) \sim F_{k,n-k-1}$ by definition~\ref{N5}.

\subsection*{Solutions to Tier B}

\textbf{Solution to Exercise~\ref{ex:B1}.} (a)~$\bH = \bH'\bH$ gives $h_{ii} = \sum_j h_{ij}^2 = h_{ii}^2 + \sum_{j \neq i} h_{ij}^2$, so $h_{ii} \geqslant 0$ and $h_{ii} \geqslant h_{ii}^2$, i.e.\ $0 \leqslant h_{ii} \leqslant 1$. (b)~$\sum_i h_{ii} = \tr(\bH) = \rank(\bH) = k+1$ (Lemma~\ref{lem:idempotent}, Theorem~\ref{thm:HM}\ref{HM:iv}). (c)~$\mathbf{e} = \bM\be$, so $\var(\mathbf{e}) = \bM(\sigma^2\bI)\bM' = \sigma^2\bM$ (Lemma~\ref{lem:var-linear}, A3), whose $(i,i)$ entry is $\sigma^2(1 - h_{ii})$.

\textbf{Solution to Exercise~\ref{ex:B2}.} (a)~By Theorem~\ref{thm:mvn-linear}, $\bc'\hat{\bb} \sim N(\bc'\bb,\; \sigma^2\bc'(\bX'\bX)^{-1}\bc)$. (b)~Under $H_0$, $Z = (\bc'\hat{\bb} - r_0)/(\sigma\sqrt{\bc'(\bX'\bX)^{-1}\bc}) \sim N(0,1)$ by (a); $V = \SSE/\sigma^2 \sim \chi^2_{n-k-1}$ is independent of $\hat{\bb}$ (Theorem~\ref{thm:sse-dist}), hence of $Z$ (\ref{N3}); so $t = Z/\sqrt{V/(n-k-1)} = (\bc'\hat{\bb} - r_0)/(s\sqrt{\bc'(\bX'\bX)^{-1}\bc}) \sim t_{n-k-1}$ (\ref{N5}). Squaring, $t^2 = (Z^2/1)/(V/(n-k-1))$ with $Z^2 \sim \chi^2_1$ (definition~\ref{N5} with $r = 1$) independent of $V$, so $t^2 \sim F_{1,n-k-1}$; it coincides with the $F$ statistic of Theorem~\ref{thm:general-F} for $\bR = \bc'$, $\br = r_0$ (Corollary~\ref{cor:t2F}).

\textbf{Solution to Exercise~\ref{ex:B3}.} Let $\bX_U$ be the full design matrix and $\bX_R$ the submatrix of its first $k_0 + 1$ columns (intercept and the $k_0$ retained regressors), with hat matrices $\bH_U = \bH$, $\bH_R = \bX_R(\bX_R'\bX_R)^{-1}\bX_R'$ and residual makers $\bM_U = \bM$, $\bM_R = \bI - \bH_R$; the reduced model is the full model under $H_0\colon \beta_{k_0+1} = \cdots = \beta_k = 0$.

(a)~The reduced model minimizes $\|\bY - \bX_U\bb\|^2$ over the subspace $\{\bb : \beta_{k_0+1} = \cdots = \beta_k = 0\}$ of $\Real^{k+1}$, whereas the unconstrained minimum $\SSE_U$ is taken over all of $\Real^{k+1}$; hence $\SSE_R \geqslant \SSE_U$. (Alternatively, by (b) below, $\SSE_R - \SSE_U = \|(\bH_U - \bH_R)\bY\|^2 \geqslant 0$.)

(b)~\emph{Nested projections.} Each column of $\bX_R$ is a column of $\bX_U$, so $\bH_U\bX_R = \bX_R$ (Theorem~\ref{thm:HM}\ref{HM:ii}); hence $\bH_U\bH_R = \bH_U\bX_R(\bX_R'\bX_R)^{-1}\bX_R' = \bH_R$ and, transposing, $\bH_R\bH_U = \bH_R$. Therefore $\bM_R - \bM_U = \bH_U - \bH_R$ is symmetric, idempotent since $(\bH_U - \bH_R)^2 = \bH_U - \bH_R - \bH_R + \bH_R = \bH_U - \bH_R$, of rank $\tr(\bH_U) - \tr(\bH_R) = (k+1) - (k_0+1) = k - k_0$ (Lemma~\ref{lem:idempotent}), and $(\bM_R - \bM_U)\bM_U = (\bH_U - \bH_R)(\bI - \bH_U) = \bH_U - \bH_U - \bH_R + \bH_R = \bzero$. \emph{Distribution.} $\SSE_R - \SSE_U = \bY'\bM_R\bY - \bY'\bM_U\bY = \bY'(\bM_R - \bM_U)\bY$. Under $H_0$, $\bY = \bX_R\bb_R + \be$ (with $\bb_R$ the first $k_0+1$ coefficients) and $\bM_R\bX_R = \bzero = \bM_U\bX_R$, so $\SSE_R - \SSE_U = \be'(\bM_R - \bM_U)\be$ and likewise $\SSE_U = \be'\bM_U\be$. With $\bz = \be/\sigma \sim N(\bzero, \bI_n)$, Theorem~\ref{thm:quadratic-chisq} gives $(\SSE_R - \SSE_U)/\sigma^2 \sim \chi^2_{k-k_0}$, and Theorem~\ref{thm:cochran} (with $\bA_1 = \bM_R - \bM_U$, $\bA_2 = \bM_U$, $\bA_1\bA_2 = \bzero$) gives independence from $\SSE_U/\sigma^2 \sim \chi^2_{n-k-1}$. Hence $F = \frac{(\SSE_R - \SSE_U)/(k-k_0)}{\SSE_U/(n-k-1)} \sim F_{k-k_0, n-k-1}$ by definition~\ref{N5}. \emph{Via Theorem~\ref{thm:general-F}.} The same $F$ is the general $F$ statistic for $\bR = [\bzero \mid \bI_{k-k_0}]$ ($q = k - k_0$) and $\br = \bzero$: by Corollary~\ref{cor:nested-F}, its numerator $(\bR\hat{\bb})'[\bR(\bX'\bX)^{-1}\bR']^{-1}(\bR\hat{\bb})/q$ equals $(\SSE_R - \SSE_U)/q$.

\textbf{Solution to Exercise~\ref{ex:B4}.} By Theorems~\ref{thm:beta-dist} and~\ref{thm:mvn-linear}, $\bR\hat{\bb} \sim N(\bR\bb, \sigma^2\boldsymbol\Sigma_R)$ with $\boldsymbol\Sigma_R = \bR(\bX'\bX)^{-1}\bR'$, which is positive definite because $\rank\bR = q$ (for $\bv \neq \bzero$, $\bR'\bv \neq \bzero$ and $(\bX'\bX)^{-1}$ is positive definite by Theorem~\ref{thm:ols}, Step~4). Under $H_0$, $\bw = \sigma^{-1}\boldsymbol\Sigma_R^{-1/2}(\bR\hat{\bb} - \br) \sim N(\bzero, \bI_q)$ (\ref{L4}, Theorem~\ref{thm:mvn-linear}), so $W = \bw'\bw = (\bR\hat{\bb} - \br)'\boldsymbol\Sigma_R^{-1}(\bR\hat{\bb} - \br)/\sigma^2 \sim \chi^2_q$ (Theorem~\ref{thm:quadratic-chisq} with $\bA = \bI_q$). $W$ is a function of $\hat{\bb}$, hence independent of $V = \SSE/\sigma^2 \sim \chi^2_{n-k-1}$ (Theorem~\ref{thm:sse-dist}, \ref{N3}), and $F = (W/q)/(V/(n-k-1))$ is the displayed statistic (the factors $\sigma^{-2}$ cancel), distributed as $F_{q,n-k-1}$ by~\ref{N5}.

\textbf{Solution to Exercise~\ref{ex:B5}.} (a)~Unbiasedness: $\expc[\hat{\bb}] = \bb + (\bX'\bX)^{-1}\bX'\expc[\be] = \bb$ (unchanged; unbiasedness does not use~\ref{A3}). Actual variance:
\[
\var(\hat{\bb}) = (\bX'\bX)^{-1}\bX'(\sigma^2\boldsymbol\Omega)\bX(\bX'\bX)^{-1} = \sigma^2(\bX'\bX)^{-1}\bX'\boldsymbol\Omega\bX(\bX'\bX)^{-1},
\]
which in general differs from $\sigma^2(\bX'\bX)^{-1}$, so the usual standard errors are wrong. (Not for \emph{every} $\boldsymbol\Omega \neq \bI$, however: if $\boldsymbol\Omega = \bI + \bv\bv'$ with $\bX'\bv = \bzero$, then $\bX'\boldsymbol\Omega\bX = \bX'\bX$ and the two expressions coincide.)

(b)~Let $\boldsymbol\Omega = \mathbf{L}\mathbf{L}'$ with $\mathbf{L}$ invertible (\ref{L4}; for instance the Cholesky factor, or the symmetric square root $\boldsymbol\Omega^{1/2}$). Transform: $\mathbf{L}^{-1}\bY = \mathbf{L}^{-1}\bX\bb + \mathbf{L}^{-1}\be$, where $\var(\mathbf{L}^{-1}\be) = \sigma^2\mathbf{L}^{-1}\boldsymbol\Omega(\mathbf{L}^{-1})' = \sigma^2\bI$. The OLS estimator applied to the transformed model is
\[
\hat{\bb}_{\mathrm{GLS}} = [(\mathbf{L}^{-1}\bX)'(\mathbf{L}^{-1}\bX)]^{-1}(\mathbf{L}^{-1}\bX)'(\mathbf{L}^{-1}\bY) = (\bX'\boldsymbol\Omega^{-1}\bX)^{-1}\bX'\boldsymbol\Omega^{-1}\bY.
\]
Since the transformed model satisfies (A1)--(A4) (its regressor matrix $\mathbf{L}^{-1}\bX$ has full column rank because $\mathbf{L}^{-1}$ is invertible), Gauss--Markov (Theorem~\ref{thm:gm}) guarantees that $\hat{\bb}_{\mathrm{GLS}}$ is BLUE in the transformed model. It is also BLUE in the original model, because the two models have the \emph{same} class of linear unbiased estimators with the \emph{same} variances: an estimator $\bC\bY$ that is linear in $\bY$ equals $(\bC\mathbf{L})(\mathbf{L}^{-1}\bY)$, which is linear in the transformed response, and conversely every $\tilde{\bC}(\mathbf{L}^{-1}\bY)$ equals $(\tilde{\bC}\mathbf{L}^{-1})\bY$; unbiasedness in the original model, $\bC\bX = \bI$, is the same condition as $(\bC\mathbf{L})(\mathbf{L}^{-1}\bX) = \bI$ in the transformed model; and $\var(\bC\bY)$ is the variance of the same random vector under both descriptions. Hence the estimator that minimizes the variance (in the positive semi-definite order) over the transformed class, namely $\hat{\bb}_{\mathrm{GLS}}$, minimizes it over the original class.

\textbf{Solution to Exercise~\ref{ex:B6}.} (a)~$\hat{Y}_0 - \expc[Y_0] = \bx_0'(\hat{\bb} - \bb)$, so by Lemma~\ref{lem:var-linear} and Theorem~\ref{thm:var-beta}, $\var(\hat{Y}_0 - \expc[Y_0]) = \bx_0'\var(\hat{\bb})\bx_0 = \sigma^2\bx_0'(\bX'\bX)^{-1}\bx_0$. (b)~$\hat{Y}_0 - Y_0 = \bx_0'(\hat{\bb} - \bb) - \varepsilon_0$, where $\varepsilon_0$ is independent of $\be$ and hence of $\bx_0'(\hat{\bb} - \bb)$; the variance of a sum of independent terms is the sum of the variances (fact~\ref{N7}), so $\var(\hat{Y}_0 - Y_0) = \sigma^2\bx_0'(\bX'\bX)^{-1}\bx_0 + \sigma^2 = \sigma^2(1 + \bx_0'(\bX'\bX)^{-1}\bx_0)$ (Theorem~\ref{thm:pred}, Step~2). (c)~The prediction variance exceeds the estimation variance by exactly $\sigma^2$, the irreducible noise variance of $\varepsilon_0$.

\subsection*{Solutions to Tier C}

\textbf{Solution to Exercise~\ref{ex:C1}.} (a)~The second block of the normal equations $\bX'\bX\hat{\bb} = \bX'\bY$ reads $\bX_2'\bX_1\hat{\bb}_1 + \bX_2'\bX_2\hat{\bb}_2 = \bX_2'\bY$, so $\hat{\bb}_2 = (\bX_2'\bX_2)^{-1}\bX_2'(\bY - \bX_1\hat{\bb}_1)$. Substituting into the first block, $\bX_1'\bX_1\hat{\bb}_1 + \bX_1'\bH_2(\bY - \bX_1\hat{\bb}_1) = \bX_1'\bY$ with $\bH_2 = \bI - \bM_2$, which simplifies to $\bX_1'\bM_2\bX_1\hat{\bb}_1 = \bX_1'\bM_2\bY$; here $\bX_1'\bM_2\bX_1 = (\bM_2\bX_1)'(\bM_2\bX_1)$ is positive definite because $\bM_2\bX_1\bv = \bzero$ would put $\bX_1\bv$ in $\calC(\bX_2)$, contradicting (A4) unless $\bv = \bzero$. (b)~The residual from regressing $\bM_2\bY$ on $\bM_2\bX_1$ is $\bM_2\bY - \bM_2\bX_1\hat{\bb}_1 = \bM_2(\bX_2\hat{\bb}_2 + \mathbf{e}) = \bM_2\mathbf{e} = \mathbf{e}$, using $\bM_2\bX_2 = \bzero$ and $\bX_2'\mathbf{e} = \bzero$ (Theorem~\ref{thm:normal-eq}\ref{ne:i}).

\textbf{Solution to Exercise~\ref{ex:C2}.} Substituting the true model into the short-regression estimator, $\tilde{\boldsymbol\gamma} = (\bX_1'\bX_1)^{-1}\bX_1'(\bX_1\bb_1 + \bX_2\bb_2 + \be) = \bb_1 + (\bX_1'\bX_1)^{-1}\bX_1'\bX_2\bb_2 + (\bX_1'\bX_1)^{-1}\bX_1'\be$; taking expectations with $\expc[\be] = \bzero$ (A2) gives the formula. Unbiasedness holds iff $\bX_1'\bX_2\bb_2 = \bzero$, which occurs when $\bX_1'\bX_2 = \bzero$ (every column of $\bX_1$ is orthogonal to every column of $\bX_2$) or $\bb_2 = \bzero$.

\textbf{Solution to Exercise~\ref{ex:C3}.} (a)~With $\calC_k \subseteq \calC_{k+1}$ the two column spaces, $\hat{\bY}_{k+1}$ is the closest point of $\calC_{k+1}$ to $\bY$ (Theorem~\ref{thm:ols-proj}) and $\hat{\bY}_k \in \calC_k \subseteq \calC_{k+1}$, so $\SSE_{k+1} = \|\bY - \hat{\bY}_{k+1}\|^2 \leqslant \|\bY - \hat{\bY}_k\|^2 = \SSE_k$; since $\SST$ is the same in both models, $R^2 = 1 - \SSE/\SST$ cannot decrease. (b)~$\bar{R}^2_{k+1} > \bar{R}^2_k \iff \MSE_{k+1} < \MSE_k \iff \SSE_{k+1}/(n-k-2) < (\SSE_{k+1} + \hat{\beta}_{k+1}^2\,\bx'\bM_k\bx)/(n-k-1)$ (Lemma~\ref{lem:sse-reduction}) $\iff \SSE_{k+1} < (n-k-2)\hat{\beta}_{k+1}^2\,\bx'\bM_k\bx \iff 1 < \hat{\beta}_{k+1}^2\,\bx'\bM_k\bx/\MSE_{k+1}$; the right-hand side is $t_{k+1}^2$ because $\SE(\hat{\beta}_{k+1})^2 = \MSE_{k+1}[(\bX_{k+1}'\bX_{k+1})^{-1}]_{k+1,k+1} = \MSE_{k+1}/(\bx'\bM_k\bx)$ by~\eqref{eq:last-diag}.

\textbf{Solution to Exercise~\ref{ex:C4}.} (a)~With $\mathcal{L}(\bb, \boldsymbol\lambda) = \|\bY - \bX\bb\|^2 + 2\boldsymbol\lambda'(\bR\bb - \br)$, Lemma~\ref{lem:matrix-diff} gives $\partial\mathcal{L}/\partial\bb' = -2\bY'\bX + 2\bb'\bX'\bX + 2\boldsymbol\lambda'\bR = \bzero'$, i.e.\ $\bX'\bX\hat{\bb}_R = \bX'\bY - \bR'\hat{\boldsymbol\lambda}$, so $\hat{\bb}_R = \hat{\bb} - (\bX'\bX)^{-1}\bR'\hat{\boldsymbol\lambda}$; imposing $\bR\hat{\bb}_R = \br$ gives $\hat{\boldsymbol\lambda} = [\bR(\bX'\bX)^{-1}\bR']^{-1}(\bR\hat{\bb} - \br)$ ($\bR(\bX'\bX)^{-1}\bR'$ is positive definite since $\rank\bR = q$). (b)~Substitute $\hat{\boldsymbol\lambda}$ into the expression for $\hat{\bb}_R$. (c)~For any $\bb$ with $\bR\bb = \br$, $\|\bY - \bX\bb\|^2 = \|\bY - \bX\hat{\bb}_R\|^2 + \|\bX(\hat{\bb}_R - \bb)\|^2$, because the cross term is $(\bY - \bX\hat{\bb}_R)'\bX(\hat{\bb}_R - \bb) = \hat{\boldsymbol\lambda}'\bR(\hat{\bb}_R - \bb) = 0$ (Theorem~\ref{thm:cls}, Step~3); so $\hat{\bb}_R$ is the constrained minimizer, and $\SSE_R \geqslant \SSE_U$ because the unconstrained minimum $\SSE_U$ is taken over a larger set.

\subsection*{Solutions to Tier D}

\textbf{Solution to Exercise~\ref{ex:D1}.} (a)~$\bX'\bX = \begin{pmatrix}5 & 15\\ 15 & 55\end{pmatrix}$, $\bX'\bY = (15, 57)'$, $\det(\bX'\bX) = 275 - 225 = 50$, $(\bX'\bX)^{-1} = \frac{1}{50}\begin{pmatrix}55 & -15\\ -15 & 5\end{pmatrix} = \begin{pmatrix}1.1 & -0.3\\ -0.3 & 0.1\end{pmatrix}$, and $\hat{\bb} = (\bX'\bX)^{-1}\bX'\bY = (16.5 - 17.1,\; -4.5 + 5.7)' = (-0.6,\; 1.2)'$. (b)~$\hat{Y}_i = -0.6 + 1.2x_i$: $\hat{\bY} = (0.6, 1.8, 3.0, 4.2, 5.4)'$ and $\mathbf{e} = (0.4, 0.2, -1.0, -0.2, 0.6)'$; $\sum e_i = 0$ and $\sum x_ie_i = 0.4 + 0.4 - 3.0 - 0.8 + 3.0 = 0$. (c)~$\SSE = 0.16 + 0.04 + 1 + 0.04 + 0.36 = 1.6$, $\MSE = 1.6/3 = 0.5333$, $s = 0.7303$. (d)~$h_{ii} = \bx_i'(\bX'\bX)^{-1}\bx_i = 1.1 - 0.6x_i + 0.1x_i^2$, giving $(0.6, 0.3, 0.2, 0.3, 0.6)$: all in $[0,1]$ and summing to $2 = k+1$. (e)~$\bar Y = 3$, $\SST = 4+1+1+1+9 = 16$, $\SSR = \sum(\hat Y_i - 3)^2 = 5.76 + 1.44 + 0 + 1.44 + 5.76 = 14.4$, and indeed $14.4 + 1.6 = 16$; $R^2 = 14.4/16 = 0.9$ and $\bar{R}^2 = 1 - (1.6/3)/(16/4) = 0.8667$. (f)~$\SE(\hat\beta_1) = s\sqrt{0.1} = \sqrt{0.05333} = 0.2309$, $t = 1.2/0.2309 = 5.196$; $F = \MSR/\MSE = 14.4/0.5333 = 27.0 = t^2$. (g)~For $\beta_1$: $1.2 \pm 3.182 \times 0.2309 = 1.2 \pm 0.735$, i.e.\ $(0.465, 1.935)$. At $\bx_0 = (1, 6)'$: $\bx_0'(\bX'\bX)^{-1}\bx_0 = 1.1 - 3.6 + 3.6 = 1.1$ and $\hat Y_0 = -0.6 + 7.2 = 6.6$; the confidence interval for the conditional mean is $6.6 \pm 3.182 \times 0.7303\sqrt{1.1} = 6.6 \pm 2.437$, i.e.\ $(4.163, 9.037)$, and the prediction interval is $6.6 \pm 3.182 \times 0.7303\sqrt{2.1} = 6.6 \pm 3.368$, i.e.\ $(3.232, 9.968)$.

\textbf{Solution to Exercise~\ref{ex:D2}.} (i)~Heteroscedasticity: (A3) fails ($\var(\varepsilon_i)$ depends on income) while (A1)/(A2) may well hold. Unbiasedness survives (Theorem~\ref{thm:unbiased} uses only (A2) and (A4)); the variance formula, the Gauss--Markov property, $\expc[\MSE] = \sigma^2$ and the exact $t$ and $F$ tests all rely on (A3) and fail in general (compare Exercise~\ref{ex:B5}, where an exception is exhibited). (ii)~The dummy variable trap: (A4) fails, because the four indicators sum to the intercept column, $\bone = D_1 + D_2 + D_3 + D_4$; $\bX'\bX$ is singular and $\hat{\bb}$ is not defined. Drop one indicator (or the intercept). (iii)~Omitted variable: the unobserved income is absorbed into the error, which is then correlated with class size, so (A2)---equivalently (A1)---fails for the regression that is estimated. Taking the four properties in turn: unbiasedness fails, the estimated coefficient vector being the short-regression estimator $\tilde{\boldsymbol\gamma}$ of Theorem~\ref{thm:ovb}, with $\expc[\tilde{\boldsymbol\gamma}] = \bb_1 + (\bX_1'\bX_1)^{-1}\bX_1'\bX_2\bb_2 \neq \bb_1$ in general; the variance formula survives as a statement about dispersion, since if the true error $\be$ satisfies (A3) then Lemma~\ref{lem:var-linear} gives $\var(\tilde{\boldsymbol\gamma}) = \sigma^2(\bX_1'\bX_1)^{-1}$, but it is dispersion around the wrong center; the Gauss--Markov property does not apply, because $\tilde{\boldsymbol\gamma}$ is not in the class of unbiased estimators that Theorem~\ref{thm:gm} compares; and the exact $t$ and $F$ tests fail, because under (A5) the statistics are still built from a normal vector, but one centered at $\bb_1 + (\bX_1'\bX_1)^{-1}\bX_1'\bX_2\bb_2$ rather than at $\bb_1$, so the null distributions of Section~\ref{sec:test} no longer hold. (iv)~Serial correlation: (A3) fails through $\cov(\varepsilon_t, \varepsilon_{t+1}) \neq 0$; the consequences are as in (i)---unbiased $\hat{\bb}$, wrong variance formula, invalid tests. (v)~A binary response: (A5) fails, since given $\bx_i$ the error $\varepsilon_i = Y_i - \bx_i'\bb$ takes only the two values $1 - \bx_i'\bb$ and $-\bx_i'\bb$; (A3) also fails, since $\var(\varepsilon_i) = p_i(1 - p_i)$ with $p_i = \bx_i'\bb$ varies with $i$ (and (A1) itself, which requires $p_i = \bx_i'\bb \in [0,1]$ for all $i$, is questionable). Taking the four properties in turn: if (A1) holds, $\hat{\bb}$ remains unbiased (Theorem~\ref{thm:unbiased} needs only (A2) and (A4)); the variance formula fails in general, since Lemma~\ref{lem:var-linear} gives $\var(\hat{\bb}) = (\bX'\bX)^{-1}\bX'\var(\be)\bX(\bX'\bX)^{-1}$ with $\var(\be)$ having diagonal entries $p_i(1-p_i)$ that vary with $i$, so $\var(\be) \neq \sigma^2\bI$; the Gauss--Markov property fails together with (A3), exactly as in (i); and the exact $t$ and $F$ tests fail, since both (A3) and (A5) are violated and the normal-theory distributions of Sections~\ref{sec:exact}--\ref{sec:test} are unavailable.

\textbf{Solution to Exercise~\ref{ex:D3}.} (a)~By Lemma~\ref{lem:SSR-chisq}, Step~2, $\bH - \bJ$ is symmetric idempotent whenever $\bone \in \calC(\bX)$, so $\bH - \bJ = (\bH - \bJ)'(\bH - \bJ)$ and, comparing $(i,i)$ entries exactly as in the proof of Theorem~\ref{thm:leverage}\ref{lev:i}, $h_{ii} - \tfrac{1}{n} = \sum_{j=1}^n (h_{ij} - \tfrac{1}{n})^2 \geqslant 0$. (b)~(The inverse below is Corollary~\ref{cor:slr}\ref{slr:ii}, and the result is the identity in the proof of Corollary~\ref{cor:slr-var} with $x_0 = x_i$.) For $\bX = [\bone \;\; \bx]$, $\bX'\bX = \begin{pmatrix} n & n\bar x\\ n\bar x & \sum_l x_l^2\end{pmatrix}$ has determinant $n\sum_l x_l^2 - n^2\bar x^2 = nS_{xx}$, where $S_{xx} := \sum_l(x_l - \bar x)^2 = \sum_l x_l^2 - n\bar x^2$; so $(\bX'\bX)^{-1} = \frac{1}{nS_{xx}}\begin{pmatrix}\sum_l x_l^2 & -n\bar x\\ -n\bar x & n\end{pmatrix}$ and
\[
h_{ii} = (1, x_i)\,(\bX'\bX)^{-1}(1, x_i)' = \frac{\sum_l x_l^2 - 2n\bar x x_i + n x_i^2}{nS_{xx}} = \frac{S_{xx} + n(x_i - \bar x)^2}{nS_{xx}} = \frac{1}{n} + \frac{(x_i - \bar x)^2}{S_{xx}}.
\]
In Exercise~\ref{ex:D1}, $\bar x = 3$ and $S_{xx} = 10$, giving $h_{ii} = 0.2 + (x_i - 3)^2/10 = (0.6, 0.3, 0.2, 0.3, 0.6)$, as found there.

\textbf{Solution to Exercise~\ref{ex:D4}.} $\bX_k'\bx = \bzero$ gives $\bH_k\bx = \bX_k(\bX_k'\bX_k)^{-1}\bX_k'\bx = \bzero$, so $\bM_k\bx = \bx$, $\bx'\bM_k\bx = \bx'\bx > 0$ (hence Lemma~\ref{lem:sse-reduction} applies) and $\bx'\bM_k\bY = \bx'\bY$. (a)~Equation~\eqref{eq:beta-new}: $\hat\beta_{k+1} = \bx'\bY/(\bx'\bx)$. (b)~In the proof of Lemma~\ref{lem:sse-reduction}, $\hat{\bb}_k^* = \hat{\bb}_k - (\bX_k'\bX_k)^{-1}\bX_k'\bx\,\hat\beta_{k+1} = \hat{\bb}_k$ because $\bX_k'\bx = \bzero$. (c)~Equation~\eqref{eq:sse-reduction}: $\SSE_k - \SSE_{k+1} = (\bx'\bM_k\bY)^2/(\bx'\bM_k\bx) = (\bx'\bY)^2/(\bx'\bx)$. (d)~Read the enlarged model as the true model and the short regression as the one omitting $\bx$: Theorem~\ref{thm:ovb} with $\bX_1 = \bX_k$ and $\bX_2 = \bx$ says that the omitted-variable bias is $(\bX_k'\bX_k)^{-1}\bX_k'\bx\,\beta_{k+1} = \bzero$; part (b) is the exact, sample-level counterpart of this statement about expectations.

\end{document}
